\documentclass[11pt,a4paper]{amsart}

\usepackage[margin=1in]{geometry}
\usepackage{amsmath,amssymb,amsthm,mathtools}
\usepackage[expansion=false]{microtype}
\usepackage{enumitem}
\usepackage[colorlinks=true,linkcolor=black,citecolor=black,urlcolor=black]{hyperref}
\usepackage{cleveref}
\hypersetup{
  pdftitle={Skew Schur equivalence by Jacobi--Trudi negative count},
  pdfauthor={John Shields},
  pdfkeywords={skew Schur function, skew-equivalence, Jacobi--Trudi determinant,
    negative index, rotation rigidity, row overlap, Toeplitz minor}
}
\pdfinfo{ /ORCID (https://orcid.org/0009-0004-3882-4571) }

\numberwithin{equation}{section}
\theoremstyle{plain}
\newtheorem{theorem}{Theorem}[section]
\newtheorem{proposition}[theorem]{Proposition}
\newtheorem{lemma}[theorem]{Lemma}
\newtheorem{corollary}[theorem]{Corollary}
\theoremstyle{definition}
\newtheorem{definition}[theorem]{Definition}
\newtheorem{example}[theorem]{Example}
\newtheorem{remark}[theorem]{Remark}
\newtheorem{problem}[theorem]{Problem}

\newcommand{\Q}{\mathbb Q}
\newcommand{\Z}{\mathbb Z}
\newcommand{\B}{\mathfrak B}

\title[Skew Schur equivalence by negative count]
{Skew Schur equivalence\\
by Jacobi--Trudi negative count}
\author{John Shields}
\date{October 4, 2026}
\thanks{Independent researcher, Ginza 2-15-2, KR Ginza II 5F, Chuo-ku, Tokyo, Japan 104-0061.\\
\hspace*{\parindent}Corresponding author: \href{mailto:math@johnnyshields.com}{math@johnnyshields.com}.  ORCID: \href{https://orcid.org/0009-0004-3882-4571}{0009-0004-3882-4571}.\\
\hspace*{\parindent}This work is archived at \href{https://doi.org/10.5281/zenodo.23133741}{doi:10.5281/zenodo.23133741}.}
\subjclass[2020]{Primary 05E05; Secondary 05A17, 15A15}
\keywords{skew Schur function, skew-equivalence, Jacobi--Trudi determinant, negative index, rotation rigidity, row overlap, Toeplitz minor}

\begin{document}

\begin{abstract}
We classify connected skew diagrams with at most six negative indices in their full row Jacobi--Trudi board, against arbitrary skew diagrams with the same Schur function.  Seven parameter families give all exceptional classes, each with four members; every other source is determined up to rotation.  When the negative pattern requires at least seven rows to realize, nine negatives force rigidity.  At ten negatives under the same minimum-row condition, exactly two seven-row quartet families are nonrigid, and every higher-row source is rigid.  Their complete classes have legal marked-composition presentations.  Classical overlap invariants recover negative count and minimum realizing row count.  Boundary coefficients identify cofactor components, and endpoint comparisons determine their placements.  A saturated first negative column with at most two negatives elsewhere forces rotation rigidity in at least six rows, without a bound on negative count.  Six-row quartets and non-dense placement resonances remain outside the classification.
\end{abstract}

\maketitle

\section{Introduction}\label{sec:introduction}

We classify connected skew diagrams with at most six negative indices in
their full row Jacobi--Trudi board.  The comparison is with an arbitrary
skew diagram having the same skew Schur function, without a bound on area
or row count and without excluding common beta entries.  When the negative
pattern requires at least seven rows to realize, rigidity holds through
nine negatives.  At ten under the same minimum-row condition, exactly
two seven-row quartet families are nonrigid; every such source with more
rows is rigid.

For a connected diagram \(D=\lambda/\mu\) with \(n\) nonempty rows,
its \emph{full row board} is
\[
 (\lambda_i-\mu_j-i+j)_{1\le i,j\le n}.
\]
Write \(\nu(D)\) for the number of strictly negative entries.  A zero
index contributes \(h_0=1\) and is not counted; common beta entries are
not canceled before computing \(\nu\).  Identify diagrams up to
translation, write \(D^*\) for rotation through \(180\) degrees, and put
\[
 [D]=\{H:s_H=s_D,\ H\text{ connected}\}.
\]
Call \(D\) \emph{rotation-rigid} when
\([D]=\{D,D^*\}\), with repetitions removed.

\begin{theorem}[Classification by negative count]\label{thm:classification}
Let \(D\) be a connected skew diagram with \(n\ge1\) nonempty rows,
and let \(H\) be any skew diagram with \(s_H=s_D\).
Common beta entries are allowed in both.
\begin{enumerate}[label=(\roman*),leftmargin=2em]
\item If \(\nu(D)\le2\), then \(H\in\{D,D^*\}\).
\item If \(\nu(D)=3\), the only exceptions to rotation rigidity have four
rows and are the complete four-member classes in
\eqref{eq:four-family}, with \(0\le c<w\) and \(b>c+w\).
\item If \(\nu(D)=4\), then \(H\in\{D,D^*\}\).
\item If \(\nu(D)=5\), the only exceptions have five rows and are the two
complete four-member families in \eqref{eq:five-families}, with
\(1\le q\le b<g\).
\item If \(\nu(D)=6\), the only exceptions have five rows and are the
complete four-member classes M1--M4 of \eqref{eq:staircase-families},
under the parameter inequalities stated there.
\end{enumerate}
\end{theorem}

Every exceptional class has two distinct rotation orbits.  The seven
parameter families are not disjoint: M1 and M2 coincide at their stated
boundary.  The six-negative exceptions have the staircase negative set
\begin{equation}\label{eq:staircase-set}
 \mathcal T_n=\{(i,j):i-j\ge n-3\},\qquad n\ge5.
\end{equation}
Sources with this pattern and at least six rows are rigid, as are all
other six-negative sources.

Read the numbers of negative entries in successive rows from bottom to
top, omitting zero parts, to obtain the \emph{negative partition}
\(\rho(D)\).  For a nonempty partition \(\rho\), let
\(n_{\min}(\rho)\) be the least number of nonempty rows in a connected
diagram realizing it.  We prove
\begin{equation}\label{eq:minimum-intro}
 n_{\min}(\rho)=1+\max_{1\le i\le\ell(\rho)}(\rho_i+i),
 \qquad n_{\min}(\varnothing)=1.
\end{equation}
This is a minimum over realizations of a pattern, not over Schur-equal
diagrams or the actual height of the source.

\begin{theorem}[Negative-index rigidity threshold]\label{thm:negative-threshold}
Let \(D\) be a connected skew diagram satisfying
\[
 \nu(D)\le9,\qquad n_{\min}(\rho(D))\ge7.
\]
For every skew diagram \(H\),
\[
 s_H=s_D\quad\Longrightarrow\quad H\in\{D,D^*\},
\]
where diagrams are identified up to translation and \(D^*\) denotes rotation through \(180\) degrees.  Common beta entries are allowed.  The bound nine is sharp: the two seven-row families in \Cref{prop:ten-negative-quartet} have ten negatives and minimum realizing row count seven, with nonrotation Schur-equal members.
\end{theorem}

\begin{theorem}[Ten-negative classification]\label{thm:ten-negative-classification}
Let \(D\) be connected with
\[
 \nu(D)=10,\qquad n_{\min}(\rho(D))=7.
\]
Its complete connected Schur class is \(\{D,D^*\}\), except for the
classes
\[
 \{F_i(a,b,d),F_i(a,b,d)^*,H_i(a,b,d),H_i(a,b,d)^*\},
 \qquad i\in\{A,B\},
\]
with the profiles in \eqref{eq:ten-family-a-profiles} and
\eqref{eq:ten-profiles}, for integers \(a,b,d\) satisfying
\(a\ge1,d\ge0,b>a+d\).
Each exceptional class has four distinct members and seven nonempty rows,
with proper legal charts in \eqref{eq:ten-chart}.
Every source in this stratum with at least eight rows is rotation-rigid.
Common beta entries are allowed, and diagrams are identified up to translation.
\end{theorem}

Together with \Cref{cor:ten-negative-high-minimum}, this classifies the
entire ten-negative domain with minimum realizing row count at least seven.
The minimum-six case is not included.
Negative count and minimum realizing row count are Schur invariants, but
the negative partition need not be.  Rotation conjugates it.  The equal
thirteen-cell profiles in \Cref{ex:negative-shape-change} have partitions
\((3,1,1)\) and \((3,2)\), which are neither equal nor conjugate.
None of these theorems prescribes a mate's negative partition.

Billera, Thomas and van Willigenburg classified ribbon equality by factor
reversal \cite{Billera2006DecomposableCompositions}.  Reiner, Shaw and
van Willigenburg established the overlap invariants used here
\cite{Reiner2007Coincidences}.  We use connected-skew-Schur irreducibility
in the corrected form of Barekat, Reiner and van Willigenburg
\cite{Barekat2009CoincidencesCorrigendum}.
McNamara and van Willigenburg's marked compositions provide a broader
source of equalities \cite{McNamara2009Classification}.  The Toeplitz-minor correspondence and row/column
reversal used below are classical \cite{GarciaGarcia2020ToeplitzMinors}.

The maximal-label coefficient recovers component functions, not their
positions or orientations.
For five negatives, opening one second-subdiagonal cell gives a
sixteen-plus-eight expansion.  For the staircase, minimum-gap
specialization recovers primitive factors before the full coefficients
assemble them.  At minimum row count six, the remaining component budget
leaves only a singleton beside a two-negative four-row component or two
dense components.  Endpoint coefficients settle both cases.

The singleton argument also gives \Cref{thm:saturated-column}: at six
or more rows, a saturated first negative column and at most two other
negative entries force rotation rigidity.  This criterion has no bound
on the full negative count, but retains its specific column geometry.

At higher minimum row count, the nine-negative budget leaves a dense
four-row component and two singletons.  A three-variable specialization
proves that its endpoint difference has a nonzero lower remainder.
At ten negatives, the four-row component can have one negative, and a
singleton can accompany an exceptional five-row core.  Rank-two and
rank-three specializations settle the former; an endpoint-coefficient
table determines every exceptional singleton context in the latter.
Connected bordering excludes further nonrotation classes at larger
heights.  The six-row quartets in \Cref{prop:six-row-quartets} remain
lower inclusions rather than complete class descriptions.

Section~\ref{sec:boards} gives polynomial readings, negative partitions,
and their overlap invariants.  Section~\ref{sec:boundary-tools} develops
homogenization and bordering.  Sections~\ref{sec:low}--\ref{sec:staircase-assembly}
classify the small-row exceptions and exclude their higher-row borders.
Section~\ref{sec:component-reductions} treats disconnected cofactors,
including the saturated-column criterion and the other minimum-six
patterns.  Sections~\ref{sec:dense-endpoints}--\ref{sec:threshold-proof}
prove the sharp threshold and finish the six-negative classification.
Section~\ref{sec:threshold-boundaries} gives the sharpness families and
boundary examples.  Section~\ref{sec:ten-negative} classifies the
ten-negative stratum at minimum row count seven and excludes its higher-row
borders.  Appendix~\ref{app:four} contains the four-row calculation;
Appendix~\ref{app:contact-first-jet} gives a further component-determination
result, not an input to the parent classification.

\section{Full boards and invariant negative count}\label{sec:boards}

\subsection{Profiles and polynomial readings}

Work in \(\Q[h_1,h_2,\ldots]\), with \(h_0=1\) and \(h_j=0\) for
\(j<0\); see \cite{Macdonald1995SymmetricFunctions}.  For one-based row
indices put \(x_i=\lambda_i-i\), \(y_i=\mu_i-i\), and translate both
strictly decreasing lists to have \(y_n=0\).  Then
\[
 s_D=\det(h_{x_i-y_j}),\qquad N=x_1,
 \qquad r_i=x_i-y_i,
 \qquad z_i=x_{i+1}-y_i.
\]
The \emph{profile} \((r;z)\) records row lengths and adjacent overlaps minus
one.  Its exact connectedness conditions are
\begin{equation}\label{eq:profile-domain}
 r_i>0,\qquad 0\le z_i<\min(r_i,r_{i+1}).
\end{equation}
It determines the diagram: starting with \(y_n=0\), recover
\(y_i=y_{i+1}+r_{i+1}-z_i\), then \(x_i=y_i+r_i\).
Rotation reverses both profile words.  A board is \emph{zero-free} if
\(X\cap Y=\varnothing\), and \emph{dense} if \(x_n>y_1\).
These are different conditions.  Ordinary degree assigns degree one to
each \(h_j\), \(j>0\); weighted degree assigns degree \(j\).
The square weight of a monomial is the sum of the squares of its generator
indices, with multiplicity.  Units contribute zero to all three readings.

\begin{lemma}\label{lem:readings}
For a connected diagram with full lists \(X,Y\), the maximum ordinary
degree of \(s_D\) is \(n\).  Its unique minimum-square-weight monomial is \(\prod_i h_{r_i}\),
with coefficient one.  The largest generator occurring is \(h_N\), and
\begin{equation}\label{eq:maximal-cofactor}
 (-1)^{n-1}[h_N]s_D
 =\det(h_{x_{i+1}-y_j})_{1\le i,j<n}.
\end{equation}
The minimum-square monomial of this cofactor reads the multiset of the
\(z_i\), padded with zeros to length \(n-1\).  The minimum ordinary degree
of \(s_D\) is \(n-|X\cap Y|\).
\end{lemma}

\begin{proof}
The maximum degree, minimum-square monomial, largest generator, and
cofactor readings are the polynomial-readings lemma of
\cite{Shields2026Cofactor}, applied to the full lists.  For the minimum
degree put \(Z=X\cap Y\), \(A=X\setminus Z\), and
\(B=Y\setminus Z\).  The exact common-beta path expansion in
\cite{Shields2026Cofactor} gives
\[
 s_D(th_1,th_2,\ldots)
 =\epsilon t^{n-|Z|}\det(h_{a_i-b_j})
   +\text{terms of higher degree in }t,
\]
where \(\epsilon\in\{1,-1\}\) is the regrouping sign.  The residual
strict lists have positive sorted diagonal, so its determinant is
nonzero by strict rearrangement.  The minimum degree is therefore
\(n-|Z|\).
\end{proof}

The polynomial readings are also consequences of the overlap theory in
\cite{Reiner2007Coincidences}.  We use the corrected irreducibility theorem
of Barekat, Reiner and van Willigenburg
\cite{Barekat2009CoincidencesCorrigendum}: the Schur polynomial of a
nonempty connected skew diagram is irreducible in the complete-function
algebra.  A disconnected diagram contributes the product of its connected
component polynomials.  Their coefficient-one minimum-square monomials
fix the scalar associates.  Thus factorization identifies the component
polynomials, not their positions or orientations.

\subsection{The invariant and the pattern}

\begin{lemma}[Negative-count invariance]\label{lem:five-mate-closure}
Connected Schur-equal diagrams have the same row count and the same
full-board negative count.  In particular, every connected mate of a
five-row five-negative board is again a five-row five-negative board.  
\end{lemma}

\begin{proof}
For consecutive rows \(j,\ldots,j+k-1\), put
\(o_{j,k}=\max(0,X_{j+k-1}-Y_j+k-1)\), their intersection size.  
For \(k\ge2\), the index \(X_{j+k-1}-Y_j\) is negative exactly
when \(o_{j,k}\le k-2\).  The diagonal and upper-triangular indices
are positive, so
\[
 \nu(D)=\sum_{k=2}^n
 \#\{j:1\le j\le n-k+1,\ o_{j,k}\le k-2\}.  
\]
The row count is read by \Cref{lem:readings}.  The unordered row-overlap
partitions are Schur invariants \cite{Reiner2007Coincidences}; padding
the \(k\)-th partition to length \(n-k+1\) recovers its zero parts.  
Hence every summand is invariant.  At a common beta the corresponding
overlap is \(k-1\), so that zero index is correctly excluded.  
\end{proof}

\begin{lemma}\label{lem:negative-partition}
Let \(D\) have \(n\) nonempty rows and full beta lists \(X,Y\).
Put \(c_i=\#\{j:x_i-y_j<0\}\).  Its negative sites are closed
southwest, and \(c_1\le\cdots\le c_n\).  Reading these counts from
bottom to top and omitting zero parts gives a partition \(\rho(D)\),
with
\[
 |\rho(D)|=\nu(D),\qquad \rho(D^*)=\rho(D)' .
\]
If \(D\) is connected, then \(c_i\le\max(0,i-2)\).
\end{lemma}

\begin{proof}
If \(x_i<y_j\), strict decrease of both lists gives
\(x_{i'}<y_{j'}\) whenever \(i'\ge i\) and \(j'\le j\).
Thus each negative row is an initial segment and these segments increase
downward.  In a connected board the first subdiagonal is nonnegative,
which gives the stated bound.

For normalized lists, rotation sends \((X,Y)\) to
\[
 \bigl((N-y_n,\ldots,N-y_1),(N-x_n,\ldots,N-x_1)\bigr).
\]
Its index at \((i,j)\) is the original index at
\((n+1-j,n+1-i)\).  In coordinates measured upward and rightward from
the southwest corner, this map interchanges the two coordinates.
It therefore conjugates the negative partition.
\end{proof}

\begin{example}\label{ex:negative-shape-change}
The thirteen-cell profiles
\[
 F=(3,4,3,2,1;1,2,1,0),\qquad
 H=(4,3,3,2,1;2,1,1,0)
\]
have \(s_F=s_H\) by \Cref{thm:five-base}, at \(b=q=1,g=2\)
in the first family of \eqref{eq:five-families}.  Their negative
partitions are
\[
 \rho(F)=(3,1,1),\qquad \rho(H)=(3,2).
\]
The former is self-conjugate, so the two partitions are neither equal nor
conjugate.  Negative partition is consequently not an invariant of a
connected Schur class, even after identifying conjugate partitions.
\end{example}

\subsection{Realizability and minimum row count}\label{sec:negative-invariants}

\begin{theorem}\label{thm:negative-realizability}
A nonempty partition \(\rho\) is the negative partition of a connected skew diagram with exactly \(n\) nonempty rows if and only if
\begin{equation}\label{eq:realizability}
 \rho_i\le n-i-1\qquad(1\le i\le\ell(\rho)).
\end{equation}
Consequently \eqref{eq:minimum-intro} holds, and every partition is realizable when the row count is unrestricted.
\end{theorem}

\begin{proof}
Connectedness gives \(x_j\ge y_{j-1}\) for \(j\ge2\).  Thus row \(j\) has at most \(j-2\) negative entries.  Reading the rows from below gives \eqref{eq:realizability}.

Conversely, pad \(\rho\) with zeros to length \(n\), put \(t_i=\rho_{n+1-i}\) and \(M=n+1\), and set
\begin{equation}\label{eq:realizing-lists}
 X_i=M(n-t_i)-i,\qquad Y_j=M(n-j)
 \qquad(1\le i,j\le n).
\end{equation}
The \(t_i\) are nondecreasing, so the \(X_i\) are strictly decreasing; the \(Y_j\) are strictly decreasing and \(Y_n=0\).  The board index is
\[
 X_i-Y_j=M(j-t_i)-i.
\]
Since \(1\le i<M\), it is negative exactly when \(j\le t_i\).  The prescribed pattern is therefore obtained.  Also \(t_1=0\) and \(t_i\le i-2\) for \(i\ge2\).  The diagonal is positive, and
\[
 X_i-Y_{i-1}=M(i-1-t_i)-i\ge M-i>0
 \qquad(i\ge2).
\]
The two strict lists consequently define nonempty rows of a connected skew diagram.  The least admissible \(n\) in \eqref{eq:realizability} is exactly \eqref{eq:minimum-intro}.
\end{proof}

For \(1\le k\le n\), the overlap of rows \(j,\ldots,j+k-1\) is
\begin{equation}\label{eq:row-overlap}
 o_k(j)=\max(0,x_{j+k-1}-y_j+k-1).
\end{equation}
The multiset of these overlaps, sorted and padded with zeros to length \(n-k+1\), is a row-overlap partition.  These partitions are Schur invariants \cite{Reiner2007Coincidences}.

\begin{proposition}\label{prop:minimal-row-invariant}
The quantity \(n_{\min}(\rho(D))\) is constant on a connected Schur class.  If \(\rho(D)\ne\varnothing\), let \(k_0\) be the first order \(k\ge2\) for which the padded overlap partition contains an entry at most \(k-2\).  Then
\begin{equation}\label{eq:overlap-minimum}
 n_{\min}(\rho(D))=n+3-k_0.
\end{equation}
\end{proposition}

\begin{proof}
Row count and negative count are invariant by \Cref{lem:five-mate-closure}; in particular the empty-pattern case stays empty.  By \eqref{eq:row-overlap}, an index in position \((j+k-1,j)\) is negative exactly when \(o_k(j)\le k-2\).  If \(t_i\) is the number of negative entries in row \(i\), the first such order is
\[
 k_0=1+\min_{t_i>0}(i-t_i).
\]
Substituting \(\rho_j=t_{n+1-j}\) into \eqref{eq:minimum-intro} gives \eqref{eq:overlap-minimum}.  Both \(n\) and the padded overlap partitions are invariant.  This argument includes zero indices and does not recover the ordering of the negative row counts within the board.
\end{proof}

\section{Boundary coefficients and propagation}\label{sec:boundary-tools}

\subsection{Faithful passage to zero-free boards}

\begin{lemma}\label{lem:homogenization}
The map
\begin{equation}\label{eq:gamma}
 \Gamma(X,Y)=(2X+1,2Y)
\end{equation}
is injective on translation classes, commutes with rotation, preserves
every negative position, and sends every connected board to a connected
zero-free board of the same order.  For \(s_D=\sum_\alpha c_\alpha h_\alpha\),
\begin{equation}\label{eq:homogenization}
 s_{\Gamma D}=\sum_\alpha c_\alpha h_1^{n-\ell(\alpha)}
                         \prod_{j\in\alpha}h_{2j+1}.
\end{equation}
Consequently any rotation-rigidity statement about a specified negative
pattern descends from zero-free boards to full boards.
\end{lemma}

\begin{proof}
An index \(j\) becomes \(2j+1\), which is negative exactly when \(j<0\).
The two lists have opposite parity.  Connectedness is preserved, and the
profile becomes \((2r+1;2z+1)\).  This also proves injectivity and rotation
covariance.  In each surviving determinant term the omitted unit factors
number \(n-\ell(\alpha)\), proving \eqref{eq:homogenization}.
Equality supplies the common \(n\) by \Cref{lem:readings}, so it transports
through this formula.  No claim about composition-irreducibility is needed.
\end{proof}

By \Cref{lem:readings}, a connected mate of a zero-free source
is zero-free: its minimum and maximum ordinary degrees both equal
its row count.  The maximal cofactor of a connected zero-free board
has positive rows.  In the component arguments below, all cofactor
components refer to this zero-free setting.

\begin{lemma}\label{lem:component-budget}
Suppose the maximal cofactor has connected components of row counts \(n_1,\ldots,n_r\) and internal negative counts \(\nu_1,\ldots,\nu_r\).  Then
\begin{equation}\label{eq:component-budget}
 \nu(D)=\sum_{i<j}n_i n_j+\sum_i\nu_i.
\end{equation}
\end{lemma}

\begin{proof}
In the ordered cofactor board every block below a diagonal component block is entirely negative, and every block above it is positive.  The diagonal blocks contribute their internal negative counts.  The deleted first row and last column of the original board have no negative indices.  This gives \eqref{eq:component-budget}.
\end{proof}

For fixed \(m=\sum_i n_i\) and \(r\), the cross cost is \((m^2-\sum_i n_i^2)/2\).  Moving one row from a smaller nonsingleton component to a largest component increases the sum of squares.  Hence the least cross cost occurs at \((m-r+1,1,\ldots,1)\) and equals
\[
 (r-1)m-\frac{r(r-1)}2.
\]

\subsection{Connected children}

For a normalized diagram \(R=(U,V)\) with nonempty rows and span \(S=U_1\), define
its positive border and its two endpoint extensions by
\begin{equation}\label{eq:border-definition}
 \B_{u,v}R=((S+u,U);(V,-v)),\qquad u,v>0.
\end{equation}
\[
 R^{\uparrow t}=((U_1+t,U_2,\ldots);V),\qquad
 R^{\downarrow t}=(U;(V_1,\ldots,V_m-t)),\qquad t>0.
\]
Lists are translated together when normalization is required.

\begin{lemma}\label{lem:border}
Let \(R\) be connected and zero-free with \(m\ge2\) rows and span \(S\).
For \(P=\B_{u,v}R\), put \(N=S+u+v\), \(q=\min(u,v)\), \(M=N-q\).
Then \([h_N]s_P=(-1)^m s_R\).  After setting \(h_N=0\), the largest
remaining label is \(M\).  If \(u<v\), its coefficient is
\((-1)^{m+1}s_{R^{\uparrow q}}\); if \(v<u\), it is
\((-1)^{m+1}s_{R^{\downarrow q}}\).  If \(u=v\),
\begin{equation}\label{eq:border-square}
 [h_M^2]s_P=\det(h_{U_i-V_j})_{2\le i\le m,\ 1\le j<m}\ne0.
\end{equation}
Thus the full polynomial and the cofactor class determine \(\{u,v\}\).
If the literal cofactor has been aligned and its opposite endpoint
extensions can be equal only when \(R=R^*\), equal borders are rotations.
\end{lemma}

\begin{proof}
The largest noncorner entries are \(N-u\) at \((2,m+1)\) and \(N-v\)
at \((1,m)\).  At unequal margins the larger is private and gives the
stated connected extension.  Its span \(S+q\) is smaller than \(M\).
At equal margins both sites contribute, with positive joint sign; their
complementary sorted diagonal is positive.  Every coefficient label in
\eqref{eq:border-square} is below \(M\).  These bounds justify extraction
after the corner term has been removed.  The span determines \(u+v\), and
\(M\) determines their minimum.  A nonliteral swap equates the two
endpoint extensions.  If \(R=R^*\), it rotates the parent.
\end{proof}

\begin{lemma}\label{lem:above-minimum-border}
If \(\rho(D)\ne\varnothing\) and \(n>n_{\min}(\rho(D))\), the zero-free maximal cofactor is connected and has negative partition \(\rho(D)\).  If \(n=n_{\min}(\rho(D))\), it is disconnected.  Positive endpoint extensions preserve the negative partition.
\end{lemma}

\begin{proof}
Deleting the first row and last column removes no negatives.  Above the minimum,
\[
 \rho_i+i\le n_{\min}(\rho)-1\le n-2,
\]
so every first-subdiagonal entry of the maximal cofactor is positive.  At the minimum some \(\rho_i+i=n-1\), placing a negative entry on that subdiagonal and separating components.  The two endpoint changes affect only the first row or last column, which are positive, so preserve the negative positions.
\end{proof}

\begin{lemma}\label{lem:small-rows}
Every connected diagram with at most three nonempty rows is rotation-rigid.
Every connected diagram whose full board is nonnegative is rotation-rigid,
in arbitrary row count.
\end{lemma}

\begin{proof}
By \Cref{lem:homogenization}, it suffices to treat zero-free boards;
\Cref{lem:readings} supplies the mate's row count and the two multisets.
One row is immediate.  Two rows have polynomial \(h_a h_b-h_{a+b-z}h_z\), recovering
\(\{a,b\}\) and \(z\).  In three rows write the profile as
\((a,b,c;s,t)\).  Its corner cofactor is
\(h_s h_t-h_bh_{s+t-b}\).
If \(s+t-b\ge0\), its top label recovers \(b\), and its diagonal recovers
\(\{s,t\}\).  After orientation only interchange of \(a,c\) remains.
For \(a\ne c\), \(s\ne t\), interchange the comparisons and rotate to
have \(a>c,s>t\).  The largest label left in their difference is
\(a+b-t\), with coefficient
\[
 h_t h_c-h_{s+t-b}h_{c+b-s}\ne0;
\]
the latter two indices are respectively below \(t\) and above \(c\).
If \(a=c\) nothing changes; if \(s=t\) the interchange is rotation.

If \(s+t-b<0\), first take \(s<t\).  After removing the row monomial
and full-run term, the two negative terms are
\[
 h_{a+b-s}h_s h_c,\qquad h_a h_{b+c-t}h_t.
\]
Only the first carries the smallest index \(s\).  Its other factors
identify \(c\) among the row multiset, since \(a+b-s>a,b\); the second
then identifies \(a\).  When \(s=t\), uniform shortening reduces to a
three-row ribbon, which is rotation-rigid by
\cite{Billera2006DecomposableCompositions}.  Shortening is performed on
the homogeneous determinant with its unit variable retained, so it is
injective.

The arbitrary-order nonnegative-board assertion is the
nonnegative-index rigidity theorem of \cite{Shields2026Cofactor}.
That theorem compares against arbitrary connected mates and allows
common beta entries.
\end{proof}

\subsection{A singleton beside a connected component}

For a normalized connected zero-free component \(Q=(C,D)\), with
\(C_1=S,D_m=0\), and \(s,u,v,g>0\), put \(a=s+g\), \(z=v+S+g\), and
\begin{equation}\label{eq:singleton-context}
 \begin{aligned}
 X&=(z+s+u,z+s,v+C_1,\ldots,v+C_m),\\
 Y&=(z,v+D_1,\ldots,v+D_m,0).
 \end{aligned}
\end{equation}
The maximal cofactor is a singleton of length \(s\) followed by \(Q\).
The separation \(g\) is a gap between these cofactor components, not a
parent overlap.  Write \(\Phi_Q(w,t)=s_{\B_{w,t}Q}\).  First-column
expansion gives the full identity
\begin{equation}\label{eq:singleton-polynomial}
 P_Q=h_{s+u}\Phi_Q(a,v)-h_s\Phi_Q(a+u,v).
\end{equation}

For any normalized zero-free diagram \(Q=(C,D)\) with \(m\ge2\) nonempty
rows, put \(S=C_1\), \(D_m=0\), and
\(D_Q(t)=s_{Q^{\uparrow t}}-s_{Q^{\downarrow t}}\), for \(t>0\).
Let \(M_{ij}\) be its unsigned minors.  Canceling the common moving
corner in the top-row and last-column Laplace expansions gives
\begin{equation}\label{eq:levels}
 D_Q(t)=\sum_{j=1}^{m-1}(-1)^{1+j}h_{S-D_j+t}M_{1j}
       -\sum_{i=2}^{m}(-1)^{i+m}h_{C_i+t}M_{im}.
\end{equation}
Every label in a first-row cofactor is at most \(C_2\), and every label
in a last-column cofactor is at most \(S-D_{m-1}\).

\begin{lemma}\label{lem:endpoint-top}
Suppose \(Q\ne Q^*\) is connected and zero-free, and either is dense of
order \(m\ge2\), or has exactly one negative and order three or four.
For every \(t>0\),
\begin{equation}\label{eq:endpoint-top}
 \max\operatorname{label}D_Q(t)=\kappa+t,
 \qquad [h_{\kappa+t}]D_Q(t)=c_Q\ne0,\quad \kappa=\max(C_2,S-D_{m-1}),
\end{equation}
where \(0<\kappa<S\) and \(c_Q\) is independent of \(t\).
Also \(D_Q(t)=0\) for symmetric \(Q\).
\end{lemma}

\begin{proof}
The cofactor bounds in \eqref{eq:levels} show that only the top level can
carry \(h_{\kappa+t}\).  Off the tie \(C_2=S-D_{m-1}\), there
is one nonzero minor.  At the tie, the coefficient is the difference of
the opposite endpoint extensions of the maximal cofactor, up to sign.
For a dense component of order at least three that child is connected and
dense; for a four-row one-negative component it is connected with three
rows.  Their rigidity, proved in \Cref{lem:small-rows}, makes the difference
nonzero unless the child is symmetric.  Equal exterior gaps would then
make \(Q\) symmetric.  In order two the same conclusion follows directly.

For a three-row one-negative component write its strict endpoint chain as
\[
 0<D_2<C_3<D_1<C_2<S,
\]
with consecutive gaps \((\alpha,\pi,\theta,\rho,\gamma)\).
At the top tie \(\gamma=\alpha\), the two minors have difference
\[
 h_\rho h_{\alpha+\pi}-h_{\rho+\gamma}h_\pi,
\]
which vanishes exactly when \(\rho=\pi\).  Together with
\(\gamma=\alpha\), that is exactly symmetry.  Off the tie only the
appropriate single minor is used, not their difference.  All cofactor
labels are at most \(\kappa\), completing the claimed extraction.
\end{proof}

\begin{lemma}\label{lem:singleton-lift}
Let \(Q\) be connected and zero-free of order at least two.  Assume that
\(Q\) and its positive endpoint extensions are rotation-rigid against
arbitrary connected mates, and that nonsymmetric \(Q\) satisfies the
nonzero, shift-independent leading-coefficient assertion of
\eqref{eq:endpoint-top}.  Every connected zero-free parent whose maximal
cofactor consists of a singleton and \(Q\) is rotation-rigid.
\end{lemma}

\begin{proof}
By \Cref{lem:readings}, an arbitrary mate is zero-free of order \(m+2\).
Its corner coefficient is \(\varepsilon h_s s_Q\), where
\(\varepsilon=(-1)^{m+1}\).  Unique factorization separates the singleton
from the degree-\(m\) irreducible factor.  Component rigidity identifies
that factor with \(Q\) or \(Q^*\).  Rotate the mate to put the singleton
first; it has exactly the form \eqref{eq:singleton-context}.

Put \(N=S+a+u+v\), \(q=\min(u,v)\), \(M=N-q\).  Expansion gives
\begin{equation}\label{eq:singleton-marker}
 -\varepsilon[h_M]P_Q=
 \begin{cases}
 h_{s+q}s_Q,&u<v,\\
 h_s s_{Q^{\downarrow q}},&v<u,\\
 h_{s+q}s_Q+h_s s_{Q^{\downarrow q}},&u=v.
 \end{cases}
\end{equation}
All coefficient labels are below \(M\).  Divisibility by
\(h_s,h_{s+q}\) distinguishes the three cases: respectively only the
second, only the first, or neither divides.  The tied sum is nonzero by
Schur positivity.  The connected degree-\(m\) factors cannot have a
linear factor.  When \(v\le u\), the endpoint rigidity of \(Q\) fixes
its orientation.  At a tie \(N\) fixes \(g\); for \(v<u\), the second
terms of \eqref{eq:singleton-polynomial} agree, and unique factorization
of the first terms fixes \(u\), then \(g\).

Suppose \(u<v\).  The polynomial fixes \(u\) and \(a+v\).
The entire first two terms of last-column expansion are the common part
\[
 J=\varepsilon s_Q(h_sh_N-h_{s+u}h_{N-u}).
\]
If two candidates have \(v_-<v_+\), subtract \(J\) and set
\(K=\max(N-v_-,S+v_+)\).  For the difference ``larger bottom margin minus
smaller bottom margin,'' its coefficient at \(K\) is
\begin{equation}\label{eq:margin-separator}
 \varepsilon\begin{cases}
 h_s s_{Q_-^{\downarrow v_-}},&N-v_->S+v_+,\\
 s_T,&N-v_-<S+v_+,\\
 h_s s_{Q_-^{\downarrow v_-}}+s_T,&N-v_-=S+v_+.
 \end{cases}
\end{equation}
Here \(T\) deletes the first component row and last column from the
larger-margin parent.  Its subdiagonal labels are
\(s,C_2-D_1,\ldots,C_m-D_{m-1}>0\), so it is connected.
Its span is \(N-v_+<K\); the product's labels are at most
\(\max(s,S+v_-)<K\); and \(K>\max(S,s+u)\).
The tied coefficient is a same-sign sum of nonzero Schur-positive
functions.  Thus \(v\), and then \(g\), are fixed.

Only reflection of a nonsymmetric component at these unchanged margins
remains.  Write
\[
 \mathcal D(w)=\Phi_Q(w,v)-\Phi_Q(v,w).
\]
Rotation identifies the second term with \(\Phi_{Q^*}(w,v)\).
This difference is zero at \(w=v\); otherwise its largest label is
\(S+\max(w,v)\), with a signed coefficient
\(D_Q(\min(w,v))\).  The full difference is
\[
 h_{s+u}\mathcal D(a)-h_s\mathcal D(a+u).
\]
A zero summand leaves the other nonzero.  For \(v<a+u\), the second
nonzero summand has strictly larger top label.  For \(v>a+u\), both tops
are \(S+v\), and their coefficient, up to sign, is
\begin{equation}\label{eq:double-top}
 h_{s+u}D_Q(a)-h_sD_Q(a+u).
\end{equation}
The second term has label \(\kappa+a+u\), greater than both
\(\kappa+a\) and \(s+u\), since \(a=s+g>s\).
It cannot cancel.  This includes equal endpoint rows of \(Q\), unchanged
margins, and a top attained twice in both border differences.
\end{proof}

\begin{corollary}\label{lem:singleton-rigidity}
Every parent in \eqref{eq:singleton-context} is rotation-rigid if \(Q\)
is dense of order \(m\ge2\), or has exactly one negative and order three
or four.  Its parameters \(s,u,v,g\) are positive integers.
\end{corollary}

\begin{proof}
Dense components and three-row components are rigid by
\Cref{lem:small-rows}.  A one-negative four-row component has a connected
three-row maximal cofactor, so \Cref{lem:small-rows,lem:border} give its
rigidity as well.  The same arguments apply to its endpoint extensions,
which preserve negative positions.  The leading-coefficient hypothesis
is \Cref{lem:endpoint-top}; apply \Cref{lem:singleton-lift}.
\end{proof}

\section{Classification through four negatives}\label{sec:low}

\subsection{Four rows and the two-negative lift}

The four-row classification needed below is stated first.  Its proof is in
Appendix~\ref{app:four}; it is an algebraic case calculation, not a census.

\begin{proposition}\label{prop:four-rows}
For a connected diagram with exactly four nonempty rows, every connected
mate is its rotation, except for the complete classes
\(\{F,F^*,H,H^*\}\) with profiles
\begin{equation}\label{eq:four-family}
 \begin{aligned}
 F&=(w,b,b,b+w-c;c,c,w),\\
 H&=(w,b,b+w-c,b;c,w,c),
 \end{aligned}
 \qquad 0\le c<w,\quad b>c+w.
\end{equation}
All four members are distinct and have exactly three negative indices.
\end{proposition}

At \((c,w,b)=(0,1,2)\), the displayed family gives the eight-cell pair
\[
 (1,2,2,3;0,0,1),\qquad (1,2,3,2;0,1,0).
\]

\begin{theorem}\label{thm:two-negatives}
A connected diagram with at most two negative indices in its full row
board is rotation-rigid, in every row count and with common betas allowed.
\end{theorem}

\begin{proof}
Work zero-free by \Cref{lem:homogenization}.  The cases through four rows
are \Cref{lem:small-rows,prop:four-rows}.  Above four rows the maximal
cofactor has \(m=n-1\ge4\) positive rows.  A split after row \(k\) would
require \(k(m-k)\ge m-1\ge3\) negatives.  Thus it is connected and has
at most two negatives; its endpoint extensions retain those same negative
positions.  Induction and \Cref{lem:border} apply, including to the
arbitrary mate after factorization of its maximal cofactor.
\end{proof}

\subsection{Columns and three corners}

\begin{theorem}\label{thm:column}
Let \(D\) be connected with \(n\ge2\) nonempty rows.  If
\(x_n\ge y_2\), or if \(x_{n-1}\ge y_1\), every connected mate is a
rotation.  Common betas are allowed, and there is no bound on the number
of negatives.
\end{theorem}

\begin{proof}
Rotate to use the first condition, and homogenize.  Through four rows
there are at most two negatives.  A connected maximal cofactor remains in
the same single-column class, as do its endpoint extensions, so induction
and \Cref{lem:border} apply.  If disconnected, only its first subdiagonal
can separate: the cofactor is a singleton followed by a dense block.
\Cref{lem:singleton-rigidity} applies.  Descend by
\Cref{lem:homogenization}.
\end{proof}

\begin{lemma}\label{lem:four-border}
Let \(R_1,R_2\) be the two displayed four-row profiles in
\eqref{eq:four-family}, with \(1\le c<w\), \(b>c+w\).
No positive border of \(R_1\) is Schur-equal to a positive border of
\(R_2\), including either orientation of each cofactor.
\end{lemma}

\begin{proof}
Put \(t=b-c\), \(d=w-c\), and \(S=3t+w\).
Both cofactor spans are \(S\).  The border readings give common unordered
margins.  In fixed orientations the parent rows are
\[
 \begin{aligned}
 &(w+u,b+d,2b-c,2b-c,b+d+v),\\
 &(w+u',b+d,2b-c,2b-c+d,b+v').
 \end{aligned}
\]
Canceling common entries in their multisets gives: at the same margins,
\(v=t\); at interchanged margins, either \(u=v=t\) or
\((u,v)=(2t,2t-d)\).  These exhaust the two possible matches for the
fixed entry \(2b-c+d\).

For \(v=t,u<t\), the marker compares top extensions of the two children.
Every row of these extensions exceeds \(w\), whereas an exceptional
four-row diagram with overlap multiset \(\{c,c,w\}\) has a row \(w\).
They cannot be rotations either, since their overlap words differ even
up to reversal.  For \(v=t,u>t\), the two bottom extensions have different
row multisets: after canceling \(w,b\), one contains \(b\) and the other
only larger values.  At the unequal swapped margins the marker compares
a bottom extension still containing \(w\) with a top extension whose
rows all exceed it.

At \(u=v=t\), put \(M=4t+w\), \(T=3b-2c\).
The eight-term four-row expansion of the endpoint cofactors gives
\[
 \begin{split}
 [h_M](s_{\B_{t,t}R_1}-s_{\B_{t,t}R_2})=-h_w\bigl(&
 (h_ch_{b+d}-h_wh_b)h_T\\
 &+(h_b^2-h_ch_{2b-c})h_{2b-c+d}\\
 &+h_wh_{2b-c}^2-h_bh_{b+d}h_{2b-c}\bigr).
 \end{split}
\]
The square coefficients cancel.  Since
\(c<w<b<b+d<2b-c<2b-c+d<T<M\), the double coefficient is
\begin{equation}\label{eq:four-border-separator}
 [h_Mh_T](s_{\B_{t,t}R_1}-s_{\B_{t,t}R_2})
 =-h_w(h_ch_{b+d}-h_wh_b)\ne0.
\end{equation}
Rotating whole parents handles either orientation of the children.
\end{proof}

\begin{proposition}\label{prop:three-corners}
Let \(D\) be connected with \(n\ge5\) rows and
\[
 x_{n-2}\ge y_1,\qquad x_{n-1}\ge y_2,\qquad x_n\ge y_3.
\]
Then \(D\) is rotation-rigid.  Common betas are allowed.
\end{proposition}

\begin{proof}
Homogenize first.  The maximal cofactor is connected, zero-free and has
\(n-1\) positive rows: its subdiagonals are \(x_{i+2}-y_i>0\).
For five rows, \Cref{prop:four-rows} identifies its complete class, and
\Cref{lem:four-border} excludes crossing between the exceptional orbits.
After aligning the literal cofactor, a margin swap would equate its
opposite endpoint extensions.  These have the same ordered overlap word,
whereas the two exceptional orbits in \eqref{eq:four-family} have
inequivalent overlap words.  Thus the extensions are rotations, forcing a
symmetric child.  For \(n\ge6\), the child retains the three displayed
inequalities one order down; so do its endpoint extensions.  Induction,
\Cref{lem:border}, and \Cref{lem:homogenization} finish the proof.
\end{proof}

\subsection{Square and hook}

A two-row block \(Q=\langle A,B,d\rangle\) has beta lists
\((A+B-d,B;B-d,0)\) and polynomial
\[
 f_Q=h_Ah_B-h_{A+B-d}h_d,\qquad 0\le d<\min(A,B).
\]
Its rotation exchanges \(A,B\).  Superscripts \(\uparrow,\downarrow\)
continue to denote endpoint extension.

\begin{theorem}[Square rigidity]\label{thm:square}
Let \(F\) be connected with \(n\ge5\) rows and negative positions exactly
\[
 \{(n-1,1),(n-1,2),(n,1),(n,2)\}.
\]
Then \(F\) is rotation-rigid.  Common betas are allowed.
\end{theorem}

\begin{proof}
First take five rows.  Its maximal cofactor consists of two two-row blocks
\(Q=\langle A,B,d\rangle\), \(R=\langle C,D,e\rangle\).
Its maximum ordinary degree is four, even when a block has a unit.
Hence an equal mate's cofactor has four positive rows.  Unique
factorization forces the same two blocks up to order and rotation.
This recovers the square pattern on the mate without assuming it.

Now work zero-free, so \(d,e>0\).  Set
\(q_1=A+B-d\), \(P=C+D-e\), \(s_1=D-e\), \(\gamma=P+g\),
\(W=v+\gamma\), with \(u,v,g>0\).  All sources have the unique form
\begin{equation}\label{eq:square-lists}
 \begin{aligned}
 X&=(W+q_1+u,W+q_1,W+B,v+P,v+D),\\
 Y&=(W+B-d,W,v+s_1,v,0).
 \end{aligned}
\end{equation}
Its row word is \((A+u,q_1,B+C+g,P,D+v)\), its overlap word is
\((A,B,C,D)\), and its span is \(N=q_1+P+u+v+g\).
Rotation sends \((Q,R,u,v,g)\) to \((R^*,Q^*,v,u,g)\).

Put
\[
 \Phi(t)=\det\begin{pmatrix}
 h_{A+u}&h_{q_1+u}&h_{t+q_1+u}\\
 h_A&h_{q_1}&h_{t+q_1}\\
 h_d&h_B&h_{t+B}
 \end{pmatrix}.
\]
For \(t>0\) this is a connected three-row Schur polynomial.  If
\(\sigma=s_{Q^{\uparrow u}}\),
\(\tau=s_{\langle A+u,q_1,A\rangle}\), expansion gives
\begin{equation}\label{eq:square-phi}
 \Phi(t)=h_{t+q_1+u}f_Q-h_{t+q_1}\sigma+h_{t+B}\tau.
\end{equation}
The three nonzero bottom-row-pair minors give
\begin{equation}\label{eq:square-star}
 s_F=f_R\Phi(W)-s_{R^{\downarrow v}}\Phi(\gamma)
       +s_{\langle P,D+v,D\rangle}\Phi(\gamma-s_1).
\end{equation}
All three arguments are positive.  The corner coefficient is \(f_Qf_R\).
The next label \(M=N-\min(u,v)\) has coefficient
\begin{equation}\label{eq:square-marker}
 [h_M]s_F=-\begin{cases}
 s_{Q^{\uparrow u}}f_R,&u<v,\\
 f_Qs_{R^{\downarrow v}},&v<u,\\
 s_{Q^{\uparrow u}}f_R+f_Qs_{R^{\downarrow v}},&u=v.
 \end{cases}
\end{equation}
Every coefficient label is below \(M\).

Orient both parents so \(u\le v\), \(u'\le v'\).  The marker fixes
\(u=u'\).  With strict margins, unique factorization identifies the
unextended factor and then the orientation of \(Q\); consequently \(Q\)
and \(W=N-q_1-u\) agree.  This remains true when \(f_Q=f_R\), by
canceling that factor.  Strict and tied margins cannot mix.  If the two
factors differ, divisibility by the old factors distinguishes them; if
they agree, cancellation leaves one endpoint extension against a sum of
two, whose coefficients at \(h_{q_1+u}h_d\) are \(-1\) and \(-2\).

At tied margins the order of unequal factors can be aligned by global
rotation.  An orientation change has marker difference
\(E_Qf_R+f_QE_R\), where
\[
 E_Q=h_{A+u}h_B-h_{B+u}h_A,\qquad
 E_R=h_Ch_{D+u}-h_Dh_{C+u},
\]
and either summand may be absent.  A nonzero \(E_Q\) is not divisible
by \(f_Q\): both have ordinary degree two and different weighted degrees.
Reduction modulo the distinct irreducible factors forces each orientation
to agree.  For equal factors put \(p=h_{A+u}h_B\), \(r=h_Ah_{B+u}\).
The four marker values, after their common term is removed, are
\(p+r,p+r,2p,2r\).  Their only nontrivial coincidence is the global
rotation pair, unless \(A=B\), when the block is already symmetric.
The span then fixes \(g\).

It remains to treat strict margins, with \(Q,u,W\) fixed.  Subtract the
common first term of \eqref{eq:square-star}.  If \(v'<v\), the largest
possible label in the remaining difference is
\(K=\max(N-v',P+v)\).  When the first is larger, its private coefficient
is \(-s_{R'^{\downarrow v'}}f_Q\ne0\).  When the second is larger,
vanishing would require
\(h_e\Phi(\gamma)=h_D\Phi(\gamma-s_1)\), impossible by unique
factorization of the connected cubics and \(e<D\).

At equality of those two labels, vanishing forces
\[
 f_Q\mid h_D\Phi(\gamma-s_1)-h_e\Phi(\gamma).
\]
The two column syzygies of the displayed \(3\)-by-\(2\) matrix give
\(h_A\sigma\equiv h_d\tau\pmod{f_Q}\), with
\(\tau\not\equiv0\).  More generally,
\(f_Q\mid h_{c_1}\Phi(t_1)-h_{c_2}\Phi(t_2)\) implies
\begin{equation}\label{eq:square-residue}
 \begin{split}
 f_Q\mid {}&h_dh_{c_2}h_{t_2+q_1}-h_dh_{c_1}h_{t_1+q_1}\\
           &-h_Ah_{c_2}h_{t_2+B}+h_Ah_{c_1}h_{t_1+B}.
 \end{split}
\end{equation}
Indeed multiply by \(h_A\) and cancel \(\tau\) in the integral domain
\(\Q[h]/(f_Q)\).  For \((c_1,t_1)=(D,C+g)\),
\((c_2,t_2)=(e,P+g)\), the unique top label on the right is
\(P+g+q_1\), with coefficient \(h_dh_e\).  This label exceeds all
labels of \(f_Q\).  A cubic divisible by the quadratic \(f_Q\) would
have at that label a scalar multiple of its two-term polynomial, not a
single nonzero monomial.  Thus the tie is impossible, and \(v'=v\).

Only \(R'=R^*\) remains.  If \(C\ne D\), comparing
\(P+v\) with \(N-v\) gives respectively
\(h_C\Phi(D+g)=h_D\Phi(C+g)\), or
\((h_Dh_{C+v}-h_Ch_{D+v})f_Q=0\); both are impossible.
At the tie use \eqref{eq:square-residue} with
\((c_1,t_1)=(D,C+g)\), \((c_2,t_2)=(C,D+g)\).
Its unique top coefficient is \(h_dh_C\) or \(-h_dh_D\), giving the
same contradiction.  Thus \(R'=R\) and \(g'=g\).

For \(n\ge6\), the maximal cofactor is connected: a negative at its
subdiagonal would require simultaneously \(i\le2\) and
\(i+2\ge n-1\), which is impossible.  It is the square one order down;
endpoint extensions preserve that exact negative set.  Induction and
\Cref{lem:border} prove rigidity.  Homogenization supplies common betas.
\end{proof}

\begin{remark}\label{rem:square-degree}
The full-board degree assertion in the proof is a maximum, not homogeneity.
For \(X=(13,12,9,6,5)\), \(Y=(9,7,4,2,0)\), the maximal cofactor is
\[
 (h_3h_2-h_5)(h_2h_3-h_4h_1),
\]
with ordinary degrees three and four.  Its maximum still recovers four
nonempty cofactor rows.
\end{remark}

\begin{proposition}\label{prop:hook}
A connected diagram with \(n\ge5\) rows and negative set
\[
 \{(n-2,1),(n-1,1),(n,1),(n,2)\}
\]
or its rotation is rotation-rigid, with common betas allowed.
\end{proposition}

\begin{proof}
For a zero-free five-row source the maximal cofactor is a singleton and a
three-row connected component with one negative.  Apply
\Cref{lem:singleton-rigidity}.  For larger sources its maximal cofactor
is connected and is the hook one order down.  Induction and
\Cref{lem:border} apply; the endpoint extensions remain in the same class.
Homogenization handles units.  In particular, at five rows the full chain
is only
\[
 0=y_5<y_4<y_3\le x_5<y_2\le x_4<x_3<y_1\le x_2<x_1;
\]
none of its three weak inequalities is silently made strict before
homogenization.
\end{proof}

\begin{proof}[Proof of \Cref{thm:classification}(i)--(iii)]
Since \(D\) is nonempty and connected, \(s_D\) is nonconstant and
irreducible by \cite{Barekat2009CoincidencesCorrigendum}.  An empty
diagram contributes \(1\), whereas every nonempty connected component,
including a single row or column, has a nonconstant irreducible Schur
function.  Component factorization and \(s_H=s_D\) therefore force
\(H\) to have exactly one nonempty component.  This reduction also
applies to parts (iv) and (v).

By \Cref{lem:negative-partition}, the negative sites form a partition,
rotation conjugates it, and \(\nu\le(n-1)(n-2)/2\).
At three negatives the shapes are a column, a row, or the three-cell
staircase.  For \(n\ge5\), use \Cref{thm:column,prop:three-corners};
at four rows use \Cref{prop:four-rows}.
At four negatives \(n\ge5\); up to rotation their shapes are a column,
a square, or a hook.  Apply \Cref{thm:column,thm:square,prop:hook}.
The first assertion is \Cref{thm:two-negatives}.
\end{proof}

\section{Five negatives}\label{sec:five-base}

\begin{theorem}[Five-row five-negative classification]\label{thm:five-base}
Let \(D\) be connected with five nonempty rows and exactly five full-board
negative indices.  Common betas are allowed.  Its complete connected Schur
class is its rotation orbit unless it is one of the following classes.
For integers \(1\le q\le b<g\), put
\begin{equation}\label{eq:five-families}
 \begin{gathered}
 X=(b+4g,b+3g,b+2g,b+g,b),\\
 \begin{array}{c|c|c}
 i&Y(F_i)&Y(H_i)\\\hline
 1&(3g,g+q,g,q,0)&(2g+q,2g,g,q,0)\\
 2&(3g+q,g+q,g,q,0)&(2g+q,2g,g+q,q,0).
 \end{array}
 \end{gathered}
\end{equation}
The complete class is \(\{F_i,F_i^*,H_i,H_i^*\}\).  All four members are
distinct.
\end{theorem}

\subsection{One opened cell and the sixteen-plus-eight expansion}

A five-row connected board has no negative entry on or above the first
subdiagonal.  There are six positions below that subdiagonal.  If exactly
five are negative, the remaining position must be one of
\((3,1),(4,2),(5,3)\): a nonnegative entry farther southwest would force
at least two of these six positions to be nonnegative.  Call the three
openings \(E,M,F\), respectively.  Rotation fixes \(M\) and interchanges
\(E,F\).  Equivalently, the negative row-count words are
\((0,0,0,2,3),(0,0,1,1,3),(0,0,1,2,2)\).  

In this section retain the unit variable.  For a board with \(n\) nonempty
rows write
\[
 P_D=\det(H_{X_i-Y_j})
     =H_0^n s_D(H_1/H_0,H_2/H_0,\ldots),
 \qquad H_j=0\quad(j<0),
\]
in \(\Q[H_0,H_1,\ldots]\), with all nonnegative-index variables
independent.  Equality of the original polynomials is equivalent to
equality of these homogeneous polynomials, since the row count is read
by \Cref{lem:readings}.  

Connected-skew-Schur irreducibility
\cite{Barekat2009CoincidencesCorrigendum} also holds after this
homogenization.  Indeed, the diagonal monomial is the unique
minimum-square-weight monomial and has coefficient one, so the ordinary
degree is exactly \(n\) and \(H_0\nmid P_D\).  A factorization of the
homogeneous polynomial has homogeneous factors.  After setting
\(H_0=1\), irreducibility over \(\Q\), obtained from the corrected
integral statement by Gauss's lemma, makes one factor constant.  A
homogeneous factor specializing to a nonzero constant is a scalar power
of \(H_0\).  That power is zero because \(H_0\nmid P_D\).  Thus \(P_D\)
is irreducible.  The coefficient-one minimum-square monomial fixes scalar
associates whenever connected factors are compared below.  

For a five-negative board let \(A^{(0)}\) be its matrix with the single
opened second-subdiagonal position \((k+2,k)\) set to zero, and put
\(\tau=X_{k+2}-Y_k\ge0\).  This is a change of one matrix position,
not a specialization of every occurrence of \(H_\tau\).  Multilinearity
gives
\begin{equation}\label{eq:five-open-decomposition}
 P_D=\det A^{(0)}+H_\tau C_{k+2,k}(A^{(0)}),\qquad k\in\{1,2,3\},
\end{equation}
where \(C_{ij}\) denotes the signed cofactor.  The first determinant is
\[
 \det A^{(0)}=
 \sum_{S\subseteq\{1,2,3,4\}}(-1)^{|S|}
 \prod_{i\in S}H_{z_i}
 \prod_{[a,b]\in\mathcal I(S)}H_{w_{a\cdots b}},
 \qquad
 w_{a\cdots b}=\sum_{i=a}^b r_i-\sum_{i=a}^{b-1}z_i.
\]
Here \(\mathcal I(S)\) joins adjacent rows exactly across the edges in
\(S\).  To prove the expansion, if the first row uses column \(j\),
rows \(2,\ldots,j\) must use columns \(1,\ldots,j-1\); remove that
interval and repeat.  These are the sixteen interval terms.  After the
opened row and column are deleted, the four remaining rows have nested
supports of sizes \(4,4,3,2\).  The last row has two choices, the next
has two remaining choices, and the first two rows have two matchings.  
Thus the second term of \eqref{eq:five-open-decomposition} has eight
supported permutation terms.  Interchanging the first two remaining
rows pairs their signs.  The full expansion has twelve terms of each
sign, before collection.  Repeated labels and common betas may merge or
cancel terms; no argument below treats the two layers as separately
invariant.  

\subsection{Corner readings and a prime obstruction}

After orienting an end opening as \(E\), use the normal forms
\begin{equation}\label{eq:five-normal-forms}
 M=(x,y,L,c,b;\ a,u,v,p),\qquad
 E=(x,L,y,c,b;\ u,v,a,p),\qquad L=u+v-\tau.
\end{equation}
Their complete validity conditions are
\[
 u,v>\tau\ge0,\qquad a,p\ge0,\qquad b>p,
\]
\[
 M:\quad x>a,\quad y>a+u,\quad c>v+p;
 \qquad
 E:\quad x>u,\quad y>a+v,\quad c>a+p.
\]
The three second-subdiagonal indices are
\(z_i+z_{i+1}-r_{i+1}\); these inequalities make exactly the indicated
one nonnegative.  Strict ordering then forces every more southwestern
entry to be negative.  

Put \(q=b-p\), \(w=c+q\), and define
\[
 \begin{aligned}
 B&=H_uH_v-H_\tau H_L,&
 F_s(z)&=H_sH_z-H_\tau H_{z+s-\tau},\\
 T(z)&=H_bH_z-H_pH_{z+q},&
 G_s(z)&=H_sT(z)-H_\tau T(z+s-\tau),\\
 Q_s(z)&=H_xF_s(z)-H_aF_s(z+x-a).
 \end{aligned}
\]
The symbol \(B\) in these subsections denotes the displayed binomial,
not a beta list.  It is primitive and linear in its unique largest
variable \(H_L\), so it is irreducible and has no coordinate-variable
factor.  The same argument applies to \(F_s(z)\) for \(s,z>\tau\)
and to \(T(z)\) for \(z>p\).  On an \(M\)-board, \(G_u(c),Q_v(y)\)
are the connected three-row determinants
\((u,c,b;\tau,p)\), \((x,y,v;a,\tau)\).  Indeed,
\(c>v+p>\tau+p\) and \(y>a+u>a+\tau\), and the remaining
endpoint inequalities are inherited from that board.  They are therefore
irreducible after homogenization.  These cubic identifications are not
asserted for arbitrary \(E\)-parameters.  Every candidate is required to
satisfy its own pattern inequalities.  The single opening supplies the
prime factor \(B\); closing it gives the product \(H_uH_v\), so this
factorization argument does not extend to the six-negative mask.  

Let \(m=X_1\), and put \(\delta=X_1-X_2\), equal to \(x-a\) for
\(M\) and \(x-u\) for \(E\).  Direct complementary expansion gives
\begin{equation}\label{eq:five-marker}
 \begin{gathered}
 C_{15}=H_aH_pB,\qquad C_{14}=-H_aH_bB,\\
 C_{25}^{M}=-H_xH_pB,\qquad
 C_{25}^{E}=-H_aH_pF_v(x),\\
 N=m-\min(q,\delta),\qquad
 [H_N](P_D-H_mC_{15})=
 \begin{cases}
 C_{14},&q<\delta,\\
 C_{25},&\delta<q,\\
 C_{14}+C_{25},&q=\delta.
 \end{cases}
 \end{gathered}
\end{equation}
The largest generator occurs only at \((1,5)\).  Its nonzero cofactor
reads \(m,\tau,L,\{u,v\},\{a,p\}\).  Only \((1,4),(2,5)\) can
carry the next index \(N\), and these positions cannot occur together:
row five would have no supported column left.  Their coefficients in
\eqref{eq:five-marker} are nonzero, including at ties.  

For \(M\), divide the second coefficient algebraically by \(B\),
then set all remaining variables to one.  The values are \(-1\) in
a strict regime and \(-2\) at a tie.  Rotation interchanges
\(q,\delta\).  For \(E\), writing \(C=C_{15}\) and \(P\) for the
second coefficient, the normalized greatest common divisors are
\[
 \gcd(C,P)\sim
 \begin{cases}
 H_aB,&q<\delta,\\
 H_aH_p,&\delta<q,\\
 H_a,&q=\delta.
 \end{cases}
\]
Their ordinary degrees are \(3,2,1\).  At a tie the parenthesis in
\(-H_a(H_bB+H_pF_v(x))\) is divisible by neither \(H_p\) nor \(B\):
\(b>p\), and \(F_v(x)\) is an irreducible quadratic distinct from
\(B\), since \(x>u\).  

\begin{lemma}[Prime obstruction]\label{lem:five-prime}
Suppose \(u>v>\tau\ge0\), \(b>p\ge0\), and \(c>v+p\).  Put
\(d=c+u-v\).  Then
\[
 B\nmid H_uT(c)-H_vT(d).
\]
\end{lemma}

\begin{proof}
Write \(P=H_uT(c)-H_vT(d)\), and suppose \(B\mid P\).  We have
\(d>u+p\).  Set \(H_u=H_\tau=0\).  The divisor becomes zero, while
\[
 P\longmapsto
 H_v\bigl(H_pH_{d+q}-H_bH_d\bigr),
\]
with \(H_p,H_b\) also specialized if their indices are \(u\) or
\(\tau\).  The factors \(H_v,H_d,H_{d+q}\) survive.  If both
monomials in parentheses survive, they are distinct: \(p<b\) and
\(p<d\), so \(\{p,d+q\}\ne\{b,d\}\).  Exactly one surviving
monomial is also impossible.  Thus both must vanish, forcing
\(p=\tau,b=u\).  Now set \(H_v=H_\tau=0\) instead.  Again \(B\)
becomes zero, but \(P\) becomes \(H_u^2H_c\ne0\), since
\(c>v+\tau\).  This contradiction proves the assertion.  

In particular, no inequality \(d>b\) is assumed.  If \(d\le b\),
then \(b\ge d>u>\tau\), so the first specialization already leaves
the nonzero term \(-H_vH_bH_d\), or \(-H_vH_b^2\) when \(d=b\).  
Its only possible companion cannot cancel it by the same multiset
comparison.  This includes \(\tau=0\), where \(H_\tau\) is the
independent variable \(H_0\).  
\end{proof}

\subsection{Same-opening rigidity}

\begin{proposition}\label{prop:five-same-opening}
Equal five-negative \(M\)-boards are identical or rotations.  Equal
five-negative \(E\)-boards, both oriented with opening \((3,1)\),
are identical after translation.  
\end{proposition}

\begin{proof}
All comparisons are between collected homogeneous polynomials.  A
supported complementary minor is nonzero: its order-preserving matching
is supported and uniquely minimizes square weight.  This justifies
using a unique largest remaining matrix index throughout the proof.  

\emph{Middle opening with both endpoints fixed.}
Suppose first that \(x,a,b,p\) agree.  The row multiset and the first
cofactor leave only
\begin{equation}\label{eq:within-candidates}
 \{y,c\}=\{\eta,d\},\qquad \{u,v\}=\{U,V\}.
\end{equation}
Put \(A(z)=H_xH_z-H_aH_{z+x-a}\).  Subtract the common contribution
\[
 B\bigl(H_aH_pH_m-H_aH_bH_{m-q}
       -H_xH_pH_{m-\delta}+H_xH_bH_{m-\delta-q}\bigr).
\]
Its four permutation sets are disjoint, since row five already uses
one of columns four and five.  The remainder is
\begin{equation}\label{eq:five-middle-remainder}
 \begin{aligned}
 \mathcal R={}&H_LA(y)T(c)-H_vA(y)T(c+u-\tau)\\
 &-H_uA(y+v-\tau)T(c)
   +H_\tau A(y+v-\tau)T(c+u-\tau).
 \end{aligned}
\end{equation}
Its maximal candidate positions are \((1,3),(3,5)\), with indices
\(y+v-\tau+\delta,c+u-\tau+q\) and cofactors
\(H_aG_u(c),H_pQ_v(y)\).  

If only the core is interchanged, orient \(u>v\).  A unique maximum
cannot vanish, so equality forces \(y+\delta=c+q\).  The coefficient
at the resulting contact is
\[
 H_pH_xF_v(y)-H_aH_bF_v(c).
\]
Unique factorization gives \(p=a,x=b,y=c\), exactly rotation.  

If only \(y,c\) are interchanged, orient \(y>c\).  The maximum
forces \(v+\delta=u+q\), with coefficient
\(H_aG_u(c)-H_pQ_v(c)\).  Irreducibility gives \(a=p\) and equality
of the cubics.  Their row multisets give \(\{u,b\}=\{x,v\}\).  
The alternative \(u=v,b=x\) is rotation.  Otherwise \(x=u,b=v\).  
The remaining cubic comparison is between
\((u,c,v;\tau,p)\) and \((u,c,v;p,\tau)\).  If \(u\ne v\) and
\(p\ne\tau\), the four noncommon high indices
\[
 c+u-p,\quad c+u-\tau,\quad c+v-p,\quad c+v-\tau
\]
have a unique maximum with nonzero coefficient.  Hence \(p=\tau\).  
For the resulting candidate pair
\[
 (u,y,L,c,v;\tau,u,v,\tau),\qquad
 (u,c,L,y,v;\tau,u,v,\tau),
\]
the full difference is \(-B\mathcal S\), where
\begin{equation}\label{eq:five-survivor}
 \begin{aligned}
 \mathcal S={}&H_u(H_cH_{y+v-\tau}-H_yH_{c+v-\tau})\\
 &-H_v(H_cH_{y+u-\tau}-H_yH_{c+u-\tau})\\
 &+H_\tau(H_{c+v-\tau}H_{y+u-\tau}
             -H_{c+u-\tau}H_{y+v-\tau}).
 \end{aligned}
\end{equation}
For \(u>v\), its largest-generator coefficient is \(-F_v(c)\) at
\(H_{y+u-\tau}\); for \(v>u\), it is \(F_u(c)\) at
\(H_{y+v-\tau}\).  Neither vanishes.  Thus this distinct candidate is
not an equality.  

If both pairs are interchanged, put \(h=y+v\), \(k=c+u\).  The four
candidate indices are
\[
 h+\delta-\tau,\quad k+q-\tau,\quad
 k+\delta-\tau,\quad h+q-\tau.
\]
For \(h\ne k\) and \(\delta\ne q\) the maximum is unique.  If
\(h\ne k\), equality therefore forces \(\delta=q\); the contact
coefficient is, up to sign,
\((H_aH_b-H_pH_x)F_u(c)\), or its counterpart with \(F_v(y)\).  
It forces \(a=p,b=x\), giving rotation.  If \(h=k\) but
\(\delta\ne q\), equality of the active cubic cofactors reads
\(\{u,c\}=\{v,y\}\).  The crossed assignment requires \(c=v,y=u\),
contrary to validity; the other assignment is the unchanged profile.  
At \(h=k,\delta=q\), the two maximal positions can occur together.  
Both determinants have the same quadratic term
\(H_aH_pH_\tau H_\nu^2\), where \(\nu=h+\delta-\tau\).  After
canceling that term, the linear coefficient of the difference is
\[
 (H_aH_b-H_pH_x)(H_uH_c-H_vH_y).
\]
The second factor is nonzero for a distinct candidate, again by
\(c>v,y>u\).  The first forces rotation.  This proves the assertion
with both endpoints fixed.  

\emph{Middle opening with only the bottom endpoint fixed.}
Now assume only \(a,b,p\) agree, and write the mate as
\[
 M'=(\xi,\eta,L,d,b;\ a,U,V,p),\qquad
 \{\xi,\eta,d\}=\{x,y,c\},\quad \{U,V\}=\{u,v\}.
\]
After subtracting the two common top-right terms, the maximal candidates
are
\[
 R=x+y+v-a-\tau,\quad R'=\xi+\eta+V-a-\tau,\quad
 Z=c+q+y-\tau,\quad Z'=d+q+\eta-\tau.
\]
Let \(\epsilon,\epsilon',\zeta,\zeta'\) be the indicators of their
attaining the maximum.  Its collected coefficient is
\[
 H_a\bigl(\epsilon G_u(c)-\epsilon'G_U(d)\bigr)
 -H_pB\bigl(\zeta H_x-\zeta'H_\xi\bigr).
\]
There is no quadratic term: \((1,3),(2,5)\) cannot occur together.  
If a top candidate is active, equality forces
\(B\mid\epsilon G_u(c)-\epsilon'G_U(d)\).  A single active top
candidate is impossible by irreducibility of the cubic.  Thus both
are active, and the row sum gives \(d=c+V-v\).  A genuine core
reversal, oriented with \(u>v\), would imply
\[
 B\mid G_u(c)-G_v(c+u-v)
       =H_uT(c)-H_vT(c+u-v),
\]
contrary to \Cref{lem:five-prime}.  Hence \(d=c,U=u,V=v\).  
Only \(x,y\) can now interchange.  First-column expansion reduces
that equality to \(H_x\mathcal A_y=H_y\mathcal A_x\), where
\(\mathcal A_z\) is the connected four-row determinant
\((z,L,c,b;u,v,p)\).  Coordinate-factor-freeness excludes \(x\ne y\).  
If no top candidate is active, both right candidates are active and
force \(x=\xi\).  Both endpoints are fixed, so the preceding argument
applies.  

For unrestricted \(M/M\), rotate both boards so \(q\le\delta\).  
The second-corner evaluation separates strict and tied regimes.  In
the strict regime, \(C_{15},C_{14}\) give the products
\(H_aH_p,H_aH_b\).  Their reduced Laurent quotient \(H_b/H_p\)
reads \(b,p\), since \(b>p\), and then \(a\).  In the tied regime,
\(m-N\) reads the common value \(q=\delta=d_0\).  Rotate one board
if necessary to align \(\{a,p\}\).  Then
\(x=a+d_0,b=p+d_0\), so both endpoints agree.  This proves the
middle-opening assertion.  

\emph{End opening: a joint-anchor identity.}
Put
\[
 A_0=H_xH_L-H_uH_{x+v-\tau},\qquad
 J_z=A_0H_z-F_v(x)H_{z+u-\tau}+BH_{x+z-\tau}.
\]
For \(z>v\), \(J_z\) is the irreducible connected three-row
determinant \((x,L,z;u,v)\).  The first three row betas are
\(X_1=m,X_2=m-x+u,X_3=m-x+\tau\).  Direct expansion gives
\begin{equation}\label{eq:five-end-block}
 P_E=J_yT(c)+H_a\sum_{i=1}^3\lambda_i
       (H_pH_{X_i}-H_bH_{X_i-q}),\qquad
 (\lambda_1,\lambda_2,\lambda_3)=(B,-F_v(x),A_0).
\end{equation}
The six indicated permutation sets are disjoint.  If
\(x,u,v,\tau,m\) and \(a,b,p\) agree, subtracting them gives
\(J_yT(c)=J_{y'}T(c')\).  The irreducible factors have degrees three
and two, so factorization reads \(c\), then \(y\).  Thus joint
agreement of these anchors forces identity.  

\emph{End opening with the upper prefix fixed.}
Assume \(x,u,v,\tau,m\) agree.  The remaining invariant data are
\(\{a,p\}\) and \(\{y,c,b\}\).  Subtract the common contribution
of the first three entries in column five.  The remainder is
\[
 \mathcal R_E=J_yT(c)-H_aH_bJ_{y+c-a}.
\]
Its largest index is \(N_1=\max(m-q,w)\), with coefficient
\[
 \begin{cases}
 -H_aH_bB,&m-q>w,\\
 -H_pJ_y,&w>m-q,\\
 -H_aH_bB-H_pJ_y,&w=m-q.
 \end{cases}
\]
These are nonzero, and divisibility by \(B\) holds exactly in the
first regime; otherwise the irreducible quadratic would divide the
irreducible cubic \(J_y\).  If both boards are top-only, the products
\(H_aH_b,H_aH_p\) recover the bottom endpoint and the joint-anchor
identity applies.  Otherwise both have \(w=w'=N_1\).  

If \(a'=a,p'=p\), common \(X_3,w\) gives \(y'=y\).  Interchanging
\(c,b\) leaves \(T(c)\) unchanged and would require
\[
 H_bJ_{y+c-a}=H_cJ_{y+b-a}.
\]
The cubic factors have no coordinate divisor, so \(b=c\), and the
board is unchanged.  If instead \(a'=p,p'=a\), with \(a\ne p\),
then \(y'=y+p-a\).  Let \(e,e'\) indicate whether the top candidate
also attains \(N_1\).  Equality gives
\[
 H_pJ_y+eH_aH_bB=H_aJ_{y'}+e'H_pH_{b'}B.
\]
Setting \(H_a=0\) forces \(e'=1\), and setting \(H_p=0\) forces
\(e=1\), by coordinate-factor-freeness.  Both boards are therefore
tied, so \(q=q'\), \(c'=c\), \(b'=b+a-p=x+y-\tau\).  
Canceling common rows gives \(\{y,b\}=\{y',b'\}\).  Since both
entries changed, \(y=b'\), forcing \(x=\tau\), contrary to \(x>u\).  
Thus the upper-prefix route gives identity.  

\emph{End opening with the bottom edge fixed.}
Assume \(a,b,p,c\) agree.  Only \(x,y\) and \(u,v\) can be
interchanged.  The top-row contributions in columns three, four, and
five agree and may be subtracted.  

For a core interchange, orient \(u>v\) and put \(z=w+y-a\).  
The maximum is \(\max(x,z)+u-\tau\), so equality forces \(x=z\).  
The contact coefficient is
\[
 K=F_v(y)T(c)-H_aH_bF_v(c+y-a).
\]
With \(J_2=c+y+v-a-\tau\), its largest index is
\(\max(w,J_2)\).  The coefficients for \(J_2>w,w>J_2,w=J_2\)
are respectively
\[
 H_aH_bH_\tau,\qquad -H_pF_v(y),\qquad
 -H_pF_v(y)+H_aH_bH_\tau.
\]
The last evaluates to one at all remaining variables equal to one;
all three are nonzero.  

For a free-row interchange, orient \(y>x\).  The maximum compares
\(y+v-\tau\) and \(w+y-a+u-\tau\).  Equality requires
\(w=a+v-u\), hence \(v>u+b\).  Its coefficient is
\[
 \begin{aligned}
 K={}&F_u(x)T(c)-H_aH_bF_u(c+x-a)\\
 &+H_aH_pH_uH_{w+x-a}-H_aH_pH_xH_v.
 \end{aligned}
\]
Its largest remaining index is the maximum of
\(x+v-u\) and \(x+v-\tau-q\), with coefficients
\(H_aH_pH_u\), \(H_aH_bH_\tau\), or their sum at a tie.  None
vanishes.  For both interchanges, the maximum compares
\(w+y-a+u-\tau\) and \(w+x-a+v-\tau\).  Each exceeds the opposite
top-column-two index by \(w-a>0\).  Equality forces
\(y+u=x+v\), with coefficient
\[
 -H_aH_p(H_xH_v-H_yH_u).
\]
Its vanishing gives only the identical candidate: the crossed assignment
is excluded by \(x>u,y>v\).  This proves bottom-edge rigidity.  

\emph{End opening with only the bottom endpoint fixed.}
Assume \(a,b,p\) agree.  Write
\[
 E'=(\xi,L,\eta,d,b;\ U,V,a,p),\qquad
 \{\xi,\eta,d\}=\{x,y,c\},\quad \{U,V\}=\{u,v\}.
\]
After the two common top-right terms are subtracted, the maximal
candidates are
\[
 R=x+y-\tau,\quad R'=\xi+\eta-\tau,\quad
 Z=c+q+y+u-a-\tau,\quad Z'=d+q+\eta+U-a-\tau.
\]
With indicators for the maximum, its coefficient is
\[
 B\bigl(\epsilon T(c)-\epsilon'T(d)\bigr)
 -H_aH_p\bigl(\zeta F_v(x)-\zeta'F_V(\xi)\bigr).
\]
If a top candidate is active, setting \(H_a=0\) excludes a single
active top term.  Both top indices agree, so the row sum gives \(c=d\)
and bottom-edge rigidity applies.  Otherwise both right candidates are
active and \(F_v(x)=F_V(\xi)\).  The multiset intersection of
\(\{u,v\}\) and \(\{x,v\}\) is exactly \(\{v\}\), since \(x>u\).  
It reads \(x=\xi,u=U,v=V\), and the joint-anchor identity applies.  

Finally, the three gcd degrees separate unrestricted \(E/E\) equality
into equal corner regimes.  For \(q<\delta\), the first two
coefficients read \(a,b,p\), so the bottom-endpoint route applies.  
For \(\delta<q\), division by the known \(-H_aH_p\) reads
\(F_v(x)\), which together with \(B\) reads \(x,u,v,\tau\) in
order.  The common \(m\) supplies the upper prefix.  At \(q=\delta\),
the degree-one gcd reads \(a\); the first coefficient then reads
\(p,B\).  Reducing the second-coefficient equality modulo \(H_p\)
reads \(b\), and subtraction reads \(F_v(x)\).  The joint-anchor
identity finishes.  We have proved
\begin{equation}\label{eq:within-inventory}
 P_M=P_{M'}\ \Longrightarrow\ M'\in\{M,M^*\},\qquad
 P_E=P_{E'}\ \Longrightarrow\ E'=E.
\end{equation}
\end{proof}

\subsection{Mixed openings and arithmetic progression}

\begin{proposition}\label{prop:five-mixed-opening}
If a five-negative \(M\)-board and a five-negative end-opening board
are Schur-equal, translations and rotations put them into exactly one
of the two pairs in \eqref{eq:five-families}.  
\end{proposition}

\begin{proof}
Orient the end opening as \(E\), and use the notation of
\eqref{eq:five-normal-forms}.  For \(M\), the second coefficient in
\eqref{eq:five-marker} is always divisible by \(B\).  For \(E\),
it is divisible by \(B\) exactly when \(q<\delta\), since
\(F_v(x)\) is an irreducible quadratic distinct from \(B\).  
Thus the \(E\)-corner is strict with \(q<\delta\).  The \(M\)-corner
cannot tie: after division by \(B\), evaluation at all variables equal
to one gives \(-2\), whereas the \(E\)-coefficient gives \(-1\).  
Rotate \(M\) if necessary so that it too has \(q<\delta\).  
The first two coefficients now give common \(a,b,p\) through the
reduced quotient \(H_b/H_p\).  Consequently the general mixed pair is
\begin{equation}\label{eq:cross-candidates}
 \begin{gathered}
 M=(x,y,L,c,b;\ a,u,v,p),\qquad
 E'=(\xi,L,\eta,d,b;\ U,V,a,p),\\
 \{\xi,\eta,d\}=\{x,y,c\},\qquad \{U,V\}=\{u,v\}.
 \end{gathered}
\end{equation}
No equality \(d=c\) is assumed.  

Subtract the common two-corner polynomial
\(H_aB(H_pH_m-H_bH_{m-q})\).  The maximal remaining candidates are
\[
 R=x+y+v-a-\tau,\quad R'=\xi+\eta-\tau,\quad
 Z=w+y-\tau,\quad Z'=d+q+\eta+U-a-\tau.
\]
Let \(J\) be their maximum, and write
\(\epsilon,\epsilon',\zeta,\zeta'\) for the corresponding indicators.  
The coefficient in the difference is
\begin{equation}\label{eq:five-mixed-coefficient}
 K=\epsilon H_aG_u(c)-\epsilon'BT(d)
       -\zeta H_xH_pB+\zeta'H_aH_pF_V(\xi).
\end{equation}
The positions \((1,3),(2,5)\) cannot occur together, so this is the
whole coefficient, not a derivative containing a further \(H_J\).  
At least one top and one right candidate must attain \(J\).  A
right-only tie would require \(H_xB=H_aF_V(\xi)\), impossible by
coordinate-factor-freeness and \(x>a\).  A top-only tie would equate
\(H_aG_u(c)\) with the coordinate-factor-free product \(BT(d)\).  
Single active terms have nonzero cofactors.  

Set \(H_p=0\).  If only \(R'\) is the active top term, the
specialization is \(-H_bH_dB\ne0\).  If only \(R\) is active, it
is \(H_aH_bF_u(c)\), which vanishes only for \(a=p\).  If both top
terms are active, their indices and the row sum give
\(d=c+a-v\).  The specialization condition is
\[
 H_p\mid
 H_aH_uH_c-H_aH_\tau H_{c+u-\tau}
 -H_dH_uH_v+H_dH_\tau H_L.
\]
Here \(c,d,c+u-\tau>p\).  Consider the possible identifications in
this order.  If \(p=a\), the residual is \(-H_dB\ne0\).  If
\(p=\tau\), equality requires \(H_aH_c=H_dH_v\), hence
\(a=v,d=c\).  If \(p=u\), it requires
\(H_aH_{c+u-\tau}=H_dH_L\), hence
\(a=L,d=c+u-\tau\).  If \(p=v\) but none of the preceding cases
holds, the three surviving terms evaluate to one at all remaining
variables equal to one.  If \(p=L\), they evaluate to minus one.  
If none of these identifications holds, the original residual would
have to vanish.  But \(H_aF_u(c)=H_dB\) is impossible by
coordinate-factor-freeness and \(d>a\).  
Thus exactly the following possibilities remain:
\begin{equation}\label{eq:five-mixed-regimes}
 \begin{array}{c|l}
 \mathrm{I}&R>R',\quad a=p,\\
 \mathrm{II}&R=R',\quad a=v,\quad \tau=p,\quad d=c,\\
 \mathrm{III}&R=R',\quad p=u,\quad a=L,\quad d=c+u-\tau.
 \end{array}
\end{equation}
The order of the specialization cases includes all coincident-index
possibilities, including \(p=0\).  

\emph{Regime I.}
Divide \eqref{eq:five-mixed-coefficient} by \(H_p\).  If \(Z\) is
absent, the result is \(G_u(c)+H_pF_V(\xi)\), nonzero modulo
\(H_p\).  Therefore \(Z\) is active, and its contact gives
\[
 w=x+v-p,\qquad x=c+b-v>b.
\]
If \(Z'\) is absent, the coefficient is \(G_u(c)-H_xB\).  
Its largest possible index is \(\max(x,x+L-p)\).  The coefficients
for \(L>p,L<p,L=p\) are respectively
\[
 H_pH_\tau,\qquad -B,\qquad 2H_pH_\tau-H_uH_v,
\]
all nonzero.  Hence \(Z'\) is active too.  The right tie and row sum
give \(\xi=x+U-p\).  The highest shifted terms cancel, leaving
\[
 K/H_p=H_bF_u(c)-H_xB+H_p(H_VH_\xi-H_uH_w).
\]
If \(U=v\), then \(\xi=w\) and the last term vanishes; unique
factorization would force \(b=x\), impossible.  Thus \(U=u,V=v\)
and \(u\ne v\).  The remaining equation is
\[
 \begin{aligned}
 0={}&H_bH_uH_c-H_bH_\tau H_{c+u-\tau}
       -H_xH_uH_v+H_xH_\tau H_L\\
    &+H_pH_vH_{x+u-p}-H_pH_uH_{x+v-p}.
 \end{aligned}
\]
Put \(A=c+u-\tau=x+L-b\), \(C=x+u-p\), \(D=x+v-p\).  
Every largest-index candidate is among \(A,C,D,x\), with uncollected
coefficients \(-H_bH_\tau,H_pH_v,-H_pH_u,-B\).  

First suppose \(\max(u,v)>p\).  If \(u<v\), the largest candidates
\(A,D\) have negative nonzero coefficients, so the equation is impossible.  
Otherwise \(u>v\) and \(u>p\), since \(u\ne v\) was already established.  
Thus \(C>D,x\); equality forces
\(A=C\) and \(H_bH_\tau=H_pH_v\).  Hence \(b=v,\tau=p\), and
the equation reduces to
\[
 H_p(H_cH_L-H_uH_{c+v-p})=0,
\]
which forces \(c=u\), since \(v>p\).  

For completeness suppose \(\max(u,v)\le p\).  If \(L>b\), then
\(A\) is a unique maximum.  If \(L<b\), the coefficient at \(x\)
is \(-B\), \(H_\tau H_L\), or \(H_\tau H_L-2H_uH_v\), according
as neither core index equals \(p\), \(u=p\), or \(v=p\).  None
vanishes.  If \(L=b\), that coefficient is nonzero except possibly
when \(u=p>v\).  In that last case the whole polynomial becomes
\[
 H_p(H_LH_c-H_pH_{c+v-\tau}),
\]
which is nonzero because \(c>p+v>p\) and \(v>\tau\).  
Therefore Regime I forces
\[
 b=v,\quad \tau=p,\quad c=u,\quad x=u.
\]
Since \(\xi=2u-p>u\), the row multiset then gives
\(y=\xi=2u-p\), \(\eta=d=u\).  

\emph{Regime II.}
Here \(a=v,\tau=p,d=c\).  Dividing the contact coefficient by
\(H_p\) gives
\[
 H_LT(c)-H_vT(c+u-p)-\zeta H_xB+\zeta'H_vF_V(\xi).
\]
Without \(Z\), reduction modulo \(H_v\) leaves \(H_LT(c)\ne0\).  
Thus \(Z\) is active and \(w=x\).  Without \(Z'\), the unique
largest variable \(H_{w+u-p}\) has coefficient \(H_pH_v\ne0\), so
\(Z'\) is active and \(\xi=w+U-v\).  The choice \(U=v\) leaves
\[
 H_b(H_LH_c-H_vH_{c+u-p})\ne0,
\]
since \(c>v\).  Thus \(U=u,V=v\), and the residual is
\[
 H_b(H_LH_c-H_vH_{c+u-p})
 +H_v(H_vH_\xi-H_uH_w).
\]
Modulo \(H_v\), this is \(H_bH_LH_c\), with \(c,L>v\), so
\(b=v\).  The residual then becomes
\(H_v(H_LH_c-H_uH_{c+v-p})\), forcing \(c=u\).  Hence
\(x=u+v-p\), \(\xi=y=2u-p\), and \(\eta=x\).  

\emph{Regime III.}
Here \(p=u,a=L,d=c+u-\tau\), and the contact coefficient divided by
\(H_u\) is
\[
 H_LT(c)-H_vT(d)-\zeta H_xB+\zeta'H_LF_V(\xi).
\]
Reduction modulo \(H_L\) forces \(Z\) to be active.  Its contact
gives \(x=c+b-\tau\).  Since \(b>u\), we have \(x>d>c\).  
As \(d\in\{x,y,c\}\), necessarily \(d=y\) and
\(\{\xi,\eta\}=\{x,c\}\).  Without \(Z'\), direct cancellation
leaves the unique largest variable \(H_x\) with coefficient
\(H_LH_\tau\ne0\).  Therefore \(Z'\) is active, giving
\(\xi=w+U-v\), where \(w=c+b-u\).  
If \(U=v\), then \(c<\xi=w<x\), impossible.  If \(U=u\), then
\(\xi=c+b-v<x\), so \(\xi=c,b=v>u,\eta=x\).  The residual is
\[
 2H_LH_vH_c-H_LH_uH_{c+v-u}-H_v^2H_{c+u-\tau}.
\]
Since \(c>u+v\) and \(v>u>\tau\), its largest index is one or both
of \(c+v-u,c+u-\tau\), with negative nonzero coefficients.  Regime
III is impossible.  

The necessary mixed parameters are therefore exactly
\begin{equation}\label{eq:cross-output}
 \begin{gathered}
 \tau=p,\quad b=v,\quad c=u,\quad y=2u-p,\quad u>v+p,\\
 (\xi,\eta,d;U,V)=(y,x,c;u,v),\\
 (a,x)=(v,u+v-p)\quad\text{or}\quad(a,x)=(p,u).
 \end{gathered}
\end{equation}
Set \(g=u-p\) and replace the family parameters \(b,q\) by
\(v,v-p\).  The domain becomes \(1\le q\le b<g\), and the profile
recurrence gives exactly \eqref{eq:five-families}.  

\end{proof}

\begin{proof}[Proof of the identities in \eqref{eq:five-families}]
Reverse the five rows, which has positive sign,
and index them by \(k=0,\ldots,4\).  For integers \(j\ge0\), put
\[
 U_j=(H_{b+(k-j)g})_{k=0}^4,\qquad
 V_j=(H_{b+(k-j)g-q})_{k=0}^4.
\]
Let \(S_0\) be the down-shift on five-vectors.  Since \(b<g\) and
\(q\le b\), the negative-index convention gives
\(S_0U_j=U_{j+1}\) and \(S_0V_j=V_{j+1}\).  For arbitrary columns,
\[
 \sum_{i=1}^5
 \det(z_1,\ldots,S_0z_i,\ldots,z_5)=0,
\]
by taking the linear coefficient in
\(\det((I+tS_0)Z)=\det Z\).  Apply this identity to the two seeds
\[
 (U_2,V_1,U_1,V_0,U_0),\qquad
 (V_2,V_1,U_1,V_0,U_0).
\]
In each case three summands have a repeated column.  The two remaining summands give
\[
 \begin{aligned}
 \det(U_3,V_1,U_1,V_0,U_0)
   &=\det(V_2,U_2,U_1,V_0,U_0),\\
 \det(V_3,V_1,U_1,V_0,U_0)
   &=\det(V_2,U_2,V_1,V_0,U_0).
 \end{aligned}
\]
These are precisely the two identities in \eqref{eq:five-families}.  
The proof retains \(H_0\), so it includes \(q=b\).  
\end{proof}

\begin{proof}[Completion of the proof of \Cref{thm:five-base}]
By \Cref{lem:five-mate-closure}, every connected mate has five rows and
five negatives.  \Cref{prop:five-same-opening} excludes every
nonrotation same-opening equality, including two end openings after
orientation.  \Cref{prop:five-mixed-opening} gives all remaining pairs
and proves their identities.  

Their parameters are unique.  From the primitive quadratic factor in
the first cofactor, recover
\[
 b=\min(u,v),\qquad
 g=\max(u,v)-\tau,\qquad q=b-\tau.
\]
The singleton pair is \(\{b,b-q\}\) in family 1 and
\(\{b-q,b-q\}\) in family 2; these are distinct since \(q>0\).  
Within each family \(F_i\) has middle opening, and \(H_i,H_i^*\)
have the two different end openings.  Also \(F_i\ne F_i^*\): its row
betas form an arithmetic progression, whereas its column betas do not,
as their first gap is larger than \(g\) and their last gap is \(q<g\).  
Thus the four members are distinct.  Same-opening rigidity and the
unique mixed parameters exclude additional members.  All arguments
were in the faithful homogeneous polynomial ring, including the
common-beta boundary.  
\end{proof}

\subsection{Endpoint and higher-row exclusion}\label{sec:five-borders}

For reference, zero-free five-negative boards have the following border
coordinates, up to rotation.  The notation in this subsection is local.  
Put \(Q=\langle A,B,d\rangle\), \(0<d<\min(A,B)\),
\(P=A+B-d\), and \(f_Q=h_Ah_B-h_Ph_d\).  Then
\begin{equation}\label{eq:five-profiles}
 \begin{aligned}
 T_1&=(s+u,p+\ell,A+a,P,B+v;s,p,A,B),\\
 &s,p,u,v>0,\quad \ell>s,\quad a>p;\\
 T_2&=(s+u,A+L,P,B+R,r+v;s,A,B,r),\\
 &s,r,u,v>0,\quad L>s,\quad R>r.
 \end{aligned}
\end{equation}
Here \(T_1^*\) has opening \(E\) and \(T_2\) has opening \(M\).  
The local \(u,v\) are exterior margins, equal to \((q,\delta)\)
on \(T_1^*\) and \((\delta,q)\) on \(T_2\).  Expansion of
\(\Phi_Q(a,v)=s_{\B_{a,v}Q}\) gives
\begin{equation}\label{eq:phi-five}
 \begin{split}
 \Phi_Q(a,v)={}&h_{P+a+v}(h_Ah_B-h_Ph_d)\\
 &-h_{P+v}(h_{A+a}h_B-h_{P+a}h_d)
   +h_{B+v}(h_{A+a}h_P-h_{P+a}h_A).
 \end{split}
\end{equation}
With \(C_Q(p;\ell,a,v)=h_{p+\ell}\Phi_Q(a,v)-h_p\Phi_Q(a+\ell,v)\),
first-column and last-row expansion yield
\begin{equation}\label{eq:t-polynomials}
 \begin{aligned}
 s_{T_1}={}&h_{s+u}C_Q(p;\ell,a,v)-h_sC_Q(p;\ell+u,a,v),\\
 s_{T_2}={}&h_{s+u}h_{r+v}\Phi_Q(L,R)-h_sh_{r+v}\Phi_Q(L+u,R)\\
 &-h_{s+u}h_r\Phi_Q(L,R+v)+h_sh_r\Phi_Q(L+u,R+v).
 \end{aligned}
\end{equation}

\begin{lemma}\label{lem:five-endpoint}
For a connected five-row diagram \(C\) with exactly five full-board
negatives, common betas allowed, and \(t>0\),
\(s_{C^{\uparrow t}}=s_{C^{\downarrow t}}\) if and only if \(C=C^*\).
\end{lemma}

\begin{proof}
Both endpoint extensions are connected and have the same negative
positions as \(C\): only the first matrix row or the last matrix
column changes, and those indices are positive before extension.  
They therefore have the same opening.  By
\Cref{prop:five-same-opening}, equality makes them identical or
rotations.  They are not identical for \(t>0\), since their ordered
first row lengths differ.  Hence
\[
 C^{\uparrow t}=(C^{\downarrow t})^*=(C^*)^{\uparrow t}.
\]
Injectivity of the first-row extension gives \(C=C^*\).  Conversely,
that symmetry and rotation invariance give equality of the extensions.  
The same-opening proposition already includes common betas.  
\end{proof}

\begin{lemma}\label{lem:five-exceptional-borders}
No positive border over either rotation orbit of \(F_i\) in
\eqref{eq:five-families} is Schur-equal to a positive border over the
\(H_i\) orbit, for \(i=1,2\) and \(1\le q\le b<g\).
\end{lemma}

\begin{proof}
Every positive endpoint extension of an exceptional five-row member is
rotation-rigid by \Cref{thm:five-base}.  To check it leaves the exceptional
list, note that one beta list of every exceptional member is arithmetic.
Its displayed column list is not arithmetic, and its first three gaps
are not all equal.  A top extension makes the row gaps
\((g+t,g,g,g)\), so neither list remains arithmetic.
After a bottom extension is normalized, its row list stays arithmetic,
but a displayed \(F\)-type member would require \(y_3=g\), whereas the
extension has \(y_3=g+t\); an \(H\)-type would require \(y_2=2g\),
whereas the extension has \(y_2=2g+t\).
The negative patterns distinguish these two types.

No extension of \(F_i\) can be an extension of \(H_i\), or its rotation:
for equal directions literal equality can be undone; for different
directions one row list is arithmetic and the other is not.
Under rotation the bottom row-beta gap of an \(H_i\) extension is its
original top column-beta gap \(q<g\), whereas an \(F_i\) extension has
bottom row-beta gap \(g\).
Thus \Cref{lem:border} excludes all unequal-margin cross-orbit borders.

At tied margins \((t,t)\), canceling common parent rows leaves respectively
\[
 \{g+t,3g-q\}=\{2g-q+t,2g\},\qquad
 \{g-q+t,3g-q\}=\{2g-q+t,2g-q\}.
\]
Each forces \(t=g\).  In the unnormalized row order \(b+kg\),
\(k=5,4,\ldots,0\), put \(a_k=h_{b+kg}\),
\(c_k=h_{b+kg-q}\); negative \(k\) give zero.
The four row vectors of the two pairs are
\[
\begin{array}{c|c}
\B_{g,g}F_1&(a_{k-3},c_{k-1},a_{k-1},c_k,a_k,a_{k+1})\\
\B_{g,g}H_1&(c_{k-2},a_{k-2},a_{k-1},c_k,a_k,a_{k+1})\\
\B_{g,g}F_2&(c_{k-3},c_{k-1},a_{k-1},c_k,a_k,a_{k+1})\\
\B_{g,g}H_2&(c_{k-2},a_{k-2},c_{k-1},c_k,a_k,a_{k+1}).
\end{array}
\]
For \(b>q\), the two alphabets consist of distinct independent
generators, since \(0<q<g\).  For \(a_0^3a_5c_2c_4\), the three
\(a_0\)-positions are forced to be
\((3,1),(5,3),(6,5)\) in the first source and
\((4,2),(5,3),(6,5)\) in the first mate.  On each side \(a_5\)
then forces \((2,6)\).  The two remaining factors in the source are
\(c_4,c_2\); in the mate the possibilities are \(c_3^2\) and
\(c_1c_5\).  Thus only the source contributes, with column word
\((2,6,1,4,3,5)\).  

For \(a_0^2a_5c_1^2c_4\), the \(a_0\)-positions are
\((5,3),(6,5)\) in the second source and \((4,2),(6,5)\) in
its mate.  Again \(a_5\) forces \((2,6)\).  The factor \(c_4\)
then forces \((1,2)\) in the source and \((1,3)\) in the mate.  
The remaining source choices give \(c_0c_2\) or a zero entry,
whereas the mate gives \(c_1^2\), with column word
\((3,6,1,2,4,5)\).  Both contributing permutations have six
inversions.  Consequently
\begin{equation}\label{eq:five-border-separators}
 \begin{aligned}
 [a_0^3a_5c_2c_4](s_{\B_{g,g}F_1}-s_{\B_{g,g}H_1})&=1,\\
 [a_0^2a_5c_1^2c_4](s_{\B_{g,g}F_2}-s_{\B_{g,g}H_2})&=-1.
 \end{aligned}
\end{equation}
At \(b=q\), only \(c_0\) becomes a unit.  The displayed monomials have
ordinary degree six, so no term containing that unit can affect them.
Alternatively \(\Gamma\) replaces the parameters by
\((2b+1,2g,2q)\) and the margins by \((2u,2v)\).
\end{proof}

\begin{proof}[Proof of \Cref{thm:classification}(iv)]
The five-row base is the algebraic \Cref{thm:five-base}.  
Work zero-free for larger rows.  If a six-row source has connected maximal
cofactor, its complete class is given by that base.  Cross-orbit borders
are excluded by \Cref{lem:five-exceptional-borders}, and same-orbit
margin swaps by \Cref{lem:five-endpoint}.

A disconnected \((n-1)\)-row cofactor with three components would use
at least \(2n-5\) negatives.  With only five negatives and \(n\ge6\),
it has two components, using at least \(n-2\) negatives.  The only
possibilities are a six-row parent with split \(1+4\) and one internal
negative, or a seven-row parent with split \(1+5\) and no internal
negative, including their reversals.  The first is
\Cref{lem:singleton-rigidity}; the second is \Cref{thm:column}.
For \(n\ge8\) no split is possible.

All other larger-row steps have a connected child.  Induction identifies
its literal diagram up to rotation, and also gives rigidity of its
endpoint extensions, which retain exactly five negatives.
\Cref{lem:border} then closes the step.  Homogenization supplies all
common-beta sources.  This proves the higher-row assertion from
\Cref{thm:five-base} and the endpoint and border lemmas, with no bound
on area or row count.
\end{proof}

\section{Staircase components and residual matching}\label{sec:staircase}

\subsection{Closure and the base statement}

The negative partition of \(\mathcal T_n\) is \((3,2,1)\).
By \Cref{lem:negative-partition} this pattern is stable under rotation.
Stability under arbitrary Schur equality requires the stronger overlap
conditions in the following lemma.

\begin{lemma}\label{lem:staircase-closure}
Let \(D\) be connected with \(n\ge5\) rows and full negative set
\(\mathcal T_n\) of \eqref{eq:staircase-set}.  Every connected
Schur-equal mate has \(n\) rows and the same negative set.  Common betas
are allowed.
\end{lemma}

\begin{proof}
Use the overlap criterion in the proof of
\Cref{lem:five-mate-closure}: an entry at row distance \(k-1\) is
negative exactly when the corresponding \(k\)-row intersection has size
at most \(k-2\).  In a staircase, every such intersection has size
at least \(k-1\) for \(2\le k\le n-3\), and at most \(k-2\)
for \(k\ge n-2\).  The invariant overlap partitions, padded using
the common row count, preserve both universal conditions.  Thus they
recover the whole staircase mask, not merely its negative count.
\end{proof}

At five rows the staircase is exactly an upper Hessenberg board.  In
profile coordinates its condition is
\begin{equation}\label{eq:hessenberg-domain}
 z_i\ge0,\qquad r_i>z_{i-1}+z_i,\qquad z_0=z_5=0.
\end{equation}
Zero-freeness is exactly \(z_i\ge1\), a condition we call \emph{thickness}.
For general order \(k\), an upper Hessenberg board has the interval expansion
\begin{equation}\label{eq:interval-expansion}
 s_{(r;z)}=\sum_{S\subseteq\{1,\ldots,k-1\}}(-1)^{|S|}
       \prod_{i\in S}h_{z_i}\prod_{[a,b]\in\mathcal I(S)}h_{w_{a\cdots b}},
 \qquad
 w_{a\cdots b}=\sum_{i=a}^b r_i-\sum_{i=a}^{b-1}z_i.
\end{equation}
Here \(\mathcal I(S)\) joins adjacent rows exactly across the edges in
\(S\).  The interval-block argument used for
\eqref{eq:five-open-decomposition} applies in every order and proves
this formula.  At five rows there are sixteen indexed summands;
coinciding monomials are collected with their signs.
A thick profile is \emph{separated} when
\[
 \{z_1,z_2,z_3,z_4\}\cap
 \{w_{a\cdots b}:1\le a\le b\le5\}=\varnothing;
\]
otherwise it has a \emph{contact}.  No separation hypothesis is imposed
on the classification below.

\begin{theorem}[Staircase base]\label{thm:staircase-base}
For a connected five-row staircase, with common betas allowed, every
connected mate is its rotation unless the pair is in one of the following
complete four-member classes.  Add \(d\) to every entry of both words in
each displayed profile:
\begin{equation}\label{eq:staircase-families}
\begin{array}{c|c|c}
 &F:\text{ rows; gaps}&H:\text{ rows; gaps}\\\hline
\mathrm{M1}&(a,b,a,a+b,b;t,0,t,t)&(a,a+b,b,a,b;t,t,0,t)\\
\mathrm{M2}&(a,b,a,a+b-t,b;0,t,0,0)&(a,a+b-t,b,a,b;0,0,t,0)\\
\mathrm{M3}&(a,b,b,b,a+b;0,0,0,a)&(a,b,a+b,b,b;0,a,0,0)\\
\mathrm{M4}&(a,b,b,a+b,a+b;0,0,a,a)&(a,b,a+b,a+b,b;0,a,a,0).
\end{array}
\end{equation}
For M1 and M2 the integer parameters satisfy
\(d,t\ge0\), \(a,b>t+d\), \(a\ne b\).
For M3 and M4 they satisfy \(d\ge0\), \(a\ge1\), \(b>a+d\).
Each class is \(\{F,F^*,H,H^*\}\).  M1 and M2 coincide at \(t=0\).

\end{theorem}

\subsection{Safe cuts and specialized content}

For a thick profile \(D=(L;g)\), write \(\mathcal G(D)\) for its gap
multiset and \(\mathcal W(D)\) for its interval-index set.  The letters
\(L,g\) denote the same row and gap data as \(r,z\) in
\eqref{eq:interval-expansion}.  A contact is an element of
\(\mathcal G(D)\cap\mathcal W(D)\).  All generator specializations
keep these numerical indices fixed.  Throughout the staircase arguments, write
\[
 R_t(a,b)=h_ah_b-h_th_{a+b-t},
\]
\[
 T_t(a,b,c)=h_ah_bh_c-h_th_{a+b-t}h_c
           -h_th_ah_{b+c-t}+h_t^2h_{a+b+c-2t}.
\]
These are the two-row and constant-gap three-row polynomials on their
thick domains.  When \(z=\min\mathcal G(D)\), put
\(J_k(D)=[h_z^k]s_D\).

\begin{lemma}[The minimum-gap interface]\label{lem:staircase-minimum}
Let \(z=\min\mathcal G(D)\), with multiplicity \(m\).  Then
\(z\notin\mathcal W(D)\).  Cutting every \(z\)-seam gives thick
components \(D_j\), and contracting those seams gives a thick profile
\(D/z\), with
\[
 s_D\big|_{h_z=0}=\prod_j s_{D_j},\qquad
 [h_z^m]s_D=(-1)^m s_{D/z}.
\]
The first factorization recovers the component polynomials and their row
counts, but not their order.  At five rows the numbers of uncollected terms
in the respective minimum-gap layers are
\begin{equation}\label{eq:task-count}
 \#\{\text{terms of }h_z\text{-degree }j\}
   =\binom mj2^{4-m},\qquad
 \sum_{j=0}^m\binom mj2^{4-m}=16.
\end{equation}
For thick profiles \(A,B\), let \(u,v\) be their adjacent endpoint
rows and let \(\alpha,\beta\) be the gaps incident to those rows inside
\(A,B\), respectively; an absent gap is zero.  A positive joining gap
\(q\) is admissible here when
\[
 u>\alpha+q,\qquad v>q+\beta.
\]
Thus the joined profile \(A\cdot_qB\) is thick.  Replacing \(u,v\)
by \(u+v-q\) gives the thick contraction \(A\odot_qB\), and
\[
 s_{A\cdot_qB}=s_As_B-h_qs_{A\odot_qB}.
\]
Connectedness of the joined profile alone is not the hypothesis of this
identity.

This qualification is substantive.  For a connected but non-Hessenberg join,
\[
s_{(2,2,2;1,1)}
-\left(h_2s_{(2,2;1)}-h_1s_{(2,3;1)}\right)
=h_3^2-h_2h_4\ne0.
\]
The applications in the staircase proof retain the stronger joining inequalities.
\end{lemma}

\begin{proof}
Every row exceeds every incident gap and therefore exceeds \(z\).
Extending an interval adds \(L_i-g\), which is positive, so every interval
index exceeds \(z\).  In \eqref{eq:interval-expansion}, setting \(h_z=0\)
keeps every minimum seam, and taking its highest power contracts every
minimum seam.  No interval factor contributes another \(h_z\).
Contraction preserves the strict row inequalities: the merged row exceeds
the sum of its two remaining exterior gaps.  Connected irreducibility
\cite{Barekat2009CoincidencesCorrigendum} and ordinary degree recover the
component polynomials and their sizes.  Choosing which of the \(m\)
minimum seams to contract proves \eqref{eq:task-count}.  Partitioning the
interval expansion according to whether the joining seam is contracted
proves the final identity.
\end{proof}

Fix a positive integer \(q\).  A \emph{\(q\)-component} is a
component obtained after every \(q\)-seam has been cut.  Orient it to
put an endpoint adjoining a cut at the bottom.

Write \(A=(a_1,\ldots,a_n;r_1,\ldots,r_{n-1})\), and put
\(r_0=r_n=0\).  The conditions used below are
\[
 a_i>r_{i-1}+r_i\quad(1\le i\le n),\qquad
 r_i\ge1,\ r_i\ne q\quad(1\le i<n),\qquad
 a_n>q+r_{n-1}.
\]
For a singleton the last condition is \(a_1>q\).  An additional
inherited boundary condition at the other endpoint is retained whenever
present.  Put \(P_A=s_A|_{h_q=0}\).

\begin{lemma}[Monomial content and one primitive factor]
\label{lem:staircase-content}
Every \(q\)-component has a unique factorization
\[
 P_A=G_AK_A,
\]
where \(G_A\) is a monic monomial in its internal gap variables and
\(K_A\) is irreducible.  Every variable in \(G_A\) has index less than
\(q\).  Through four rows the complete content table is as follows.
For two rows,
\[
 G_{(a,b;r)}=\begin{cases}h_r,&a=q,\\1,&a\ne q.\end{cases}
\]
For three rows, put \(u=a+b-r\).  Then
\[
 G_{(a,b,c;r,s)}=
 \begin{cases}
 h_r,&a=q,\\
 h_s,&a\ne q\text{ and }[(b=q,\ r=s)\text{ or }(u=q,\ a=s)],\\
 1,&\text{otherwise}.
 \end{cases}
\]
For four rows,
\[
 G_{(a,b,c,d;r,s,t)}=
 \begin{cases}
 h_rh_s,&a=q,\ c=q,\ s=t,\\
 h_r,&a=q\text{ and not }(c=q,\ s=t),\\
 h_s,&a\ne q\text{ and at least one of }\left\{
 \begin{array}{l}
 b=q,\ r=s,\\ c=q,\ s=t,\\ a+b-r=q,\ a=s,
 \end{array}\right.\\
 1,&\text{otherwise}.
 \end{cases}
\]
The three conditions in the third line are alternatives; simultaneous
occurrence does not multiply their contributions.  A singleton has
\(G_A=1\).
\end{lemma}

\begin{proof}
Let \(T\) be the full interval index.  It exceeds every other index and,
by the inherited boundary, exceeds \(q\).  Thus
\[
 P_A=(-1)^{n-1}\left(\prod_i h_{r_i}\right)h_T+B,
 \qquad h_T\text{ does not occur in }B.
\]
Set \(G_A=\gcd(\prod_i h_{r_i},B)\).  After division, the polynomial is
primitive and linear in \(h_T\), so Gauss's lemma proves irreducibility.
This also covers \(B=0\).

For the content calculation at four rows use
\begin{gather*}
 u=a+b-r,\quad v=b+c-s,\quad w=c+d-t,\\
 x=a+b+c-r-s,\quad y=b+c+d-s-t,\quad
 T=a+b+c+d-r-s-t.
\end{gather*}
The exact expansion is
\begin{equation}\label{eq:staircase-four-expansion}
\begin{aligned}
 s_A={}&h_ah_bh_ch_d-h_rh_uh_ch_d-h_sh_ah_vh_d-h_th_ah_bh_w\\
 &+h_rh_sh_xh_d+h_rh_th_uh_w+h_sh_th_ah_y-h_rh_sh_th_T.
\end{aligned}
\end{equation}
The analogous three-row expansion has four terms.  Since
\(c>q+s\) in three rows, its only possible \(q\)-indices are
\(a,b,a+b-r\).  Setting an internal gap variable to zero in those four
terms gives exactly the three-row table.

If \(a=q\) in four rows, near-concatenation gives
\[
 P_A=-h_rP_{(q+b-r,c,d;s,t)}.
\]
The new first row exceeds \(q\), so the three-row table gives the first
two lines of the four-row table.

Suppose \(a\ne q\).  Every possible variable divisor is a gap variable.
A gap value \(r\ne s\) cannot divide \(P_A\): after setting \(h_r=0\),
the first-row cut leaves a three-row factor whose largest term survives.
If \(r=t\ne s\), both end seams are cut and the middle two-row factor
is still nonzero.  Similarly a value \(t\ne s\) cannot divide.
For the latter assertion the only exceptional largest-term case is
\(x=q\); the remaining prefix is
\[
 h_ah_bh_c-h_rh_uh_c-h_sh_ah_v.
\]
Its positive monomial cannot cancel a negative one because
\(u>\max(a,b)\), \(v>\max(b,c)\).  If setting \(h_t=0\) kills the
positive term, then \(a=t\) or \(b=t\), since \(c>t\); the first
negative term survives.  Hence only \(h_s\) can divide.

Cutting the middle seam gives
\[
 P_A\big|_{h_s=0}
 =\left(R_r(a,b)\big|_{h_q=h_s=0}\right)
  \left(R_t(c,d)\big|_{h_q=h_s=0}\right).
\]
The first factor vanishes exactly when \(b=q,r=s\) or \(a=s,u=q\).
The second vanishes exactly when \(c=q,s=t\), because
\(d>q+t\) and \(w>d\).  In each two-term factor a surviving positive
term cannot cancel its negative term, whose full interval exceeds both
rows.  This proves the divisibility conditions, including coincident gaps.

Only repeated gap values require a multiplicity check.  If
\(r=s\ne t,b=q\), the coefficient of the first power of \(h_s\) is
\[
 h_u(-\overline h_ch_d+h_th_w)
 +h_a(-h_vh_d+h_th_y).
\]
Compare \(u\) and \(y\).  In their tie, \(a=w\), and the coefficient
at that common largest index is
\(-\overline h_ch_d+2h_th_w\ne0\).
If \(s=t\ne r,c=q\), that coefficient is
\[
 h_d(-\overline h_ah_v+h_rh_x)
 +h_w(-\overline h_a\overline h_b+h_r\overline h_u).
\]
When \(a=s\) the surviving terms are positive.  Otherwise compare
\(x,w\); their tie gives \(u=d\) and the nonzero coefficient
\(2h_rh_u-h_a\overline h_b\), since \(u>\max(a,b)\).
For \(u=q,a=s=t\ne r\), the first coefficient is
\(h_d(h_b\overline h_c+h_rh_x)\ne0\).
Bars here mean specialization at both \(h_q=h_s=0\).
If all three gaps coincide, a surviving one-seam term has the common
negative sign of all such terms; hence the first coefficient is nonzero.
This proves the multiplicities in the table.

Finally, suppose a variable \(h_v\) with \(v>q\) divided \(P_A\).
Cut all internal \(v\)-seams after setting \(h_v=0\).  Every resulting
component has an endpoint adjacent to such a seam, so its full interval
exceeds \(v>q\).  Its largest-index coefficient is an internal gap
monomial containing neither \(h_v\) nor \(h_q\).  Each specialized
factor is therefore nonzero, contrary to the asserted divisibility.
Internal gaps cannot equal \(q\), proving the low-gap assertion.
\end{proof}

Normalize primitive factors by
\[
 \widehat K_A=(-1)^{n-1}P_A/G_A.
\]
Their largest-index coefficients are positive monomials with numerical
coefficient one.  Thus normalized associates are equal.  In the remaining
staircase arguments, a bar denotes specialization at \(h_q=0\), and
\[
 \varepsilon(P)=P\big|_{h_j=1\ (j\ne q)},\qquad
 \|P\|_{h,1}=\sum_\lambda\bigl|[h_\lambda]P\bigr|
 \quad\text{for }P\in\Q[h_j:j\ne q].
\]
The sum is over distinct commutative monomials after collection.
If all \(q\)-seams of a five-row profile are cut, its aggregate
extracted content is recovered directly from the restriction:
\[
 \prod_A G_A=
 \prod_{1\le j<q}h_j^{\nu_{h_j}(s_D|_{h_q=0})}.
\]
Here \(\nu_{h_j}\) denotes the exponent of the variable divisor.  A linear
primitive has full interval index greater than \(q\),
and a nonlinear primitive has no variable divisor.  This identifies the
aggregate content, not its ownership by individual components.

\subsection{Matching the recovered primitives}

\begin{lemma}[Endpoint coefficients of a four-row restriction]
\label{lem:staircase-endpoint-coefficients}
Let \(A=(a,b,c,d;r,s,t)\) be a \(q\)-component with \(a\ne q\),
and put
\[
 x=a+b+c-r-s,\qquad y=b+c+d-s-t,\qquad
 T=a+b+c+d-r-s-t,\qquad M=\max(x,y).
\]
The two largest indices occurring in \(P_A\) are \(T,M\), and
\[
 -\frac{[h_M]P_A}{[h_T]P_A}
 =\begin{cases}
 h_a/h_r,&y>x,\\
 h_d/h_t,&x>y,\\
 h_a/h_r+h_d/h_t,&x=y.
 \end{cases}
\]
The strict case recovers one ordered endpoint row/gap pair; the tie
recovers the unordered two such pairs, with multiplicity.  The same
ratio is obtained from \(\widehat K_A\).
\end{lemma}

\begin{proof}
Every proper interval index is at most \(M\), and all shorter
intervals and gap indices are strictly below \(M\).
Moreover \(T>M\), \(y>d>q\), and \(x=q\) therefore implies
\(y>x\).  The coefficients in
\eqref{eq:staircase-four-expansion} are
\[
 [h_T]P_A=-h_rh_sh_t,\qquad
 [h_M]P_A=
 \mathbf1_{M=x}h_rh_sh_d+\mathbf1_{M=y}h_sh_th_a.
\]
Both displayed contributions survive and have the same sign.
Each reduced Laurent monomial identifies its ordered pair because
\(a>r\) and \(d>t\); equal contributions retain coefficient two.
The content has all its indices below \(q<M\), so division by it
does not change the ratio.  Finally, if \(x>y\), then
\(a-r>d-t>q\).  Rotation still satisfies the inherited boundary
condition and puts the recovered endpoint first.  This also holds at
a tie, when both endpoint excesses exceed \(q\).
\end{proof}

\begin{lemma}[Residuals through three rows and cross-size matching]
\label{lem:staircase-small-residuals}
Among \(q\)-components with at most four rows, the quadratic normalized
primitives are exactly
\begin{align*}
 B_r(a,c)&=h_rh_{a+c-r}-h_ah_c,\\
 M_{q,r}(a,c)&=h_rh_{a+c+q-2r}
              -h_{a+q-r}h_c-h_ah_{c+q-r},\\
 I_{q;r,s}(c)&=h_{q-s+r}h_c
              -h_sh_{q+c-2s+r}+h_rh_{q+c-s}.
\end{align*}
Their matching keys are respectively \((r,\{a,c\})\),
\((r,\{a,c\})\), and \((r,s,c)\).  Different types cannot match.
Their possible origins, through four rows, are
\[
\begin{array}{c|c|c}
\text{component condition}&G_A&\widehat K_A\\\hline
(a,b;r),\ a\ne q&1&B_r(a,b)\\
(q,b,c;r,s)&h_r&B_s(q+b-r,c)\\
(a,q,c;r,r),\ a\ne q&h_r&M_{q,r}(a,c)\\
(s,q-s+r,c;r,s)&h_s&I_{q;r,s}(c)\\
(q,b,q,d;r,s,s)&h_rh_s&M_{q,s}(q+b-r,d).
\end{array}
\]
All origins retain their thick and inherited boundary conditions.
A linear primitive \(h_T\) comes from the singleton \((T)\), with
\(T>q\), or from \((q,T-q+r;r)\), with \(1\le r<q\) and
\(T>2q\).  The stronger bound belongs only to the two-row origin.

A primitive three-row cubic determines its profile up to rotation.
A three-row primitive cubic and a four-row primitive cubic can match
only by endpoint contraction: the four-row component is
\[
 (q,b,c,d;r,s,t),
\]
and the three-row component is
\((q+b-r,c,d;s,t)\) or its rotation.  The four-row content is \(h_r\).
\end{lemma}

\begin{proof}
Sending every surviving generator to one gives the quadratic values
\(0,-1,1\), respectively.  Their largest coefficients recover \(r\).
The negative monomial of \(B\) recovers its unordered row pair.
For \(M\), after removing the largest-index monomial, the smallest
index is \(\min(a,c)\) and the largest is \(\max(a,c)+q-r\).
For \(I\), the smaller index in its negative monomial is \(s\), and
its largest index is \(q+c-s\).  The content table proves the origin
list.  These readings allow coincident same-sign monomials.

For a primitive three-row cubic the possible contact types are untouched,
middle-row contact \(b=q,r\ne s\), and interval contact
\(a+b-r=q,a\ne s\).  Their pairs (augmentation, norm) are
\((0,4),(-1,3),(1,3)\).  An untouched cubic is a full Schur function,
so \Cref{lem:small-rows} applies.  For a middle contact, the unequal gap
labels distinguish the two negative monomials.  After ordering the gaps,
a nontrivial match would require both
\[
 a'=c-(q-r),\qquad a'=c+(q-s),
\]
which is impossible.  For an interval contact, its surviving diagonal
recovers \(c>q\) and \(\{a,b\}\), where \(a,b<q\).  Then
\(r=a+b-q\), the largest coefficient recovers \(s\), and the negative
monomial recovers \(a\).  This proves cubic rigidity.

For a four-row cubic beginning with \(q\), the endpoint identity in
\Cref{lem:staircase-content} proves the asserted contraction rule.
Suppose instead \(a\ne q\).  The content table leaves
\(b=q,r=s\), \(c=q,s=t\), or \(u=q,a=s\), where \(u=a+b-r\).
Expanding \eqref{eq:staircase-four-expansion} after division by the
content gives the following norms:
\[
\begin{array}{c|c}
 b=q,\ r=s&6,\text{ or }5\text{ if }c=q\\
 c=q,\ s=t&6,\text{ or }5\text{ if }b=q\text{ or }u=q\\
 u=q,\ a=s&6,\text{ or }5\text{ if }c=q\text{ or }v=q.
\end{array}
\]
There is one exception to this display:
\[
 c=q,\quad s=t,\quad d=r,\quad a=q+r-s.
\]
To check the exception and completeness of the norm reading, in the
right-middle case the six normalized terms are
\begin{align*}
 &h_ah_{b+q-s}h_d+h_a\overline h_bh_{q+d-s}
 -h_rh_{a+b+q-r-s}h_d\\
 &\qquad-h_r\overline h_{a+b-r}h_{q+d-s}
 -h_sh_ah_{b+q+d-2s}+h_rh_sh_{a+b+q+d-r-2s}.
\end{align*}
The only possible opposite-sign coincidence pairs the first term with
the fourth.  Since \(a,b+q-s>r\), it forces \(d=r\); then
\(a+b-r>a\) forces \(a=q+r-s\).

For the left-middle case \(b=q,r=s\), the normalized primitive is
\[
 \overline h_ch_uh_d+h_ah_vh_d-h_rh_xh_d
 -h_th_uh_w-h_th_ah_y+h_rh_th_T.
\]
The full-interval term is private.  Canceling shared factors in the
other positive--negative comparisons excludes all but
\(h_ah_vh_d=h_th_uh_w\): one uses
\(x>\max(u,c,a,v)\), \(w>\max(c,d)\),
\(y>\max(v,d)\), and \(c,d,v>t\).
In the remaining comparison \(t\) must match \(a\), because
\(v,d>t\).  Then \(w>d\) forces \(w=v\), giving
\(d=q+t-r\), contrary to \(d>q+t\).
Only the first term can be deleted by an additional \(q\)-contact,
namely \(c=q\).

For interval extraction \(u=q,a=s\), the normalized primitive is
\[
 -h_b\overline h_ch_d+h_s\overline h_vh_d+h_th_bh_w
 -h_rh_xh_d-h_sh_th_y+h_rh_th_T.
\]
Here \(r<s<b<q\), \(c>s+t\), \(d>q+t\), and
\(w>\max(x,d)\).  The full-interval term is private.
For the term \(h_sh_vh_d\), canceling its shared factor with each
negative term would respectively require a row \(b\) or \(c\)
to equal \(s\), would equate the distinct least indices \(s,r\),
or would equate \(\{v,d\}\) with \(\{t,y\}\).
For \(h_th_bh_w\), the comparisons with the first and fifth terms
would require \(\{c,d\}=\{t,w\}\) or
\(\{b,w\}=\{s,y\}\); the comparison with the fourth is
impossible because \(w>\max(r,x,d)\).
The only possible deletions are \(c=q\) or \(v=q\), not both.
Thus these two cases have the stated norms even when same-sign
monomials coincide.

On the exceptional locus put \(\delta=q-s\).  The normalized cubic is
\[
 h_a^2h_b-h_r^2h_{b+2\delta}
 -h_sh_ah_{b+r+q-2s}+h_rh_sh_{b+r+2q-3s}.
\]
Its norm is four, but it cannot be a three-row cubic.  Its augmentation
would force the untouched type, whose row multiset would be
\(\{a,a,b\}\).  The negative monomial
\(h_r^2h_{b+2\delta}\) contains none of those row values:
\(r<a,b<b+2\delta\) and \(a<b+2\delta\).
Every negative three-row term has an endpoint-row factor.  This
contradiction disposes of the exception and proves the cross-size result.
\end{proof}

\begin{theorem}[Size-preserving four-row residual matching]
\label{thm:staircase-residual-matching}
Let \(A,B\) be four-row \(q\)-components whose normalized primitive
factors have the same degree, either three or four.  Their factors agree
exactly in the corresponding cases below.  The quadratic case is already
covered by \Cref{lem:staircase-small-residuals}.

For cubic primitives, either both components are endpoint lifts
\[
 (q,U-q+r,c,d;r,s,t),\qquad 1\le r<q,
\]
of the same primitive three-row profile \((U,c,d;s,t)\), up to rotation
of that contracted profile, or the components are identical, rotations,
or members of the rotation quadruple
\[
 (a,q,a,a+2q-2z;z,z,z)
 \sim(a,a+q-z,q,a+q-z;z,z,z),
 \quad a,q>2z,\quad a\ne q.
\]
The endpoint lifts may use different values of \(r\).  All inherited
boundary conditions remain in force.

For quartic primitives, the possibilities are identity, rotation,
a full four-row equality from \eqref{eq:four-family} untouched by
\(h_q=0\), or the rotation quadruple
\begin{equation}\label{eq:staircase-residual-family}
\begin{aligned}
 A&=(a,q,a,a+2q-r-s;r,s,r),\\
 B&=(a,a+q-s,q,a+q-r;r,r,s),
\end{aligned}
\qquad a,q>r+s,\quad r\ne s,\quad a\ne q.
\end{equation}
For the full-four-row alternative, ``untouched'' means that \(q\)
is not an interval index of either component.
\end{theorem}

\begin{proof}
Use the indices \(u,v,w,x,y,T\) of
\eqref{eq:staircase-four-expansion}.  For \(a\ne q\),
\Cref{lem:staircase-endpoint-coefficients} recovers \(a,r\) after
orienting a strict comparison with \(y>x\), and both endpoint pairs
at a tie.

For the remaining readings, let \(e_j=\xi^j\) for \(j\ne q\),
\(e_q=0\), and define
\[
 \mathcal D(P)=\sum_{j\ne q}
 \left.\frac{\partial P}{\partial h_j}\right|_{h_i=1\ (i\ne q)}e_j,
 \qquad \mathcal D_0=e_x+e_y-e_v-e_T.
\]
The contact set records exactly which of \(a,b,c,u,v,x\) equal \(q\).
Expansion gives
\[
\begin{array}{c|l}
\text{contact set}&\mathcal D(P_A)\\\hline
\varnothing&\mathcal D_0\\
b&\mathcal D_0+e_t+e_w-e_c-e_d\\
c&\mathcal D_0+e_r+e_u-e_a-e_b\\
u&\mathcal D_0+e_c+e_d-e_t-e_w\\
uc&e_x+e_y-e_T-e_a-e_b-e_t-e_v-e_w\\
bc&e_r+e_t+e_x+e_y+e_u+e_w-e_v-e_T\\
uv&e_x+e_y-e_T+e_a+e_s+e_c+2e_d-e_t-e_w.
\end{array}
\]
If \(G_A=h_g\), then
\[
 \mathcal D(\widehat K_A)
 =-\mathcal D(P_A)-\varepsilon(\widehat K_A)e_g.
\]
The possible nonendpoint contact sets and augmentations of \(P_A\) are
\[
\begin{array}{c|c}
\varnothing,b,c,u&0\\ v,uv,bc&1\\ x,uc&-1.
\end{array}
\]
These are identities of collected polynomials, not distinct-index
approximations.

\emph{Cubic primitives.}
The endpoint class and its separation from nonendpoint components follow
from \Cref{lem:staircase-small-residuals}.  For nonendpoint cubics the
content table and the preceding derivative give
\[
\begin{array}{c|c|c|l}
\text{contact}&\text{extraction condition}&\varepsilon(\widehat K_A)
 &\text{least nonzero derivative coefficient}\\\hline
b&r=s&0&-1\text{ at }t\\
c&s=t&0&-1\text{ at }r\\
u&a=s&0&+1\text{ at }t\\
bc&r=s\text{ or }s=t&-1&-1\text{ at the nonduplicated outer gap}\\
uv&a=s&-1&+1\text{ at }t.
\end{array}
\]
In the all-equal \(bc\) case the indicated index is that common gap.
For \(uc\), the augmentation is one and the part below \(q\) is
\[
 \mathcal D(\widehat K_A)_{<q}=e_a+e_b+e_t-e_s.
\]
Thus interval-extraction and middle-row-extraction mechanisms cannot
cross-match.

For the \(uc\) case, the endpoint reading is strict and recovers
\(a,r\), hence \(b=q+r-a,c=q\).  Subtracting \(e_a+e_b\) from
the low derivative leaves \(e_t-e_s\).  Its nonzero case recovers both
gaps; if it is zero, the largest-index coefficient \(h_rh_t\) recovers
their common value.  Then \(d=T-2q+s+t\).  This covers both extraction
conditions \(a=s\) and \(s=t\), including their intersection.

For interval extraction \(u=q,a=s\) outside that case, the strict
endpoint reading recovers \(a=s,r\).  If also \(v=q\), then
\(b=q-s+r,c=2s-r\), and \(T\) fixes \(d\).
For an isolated \(u\)-contact, the least derivative recovers \(t\).
The restored polynomial \(P_A=h_s\widehat K_A(-1)\) has its surviving
diagonal as unique minimum-square monomial, recovering the row multiset.
Only \(c,d\) can exchange.  Their difference is
\[
 h_s\left[
 h_d(h_rh_{q+c-s}-h_ah_{b+c-s})
 -h_c(h_rh_{q+d-s}-h_ah_{b+d-s})\right].
\]
If \(c>d\), its largest index is \(q+c-s\), with nonzero
coefficient \(h_sh_rh_d\); the reverse inequality is symmetric.

For middle extraction with a nonsquare largest coefficient
\(h_\alpha h_\beta\), the endpoint gap \(r\) and the least derivative
index identify whether the duplicated gap is on the left or right.
If that least index differs from \(r\), the contact is on the left,
\(r=s\), and the index is \(t\).  Otherwise the contact is on the
right, \(s=t\), and the other factor is the duplicated gap.  Thus
\(G_A\) and the ordered gap triple are recovered.  In the left-middle
case with \(c\ne q\), direct expansion gives
\[
 \widehat K_A=-h_dM_{q,r}(a,c)
 -h_th_uh_w-h_th_ah_y+h_rh_th_T.
\]
Here \(u=a+q-r\), \(w=T-a-q+2r\), and \(y=T-a+r\)
are known.  Subtracting the last three terms therefore recovers
\(-h_dM_{q,r}(a,c)\).  Its quadratic is primitive and linear in its
largest generator \(h_x\), with coefficient \(h_r\): all indices in
its other two terms exceed \(r\).  Gauss's lemma proves irreducibility
using only \(a,c>r\) and \(q>2r\); the additional boundary
\(c>q+r\) is not required here.  Hence \(d\), then \(c\), is
recovered.  In the right-middle case put
\(\eta=a-r>0\).  The exact derivative identity
\[
 \mathcal D(P_A)-e_r+e_a-e_y+e_T
 =(\xi^\eta-1)(\xi^b+\xi^v)
\]
recovers \(b,v\), since \(v>b\), and then \(s,t,d\).
The \(bc\) case is immediate once its gaps are recovered.
At an endpoint tie the two endpoint pairs are already known, and the
same readings fix the middle data.  This includes the norm-four locus
\(c=q,s=t,d=r,a=q+r-s\): its derivative and the last identity are
unchanged by the cancellation.

The remaining cubic case has all gaps equal to \(z\) and a single
middle-row contact.  For this calculation, rotate to
\((A,q,C,D;z,z,z)\), with \(A,C,D\ne q\).  This rotation preserves
the specialized determinant, but need not preserve \(D>q+z\).
We use only \(A,D>z\), \(C,q>2z\), and the fact that the
second row is the only index equal to \(q\).  No \(q\)-component
lemma is applied in this temporary orientation.  Put \(\delta=q-z\).  Direct expansion gives, at
\(h_z=0\),
\[
 h_D(h_{A+\delta}h_C+h_Ah_{C+\delta}).
\]
For \(A\ne C\), the quadratic is primitive and linear in its largest
generator, so it recovers \(\{A,C\}\), and the linear factor recovers
\(D\).  Swapping \(A,C\) changes the remaining primitive by
\begin{align*}
-h_z\bigl[&h_{A+\delta}h_{C+D-z}+h_Ah_{C+D+\delta-z}\\
          &-h_{C+\delta}h_{A+D-z}-h_Ch_{A+D+\delta-z}\bigr].
\end{align*}
For \(A\ne C\), its private largest index is
\(\max(A,C)+D+\delta-z\); the strict bound needed here is exactly
\(D>z\), not the discarded inherited boundary.
For \(A=C=\lambda\), the specialization is
\(2h_\lambda h_{\lambda+\delta}h_D\).  Another decomposition into
such a shifted pair is possible only when
\(D=\lambda+2\delta\), replacing the repeated row by
\(\lambda+\delta\) and the endpoint by \(\lambda\), or conversely.
These are precisely the constant-gap quadruple in the theorem, retaining
only the orientations satisfying the original component hypotheses.

\emph{Quartic primitives.}
Now \(G_A=1\).  If no row equals \(q\), the diagonal monomial
survives as the unique minimum-square monomial, with coefficient one.
Merging two neighboring interval blocks of indices \(U,V\) across a
gap \(g\) increases square weight by
\[
 g^2+(U+V-g)^2-U^2-V^2=2(U-g)(V-g)>0.
\]
If exactly one row is deleted at \(h_q=0\), every minimum surviving
term has negative sign.  This distinguishes \(x\) from \(uc\),
and untouched or \(u\)-contact profiles from single-middle contacts.
For \(bc\), the two minimum candidates are
\[
 -h_sh_ah_{2q-s}h_d,\qquad h_rh_th_uh_w.
\]
They cannot cancel because quartic primitivity requires \(s\ne r,t\).
A negative minimum, or a tie, cannot match a surviving diagonal.
If the positive candidate is uniquely minimal, matching a \(v\)- or
\(uv\)-contact would force its middle-row multiset to be \(\{r,t\}\),
the two factors below \(q\).  Its thickness inequality
\(L_2+L_3>g_1+2g_2+g_3\) contradicts that equality.

For \(x=q\), the endpoint reading recovers \(a,r\), and the diagonal
recovers the unique row \(d>q\).  Then \(t=q+d-T\) and the gap
multiset recover the gaps.  After subtracting the fixed terms, equality
reduces to equality of
\[
 R_r(a,b)R_t(c,d),
 \qquad
 P_A=R_r(a,b)R_t(c,d)-h_sh_aR_t(v,d)-h_rh_sh_th_T.
\]
Quadratic factorization fixes the ordering; a crossed match requires
\(r=t,a=d\), impossible because \(a<q<d\).
For \(uc\), the endpoint reading gives \(a,r\), hence
\(b=q+r-a,c=q\).  The low derivative
\(-e_a-e_b-e_t\) recovers \(t\), and then \(s,d\).

For \(v=q\), including \(uv\), the surviving diagonal and strict
endpoint reading recover \(a,r,d\), since \(b,c<q<d\).  Then
\(t=a+q+d-r-T\), the gap multiset, and the row multiset recover
\(s,\{b,c\}\).  Removing fixed terms leaves
\[
 \overline{R_r(a,b)}R_t(c,d)
 =\overline{R_r(a,c)}R_t(b,d).
\]
Only the first quadratic can specialize.  If only one side does so, the
factor degrees are \(1+1+2\) against \(2+2\), impossible.  Otherwise
quadratic matching gives identity or \(r=t,a=d\), which is rotation.
The tied endpoint case already has both endpoint pairs.  For \(bc\),
the low derivative is \(e_r+e_t\); together with the endpoint reading
it recovers both endpoints and outer gaps, and the middle rows are \(q,q\).

An untouched polynomial has least derivative index \(v\), exceeding
every gap.  An isolated \(u\)-contact has least derivative index \(t\),
with coefficient \(-1\), so these cases cannot mix.  For the latter,
the endpoint reading and derivative recover \(a,r,t\), then \(s,b\).
The diagonal recovers \(\{c,d\}\), and the displayed interval-contact
swap defect excludes a nontrivial exchange.  For two untouched profiles,
\Cref{prop:four-rows} supplies exactly the full-function alternatives.

It remains to compare single-middle contacts.  In the strict orientation,
\(a,r\) are recovered.  For \(b=q\), the least derivative index is
\(t\), with coefficient one; for \(c=q\), it is \(r\), also with
coefficient one.  Two left-middle profiles therefore have the same gap
triple.  Removing their fixed terms leaves
\[
 h_d(h_rh_sh_x-h_rh_uh_c-h_sh_ah_v).
\]
The cubic factor is primitive and linear in \(h_x\): here
\(r\ne s\), \(u>q>s\), \(v>q>r\), and \(a>r,c>s\).
It is irreducible, so \(d\), then \(c\), are fixed.
Two right-middle profiles are fixed by the preceding derivative identity
\((\xi^{a-r}-1)(\xi^b+\xi^v)\).

For a mixed comparison, write the source as
\((a,q,c,d;r,s,r)\).  Equality of the least derivative indices forces
its last gap to be \(r\).  The mate is \((a,b,q,e)\), with gaps
\((r,s,r)\) or \((r,r,s)\), and
\(c+d=b+e\).  Remove the common full and suffix terms.
For the first gap order the only possible largest residual indices and
coefficient contributions are
\[
\begin{array}{c|cc}
&\text{first index}&\text{second index}\\\hline
\text{source}&a+q+c-r-s:\ h_rh_sh_d&c+d-r:\ h_r^2h_{a+q-r}\\
\text{mate}&a+b+q-r-s:\ h_rh_sh_e&q+e-r:\ -h_rR_r(a,b).
\end{array}
\]
A strict second mate maximum has augmentation zero, unlike the positive
source contribution.  At a mate tie, canceling its negative monomial
requires \(a=s,b=e\), and matching the result forces \(b=q\), contrary
to a single-middle contact.  The first mate maximum therefore dominates;
matching forces \(d=e,c=b\).  The remaining equality is
\[
 h_{a+q-r}R_r(c,d)=h_{d+q-r}R_r(a,c),
\]
forcing \(a=d\), contrary to the strict orientation.

For the swapped gap order the mate contributions are
\[
 a+b+q-2r:\ h_r^2h_e,
 \qquad q+e-s:\ -h_sR_r(a,b).
\]
If the second is maximal, including a tie, its coefficient is nonzero
modulo \(h_r\), unlike every source coefficient.  Thus the first
dominates.  The source's first contribution cannot match it modulo
\(h_r^2\), so the source's second dominates and
\(e=a+q-r\).  In the remainder, setting \(h_s=0\) and \(h_r=0\)
respectively gives
\[
 \{c,d\}=\{a,b+q-r\},\qquad
 \{q+c-s,d\}=\{b,q+e-s\}.
\]
The first matching \(c=a,d=b+q-r\) forces \(b=a+q-s\), giving
\eqref{eq:staircase-residual-family}.  The other forces \(a=b,d=a\),
after which the remaining terms require
\[
 \{2a+2q-2r-s,a\}=\{2a-r,a+2q-r-s\}.
\]
Both entries on the right exceed \(a\), a contradiction.
At a tie of the endpoint readings, the same derivative forces equal
outer gaps, and the known endpoint pairs make the mixed comparison a
rotation.

Finally, for \(U=a+q-r\), \(V=a+q-s\), and
\(D=a+2q-r-s\), direct expansion gives the full identity
\[
 s_{(a,q,a,D;r,s,r)}-s_{(a,V,q,U;r,r,s)}
 =h_q\bigl(h_aR_r(a,D)-h_UR_r(a,V)\bigr).
\]
It proves sufficiency of the quartic family and, at \(r=s\), of the
nonendpoint cubic family.  The endpoint identities prove the remaining
sufficiency statements.
\end{proof}

These matching results recover possible components and their contractions.
They do not fix the order or orientations of components in a parent.

\section{Staircase assembly and higher-row exclusion}\label{sec:staircase-assembly}

\subsection{Compatibility across cuts}

For the following illustration put \(L=w_{1234}\),
\(R=w_{2345}>L\), \(T=w_{345}\), and \(\kappa=\max(L,T)\).
The third supported label has collected coefficient
\begin{equation}\label{eq:third-marker}
\begin{aligned}
 [h_\kappa]s_D={}&-\mathbf1_{\kappa=L}
 h_{z_1}h_{z_2}h_{z_3}h_{r_5}\\
 &+\mathbf1_{\kappa=T}h_{z_3}h_{z_4}
 (h_{r_1}h_{r_2}-h_{w_{12}}h_{z_1}).
\end{aligned}
\end{equation}
The binomial is nonzero because \(w_{12}>r_1,r_2,z_1\).
At \(L=T\), both contributions are retained and their sum has
augmentation \(-1\).  No ordering regime is assumed invariant.

\begin{example}\label{ex:mixed-stratum}
The equal profiles
\[
 (2,4,4,4,5;1,1,1,2),\qquad(2,4,5,4,4;1,2,1,1)
\]
have \((L,R,T)=(11,13,10)\) and \((11,13,11)\), respectively.
On the second side a positive and a negative third-label term cancel,
leaving \(-h_5h_1^3\), the coefficient on the first side.
The proofs below use the collected coefficients and permit this change
of regime.
\end{example}

\begin{lemma}[End-to-interior seam matching]
\label{lem:staircase-seam-crossing}
Suppose a gap \(q\) occurs once in the common gap multiset of two equal
thick five-row profiles.  If its seam is at an end on one side and in the
interior on the other, then, after independent rotations, the pair is
\[
 (q,b,b,b,u;z,z,z,q)
 \sim(q,b,u,b,b;z,q,z,z),
 \qquad u=b+q-z,\quad q>z,\quad b>q+z.
\]
No restriction on the other three gaps is assumed.  In particular this
classifies the \(4+1\) against \(2+3\) comparison, including its
endpoint-contact subcase.
\end{lemma}

\begin{proof}
Orient the \(q\)-cuts as \(4+1\) and \(2+3\), and write the first
profile as \((a,b,c,d,e;r,s,t,q)\).  The four-row primitive cannot
have degree four.  The two-row component on the other side must extract
its gap and become linear: otherwise that side has two nonlinear
primitives, whereas the first side has one.  Thus its first row is \(q\).

If the common nonlinear primitive is cubic,
\Cref{lem:staircase-small-residuals} forces \(a=q\) and identifies the
other three-row component with
\((u,c,d;s,t)\) or its rotation, where \(u=q+b-r\).
The low-gap content fixes the mate's first gap as \(r\), and its linear
primitive gives \(e=q+B-r\), where \(B\) is its second row.
The row multisets reduce to \(\{b,e\}=\{B,u\}\).  Since \(u>b\),
this forces \(B=b,e=u\).
If the common primitive is quadratic, the four-row table forces
\(a=c=q,s=t\).  Its type is \(M_{q,s}(u,d)\), whose three-row
origin has middle row \(q\) and both gaps \(s\).
The same content and row-multiset argument again gives \(B=b,e=u\).

Every candidate is therefore one of
\begin{align*}
 F&=(q,b,c,d,u;r,s,t,q),\\
 H_+&=(q,b,u,c,d;r,q,s,t),\\
 H_-&=(q,b,d,c,u;r,q,t,s),\qquad u=b+q-r,
\end{align*}
with
\[
 q>r,\quad b>r+\max(q,s),\quad c>s+t,\quad d>q+t,
 \qquad q\notin\{s,t\}.
\]
Put \(J=R_r(q,b)\), and
\[
 A_0=b+c+d+q-r-s-t,\quad B_0=b+c+d-r-t,
 \quad D_0=2b+c+q-2r-s.
\]
After division by \(h_q\), the only potentially largest indices in
\(s_F-s_{H_-}\) are \(A_0,B_0\), with coefficient contributions
\(-h_bh_sh_t\), \(h_bh_qh_t\).  All others are strictly below their
maximum.  Since \(A_0-B_0=q-s\ne0\), this difference is nonzero.

For \((s_F-s_{H_+})/h_q\), the complete list of possible largest
contributions is
\[
\begin{array}{c|c}
 A_0&-h_bh_sh_t\\ B_0&h_tJ\\ D_0&h_dh_rh_s.
\end{array}
\]
Coincident indices require addition of their displayed contributions.
If \(s>q\), then \(B_0>A_0\).  Its only possible cancellation is a
tie with \(D_0\), whose coefficient has all-ones value one and is
nonzero.  Hence \(s<q\).  Now \(A_0>B_0\), so equality forces
\(A_0=D_0\) and \(h_rh_d=h_th_b\).  Since \(b>r,d>t\),
\(t=r,d=b\).

After this substitution the possible largest contributions are
\[
\begin{array}{c|c}
 2b+c-2r&h_rJ\\
 2b+c-r-s&-h_sJ\\
 2b+q-2r&-h_rR_r(b,c).
\end{array}
\]
If \(s<r\), a necessary tie gives \(c=q+s-r\); the tied coefficient
has private largest factor \(h_{b+q-r}\) with coefficient \(h_sh_r\).
If \(s>r\), the necessary tie gives \(c=q\); after that cancellation,
the next largest index is \(2b+q-r-s\), with coefficient \(-h_sJ\).
Both are nonzero.  Thus \(s=r\).

Finally put \(\alpha=2b+q-2r\), \(\beta=b+c+q-2r\).  The entire
remaining difference is
\[
 -h_rh_q\left[
 h_\alpha R_r(b,c)-h_\beta R_r(b,b)
 +h_u(h_bh_{b+c-r}-h_ch_{2b-r})\right].
\]
If \(c>b\), the index \(\beta\) is uniquely largest; if \(c<b\),
\(\alpha\) is uniquely largest.  Both coefficients are nonzero.
For \(c=b\), the bracket vanishes identically.  This proves necessity
and sufficiency, with \(z=r\).
All the displayed contributions follow by grouping
\eqref{eq:interval-expansion}; the strict inequalities displayed above
bound every omitted index below the respective maximum.
\end{proof}

\begin{lemma}[The two-and-two cross-pattern]
\label{lem:staircase-a2}
Let \(q>z\ge1\).  A thick equality between profiles with gap words
\((z,z,q,q)\) and \((z,q,q,z)\) is, after rotating the second profile,
exactly
\[
 (q,b,b,u,u;z,z,q,q)
 \sim(q,b,u,u,b;z,q,q,z),
 \qquad u=b+q-z,\quad b>q+z.
\]
\end{lemma}

\begin{proof}
Write the source as \((a,b,c,d,e;z,z,q,q)\).  The safe \(z\)-cut and
three-row rigidity leave only
\[
 H_1=(a,c,d,e,b;z,q,q,z),\qquad
 H_2=(b,c,d,e,a;z,q,q,z).
\]
Their common validity gives \(c,e>q+z,d>2q\).
The \(q\)-cut on the source has a three-row component and two
singletons; the mate has two two-row components and a singleton.
A primitive cubic on the source cannot match.  The three-row content
table therefore forces \(a=q\) or \(b=q\).
If both equal \(q\), the source retains a quadratic while the mate
has only linear factors.  If \(b=q,a\ne q\), its primitive quadratic
is of type \(M\), with augmentation \(-1\); the mate's surviving
quadratic is of type \(B\), with augmentation zero.
Thus \(a=q,b\ne q\).

Put \(u=q+b-z\), \(v=q+c-z\).  For \(H_2\), the specialization
would give
\[
 h_eR_z(u,c)=h_{q+e-z}R_z(b,c),
\]
forcing the unequal linear indices \(e=q+e-z\), impossible.
For \(H_1\), it gives
\[
 h_eR_z(u,c)=h_vR_z(e,b).
\]
Hence \(e=v\) and \(\{u,c\}=\{e,b\}\).  Since \(u>b\),
\(c=b,e=u\).  Only the row \(d\) remains:
\[
 F_d=(q,b,b,d,u;z,z,q,q),\qquad
 H_d=(q,b,d,u,b;z,q,q,z).
\]
Set
\(\alpha=2b+q-2z\), \(\lambda=b+d-z\),
\(\mu=b+d-q\), and \(v_0=2b-z\).  The exact defect is
\[
 s_{F_d}-s_{H_d}
 =-h_zh_q\left[
 -h_\alpha R_q(b,d)+h_\lambda R_q(b,u)
 +h_u(h_dh_{v_0}-h_uh_\mu)\right].
\]
For \(d>u\), \(\lambda\) is uniquely largest, with coefficient
\(R_q(b,u)\ne0\).  For \(d<u\), \(\alpha\) is uniquely largest,
with coefficient \(-R_q(b,d)\ne0\).  At \(d=u\),
\(\lambda=\alpha,\mu=v_0\), and the bracket vanishes identically.
\end{proof}

\begin{lemma}[Two terminal cuts]\label{lem:staircase-four-plus-one}
An equality for which a unique \(q\)-seam is at an end on both sides is
rotation-trivial.  Equivalently, after putting both
\(q\)-seams last, equal five-row profiles are identical.
\end{lemma}

\begin{proof}
Write the profiles as \(A\cdot_q(e)\), \(B\cdot_q(f)\), with
\(e,f>q\).  Every four-row primitive has degree at least two.
Low-gap content and unique factorization of the restriction therefore
give
\[
 e=f,\qquad G_A=G_B,\qquad\widehat K_A=\widehat K_B.
\]
The full row multiset also recovers the component row multisets.
The residual family \eqref{eq:staircase-residual-family}, including its
constant-gap cubic case, cannot satisfy those row multisets: after
removing \(a,q\), it would require
\[
 \{a,a+2q-r-s\}=\{a+q-r,a+q-s\},
\]
whose right entries both exceed \(a\).

If the four-row components are full equals, near-concatenation requires
their \(q\)-contractions with \((e)\) to be full equals too.
A nonrotation pair from \eqref{eq:four-family} cannot persist: its gap
values and unchanged middle rows fix its parameters, while contraction
increases the last row by \(e-q>0\).  For a component and its rotation,
the contraction gap words are reverses, so they cannot give that
end/middle exception.  Direct matching makes the component symmetric;
matching by rotation would require \(e=q\).

For endpoint cubic lifts,
\Cref{lem:staircase-small-residuals,thm:staircase-residual-matching}
and the row multiset leave only
\[
 (q,b,c,u,e;r,s,t,q),\qquad
 (q,b,c,u,e;r,t,s,q),\qquad u=b+q-r.
\]
Put \(J=R_r(q,b)\) and
\begin{align*}
 \Psi(v)=h_v\bigl[&h_{b+c+q-r-v}R_r(b,e)
 -h_{b+c-v}R_q(u,e)\\
 &-h_{b+c+e-r-v}J\bigr].
\end{align*}
The full difference is \(h_q(\Psi(s)-\Psi(t))\).
If \(s<t\), its unique largest index is \(b+c+e-r-s\), with
coefficient \(-h_qh_sJ\ne0\); the other inequality is symmetric.
Thus \(s=t\).
For quadratic primitives the four-row origin is
\((q,b,q,d;r,s,s)\).  Its matching key recovers
\(s,\{b+q-r,d\}\), and content recovers \(r\).
The additional row pair \(\{b,d\}\) fixes \(b,d\) in order by
subtracting sums of squares.  All possibilities are now identical.
\end{proof}

\subsection{Interior matching}

\begin{lemma}[Adjacent endpoint exchange]\label{lem:staircase-adjacent-swap}
For thick profiles of at least three rows, exchanging just the first two
rows while keeping the gap word fixed preserves the Schur function only
when those rows are equal.  The same holds for the last two rows.
\end{lemma}

\begin{proof}
Two applications of near-concatenation give
\[
 s_{(a,b,c,\ldots;r,s,\ldots)}-s_{(b,a,c,\ldots;r,s,\ldots)}
 =-h_s\left[h_as_{(b+c-s,\ldots)}-h_bs_{(a+c-s,\ldots)}\right].
\]
When \(a>b\), the second determinant has a strictly larger full-interval
index than every index in the first summand and than \(a,b,s\).
Its largest coefficient is nonzero.  The other inequality is symmetric;
rotation proves the last-two assertion.
\end{proof}

\begin{lemma}[A contact-blind quadratic operator]
\label{lem:staircase-quadratic-injectivity}
Let \(e>q\ge1\) and \(K>q\) be integers, and put \(\eta=e-q\).
Let \(\mu:\Z\longrightarrow\Q\) have finite support satisfying
\[
 \operatorname{supp}\mu\subseteq\{s\in\Z:0<s<K-q\},\qquad
 H_N(\mu)=\sum_s\mu_s h_sh_{N-s}.
\]
Then
\[
 h_eH_K(\mu)-h_qH_{K+\eta}(\mu)=0
 \quad\Longrightarrow\quad\mu=0.
\]
\end{lemma}

\begin{proof}
Reduce modulo \(h_q\).  Since the quotient ring is an integral domain
and \(e\ne q\), \(H_K(\mu)\) is divisible by \(h_q\).
For every support point, \(K-s>q\).  Thus its only possible summand
containing \(h_q\) has \(s=q\); the reflected point \(K-q\)
is excluded.  Substitution and cancellation give
\[
 H_K(\mu)=\mu_qh_qh_{K-q},\qquad
 H_{K+\eta}(\mu)=\mu_qh_eh_{K-q}.
\]
If \(q\) is outside the support range, \(\mu_q=0\) already.
Otherwise \(K>2q\).  In the second equality, the coefficient of
\(h_qh_{K+\eta-q}\) is \(\mu_q\) on the left and zero on the
right: its reflected support point exceeds \(K-q\), and neither
\(e\) nor \(K-q\) equals \(q\).  Hence \(\mu_q=0\).
Both quadratic polynomials vanish.  Comparing their monomials, with
\(\mu\) extended by zero, gives
\[
 \mu_s=-\mu_{K-s},\qquad \mu_s=-\mu_{K+\eta-s}.
\]
At a reflection midpoint the coefficient equation gives \(\mu_s=0\).
Composition gives \(\mu_{s+\eta}=\mu_s\), and finite support implies
\(\mu=0\).
\end{proof}

The strict upper support bound cannot be dropped.  At \(K=q+e\),
the excluded measure \(\delta_e\) gives
\(h_e(h_eh_q)-h_qh_e^2=0\).  Below the bound follows from the
actual thickness inequalities; no support point equal to \(q\)
is excluded.

\begin{lemma}[Interior-to-interior reconstruction]
\label{lem:staircase-interior}
Let \(q>z\ge1\), and let
\[
 F=(a,b,c,d,e;z,q,z,z),\qquad
 H=(A,B,C,D,E;z,q,z,z)
\]
be thick five-row Hessenberg profiles.  If \(s_F=s_H\), then \(F=H\),
or, after interchanging the profiles,
\[
 \begin{aligned}
 F&=(p,m,p,t,m;z,q,z,z),\\
 H&=(m,p,m,t,p;z,q,z,z),
 \end{aligned}
 \qquad t=p+m-q,
\]
where \(p\ne m\) and \(p,m>q+z\).  The nonrotation alternative is
separated.  No band condition or specialized four-row matching theorem
is used in the proof.
\end{lemma}

\begin{proof}
Thickness means exactly
\[
 a>z,\qquad b,c>q+z,\qquad d>2z,\qquad e>z,
\]
and the capitalized counterparts.  Set
\[
 x=a+b-z,\quad y=c+d+e-2z,\qquad
 x'=A+B-z,\quad y'=C+D+E-2z.
\]
The minimum-cut coefficients are
\begin{align}
 J_0(F)&=h_ah_dh_eR_q(b,c),\label{eq:staircase-interior-zero}\\
 J_3(F)&=-R_q(x,y),\label{eq:staircase-interior-top}\\
 J_2(F)&=h_eR_q(x,c+d-z)
       +h_{d+e-z}R_q(x,c)+h_aR_q(b,y).
       \label{eq:staircase-interior-layer}
\end{align}
The three terms in the last line contract the pairs \(\{1,3\}\),
\(\{1,4\}\), and \(\{3,4\}\) of minimum seams.  The two monomials
of \(R_q(b,c)\) are relatively prime, since
\(q<\min(b,c)\) and \(b+c-q>\max(b,c)\).  Its monomial gcd and
quotient, and then the highest coefficient, therefore give
\begin{equation}\label{eq:staircase-interior-data}
 \{b,c\}=\{B,C\},\quad
 \{a,d,e\}=\{A,D,E\},\quad
 \{x,y\}=\{x',y'\}.
\end{equation}
All braces denote multisets.  The local bounds used for extraction are
\[
 x>\max(a,b,q),\qquad
 y>\max(c,d,e,c+d-z,d+e-z,q).
\]
In the notation \(\bar P=P|_{h_q=0}\), full equality also gives
\begin{equation}\label{eq:staircase-interior-specialization}
 \overline{s_F}=\overline{R_z(a,b)}\,
                    \overline{T_z(c,d,e)}
 =\overline{R_z(A,B)}\,\overline{T_z(C,D,E)}.
\end{equation}
This is the cut at the \(q\)-seam.  The proof uses this equality of
full specializations, not just equality of \(J_0,J_2,J_3\).

\emph{Ordered top pair.}
Suppose \(x'=x\), \(y'=y\).  If \(B=b\), then \(A=a,C=c\).
After subtracting the fixed last term in
\eqref{eq:staircase-interior-layer}, the remainder is
\[
 h_xH_{y+z}(\mu)-h_qH_{y+z+x-q}(\mu),\qquad
 \mu=\delta_e+\delta_{d+e-z},
\]
where \(\delta_v\) denotes unit mass at \(v\).
With \(K=y+z\), the two complements are
\[
 K-e=c+d-z>q,\qquad K-(d+e-z)=c>q.
\]
Thus \Cref{lem:staircase-quadratic-injectivity} applies to the
difference of the source and mate measures, with multiplier \(h_x\).
It recovers \(\{e,d+e-z\}\).  Since \(d>2z\), the smaller value
is \(e\); hence \(E=e,D=d\).  This includes \(x=y\).

Otherwise \(B=c\ne b,C=b,A=a+b-c\ne a\).
The singleton multiset makes \(A=d\) or \(A=e\), so
\(y=x+e-z\) or \(y=x+d-z\); in particular, \(y>x\).
Reduce \eqref{eq:staircase-interior-layer} modulo \(h_q\) before
extracting \(h_y\).  The other factors have indices below \(y\),
and
\[
 \bar h_ah_b=\bar h_Ah_c.
\]
Both sides cannot vanish without \(a=A=q\).  Monomial matching
therefore gives \(a=c=p,\ b=A=m\), with
\(p\ne m\) and \(p,m>q+z\).
Up to interchanging the profiles, the remaining singleton allocations are
\[
\begin{array}{c|c|c}
 &F&H\\\hline
 T_1&(p,m,p,m,t)&(m,p,m,p,t)\\
 T_2&(p,m,p,t,m)&(m,p,m,t,p)\\
 T_3&(p,m,p,m,t)&(m,p,m,t,p).
\end{array}
\]
All have gap word \((z,q,z,z)\); coincidences of \(t\) with
\(p,m,q\) are retained.  For \(X=p+m-z\), their defects are
\begin{align*}
 T_1:\quad\overline{J_2(F)-J_2(H)}
 &=h_X(h_ph_{m+t-z}-h_mh_{p+t-z}),\\
 T_2:\quad J_2(F)-J_2(H)
 &=h_q\bigl(h_{p+2m-q-z}h_{p+t-z}
            -h_{2p+m-q-z}h_{m+t-z}\bigr),\\
 T_3:\quad\overline{J_2(F)-J_2(H)}
 &=h_X(h_X\bar h_t-h_mh_{p+t-z}).
\end{align*}
The first is nonzero because \(t-z>0,p\ne m\).
The second vanishes only for \(p=m\) or \(t=p+m-q\).
In the third, \(t=q\) leaves a nonzero term; otherwise matching
forces \(t=m\), since \(X>m\).  This is \(T_2\) with \(t=m\),
which would force \(p=q\).  Only the stated family remains.

\emph{Swapped top pair.}
The locus \(x=y\) was already treated.  Otherwise interchange the
profiles so \(x>y\), and suppose \(x'=y,y'=x\).
Modulo \(h_q\), extraction of \(h_x\) gives
\begin{equation}\label{eq:staircase-interior-swap}
 \bar h_eh_{c+d-z}+\overline{h_{d+e-z}}h_c=\bar h_Ah_B.
\end{equation}
Here \(c,c+d-z,B>q\).  Evaluate the surviving variables at one.
Both left terms cannot vanish, since that would give
\(e=q=d+e-z\), contrary to \(d>2z\).  If their monomials
coincide they still contribute twice.  Hence
\[
 A\ne q,\qquad e=q\ \text{or}\ d+e-z=q,
\]
with exactly one alternative true.

Put \(\delta=q-z>0\), and define
\begin{align*}
 N_\delta(c,e)&=h_{c+\delta}h_e+h_ch_{e+\delta}
                         -h_zh_{c+e+\delta-z},\\
 I(c,d,e)&=h_eR_z(c,d)+h_z^2h_{c+q-z}.
\end{align*}
The entire specialization needed here is
\begin{equation}\label{eq:staircase-interior-factor-table}
\begin{array}{c|c}
 \text{condition}&\overline{T_z(c,d,e)}\\\hline
 e=q&-h_zR_z(c,d+\delta)\\
 e\ne q,\ d=q&-h_zN_\delta(c,e)\\
 d+e-z=q&I(c,d,e)\\
 \text{none of these}&T_z(c,d,e).
\end{array}
\end{equation}
The first row includes \(d=e=q\).  The table is exhaustive because
\(c,c+d-z,c+d+e-2z>q\): only \(d,e,d+e-z\) can equal \(q\).
Its primitive factors are irreducible in \(\Q[h_j:j\ge1,j\ne q]\).
Each is linear in its private largest variable, with coefficient
\(-h_z\) for \(R_z,N_\delta\), and \(h_z^2\) for \(I,T_z\).
The remaining polynomial is relatively prime to \(h_z\): modulo
\(h_z\) it is a nonzero product, or for \(N_\delta\) a sum of
two positive monomials.  Gauss's lemma proves irreducibility after
specialization.  No preservation-of-irreducibility assumption is used.
The respective primitive degrees are \(2,2,3,3\), and
\[
 \varepsilon(R_z)=0,\qquad \varepsilon(N_\delta)=1.
\]
These readings also distinguish the quadratic types when smaller
indices coincide.

If \(d+e-z=q\), then \(d,e<q\), and
\eqref{eq:staircase-interior-swap} forces
\[
 A=e,\qquad B=c+d-z=b,\qquad C=c.
\]
Indeed \(B>c\), so \eqref{eq:staircase-interior-data} identifies
it with \(b\).  The strict inequality \(x>y\) gives \(a>e\),
and \(\{D,E\}=\{a,d\}\).
If \(a\ne q\), the mate's tail is untouched by specialization:
its only possible interval contact is
\(a+d-z>e+d-z=q\).  In
\eqref{eq:staircase-interior-specialization}, factor degrees force
\[
 R_z(a,b)=R_z(e,b),
\]
since each side has exactly one irreducible quadratic and one
irreducible cubic.  The factors have the same leading normalization.
It follows that \(a=e\), a contradiction.
If \(a=q\), the source's left factor is \(-h_zh_{q+b-z}\),
and the mate's tail is \(-h_z\) times an irreducible quadratic.
After canceling \(-h_z\), irreducible factor degrees would be
\(1+3\) against \(2+2\), which is impossible.  This separately
excludes the row contact that accompanies the interval contact.

It remains that \(e=q\).  First \(a\ne q\): otherwise the source
specialization has monomial gcd \(h_z^2h_x\), while the mate's
left factor, with \(A\ne q\), is an irreducible quadratic and its
tail has monomial gcd at most \(h_z\).
The source specialization is consequently
\[
 -h_zR_z(a,b)R_z(c,d+\delta).
\]
The mate must have monomial content \(h_z\) and two quadratic
primitives.  Its middle-contact primitive is excluded by
\(\varepsilon(N_\delta)=1\), whereas both source quadratics have
augmentation zero.  Thus \(E=q\).  Canceling \(-h_z\), the
distinct largest indices \(x,y\), swapped on the mate, match the
quadratic factors in
\eqref{eq:staircase-interior-specialization} and give
\[
 \{A,B\}=\{c,d+\delta\},\qquad
 \{C,D+\delta\}=\{a,b\}.
\]
Together with \eqref{eq:staircase-interior-data}, these imply
\[
 \{A,D\}=\{a,d\},\qquad
 \{A,D+\delta\}=\{a,d+\delta\}.
\]
Subtracting square sums gives \(D=d\), then \(A=a\).
If \(B=b\), the top pair was not genuinely swapped.  Otherwise
\(B=c,C=b,b\ne c\), and the factor matching gives \(a=d+\delta\).
The sole remaining pair is
\[
 (d+q-z,b,c,d,q;z,q,z,z),\qquad
 (d+q-z,c,b,d,q;z,q,z,z).
\]
Its full second-coefficient defect is
\begin{equation}\label{eq:staircase-interior-final-defect}
 J_2(F)-J_2(H)=h_q\bigl(
 h_{b+d+q-2z}h_{c+d-z}-h_{b+d-z}h_{c+d+q-2z}\bigr).
\end{equation}
For \(v=b+d-z,w=c+d-z\), its bracket is
\(h_{v+\delta}h_w-h_vh_{w+\delta}\).  Since \(\delta>0\),
monomial matching forces \(v=w\), hence \(b=c\), a contradiction.

For sufficiency in the ordered family, join \((p,m;z)\) and
\((p,t,m;z,z)\) in either order across \(q\).
Both functions equal
\[
 s_{(p,m;z)}s_{(p,t,m;z,z)}
 -h_qs_{(p,t,t,m;z,z,z)}.
\]
Rotating the second joining gives the stated profile.
Every row exceeds \(q\), since \(t=p+m-q>p,m>q\);
every nonsingleton interval exceeds its constituent rows.
Thus this nonrotation family is separated.
\end{proof}

\begin{remark}
The extractions of \(h_x\) and \(h_y\) above follow the specialization
\(h_q=0\); before it, terms carrying \(h_q\) can have larger indices.
The cases \(x=y\), \(a=q\), \(e=q\), \(d=e=q\),
\(d+e-z=q\), and \(t=q\) are retained.  The injectivity bound
comes from \(c>q\) and \(c+d-z>q\), not from a comparison of
\(d\) with \(b,c\) or from any global band hypothesis.
For example, the simultaneous row and interval contact occurs in
\[
 (7,13,9,5,3;1,7,1,1),\qquad
 (3,13,9,5,7;1,7,1,1).
\]
These have equal \(J_0,J_3\), but after setting \(h_7=0\) and
canceling \(-h_1\) their full-function comparison would be
\[
 h_{19}\bigl(h_3h_5h_9-h_1h_3h_{13}+h_1^2h_{15}\bigr)
 =R_1(3,13)R_1(9,11).
\]
The irreducible factor degrees are \(1+3\) and \(2+2\), respectively.
This is the \(a=q,d+e-z=q\) branch, not an excluded degeneration.
\end{remark}

\subsection{Assembly of the five-row classes}

\begin{lemma}[All remaining five-row placements]
\label{lem:staircase-thick-classification}
Every equality of thick five-row profiles is identity, rotation, or a
member of the following rotation quadruples:
\begin{equation}\label{eq:thick-templates}
\begin{array}{c|c|c}
 &F&H\\\hline
\mathrm{A1}&(p,m,m,m,m+p-z;z,z,z,p)&(p,m,m+p-z,m,m;z,p,z,z)\\
\mathrm{A2}&(p,m,m,m+p-z,m+p-z;z,z,p,p)&(p,m,m+p-z,m+p-z,m;z,p,p,z)\\
\mathrm{B}&(p,m,p,p+m-w,m;z,w,z,z)&(p,p+m-w,m,p,m;z,z,w,z).
\end{array}
\end{equation}
For A1 and A2, \(z\ge1,p>z,m>p+z\).  For B,
\(z,w\ge1,p,m>z+w,p\ne m\).
\end{lemma}

\begin{proof}
Classify by the multiplicity of the minimum gap \(z\).
Every safe cut below identifies irreducible full component polynomials,
not merely their specialized residuals.

\emph{A unique minimum.}
If the unique minimum is at an end, its safe cut has factor degrees
\(1,4\).  An interior cut has degrees \(2,3\), so irreducibility
forces the minimum seam of an equal mate to be at an end as well.  
\Cref{lem:staircase-four-plus-one}, with \(q=z\), now gives
rotation-triviality.  

If the unique minimum is interior, normalize
\[
 F=(a,b,c,d,e;r,z,s,t),\qquad r,s,t>z.
\]
The cut gives a two-row and a three-row component, each up to reversal.
A left-only reversal is excluded by
\Cref{lem:staircase-adjacent-swap}.  The other candidates are
\[
 (a,b,e,d,c;r,z,t,s),\qquad (b,a,e,d,c;r,z,t,s).
\]
Their contractions have gap words \((r,s,t)\), \((r,t,s)\).
A nonrotation four-row comparison would be
\eqref{eq:four-family}.  Its unique larger gap would then move between
positions three and four in the original profiles.
\Cref{lem:staircase-seam-crossing} would require the other three gaps,
including \(z\), to coincide, contrary to their strict inequalities.
Thus contractions match directly or by rotation.  Direct matches are
trivial.  A tail-only reversal followed by rotation forces a symmetric
tail.  Reversing both pieces and rotating the contractions gives exactly
\[
 r=s=t,\quad c=a,\quad e=b,\quad d=a+b-z.
\]
These are B, with repeated gap \(r\), exceptional gap \(z\), and
rows \(a,b>r+z\).  They are separated.

\emph{Two minima and distinct other gaps.}
Let the gaps be \(\{z,z,r,q\}\), with \(r,q>z\) distinct.
Neither unique upper gap can change end/interior status by
\Cref{lem:staircase-seam-crossing}.  Together with the degrees in the
safe \(z\)-cut, this fixes the gap word up to rotation.  The four
orders, allowing the names \(r,q\) to exchange, are
\[
\begin{array}{c|c|c}
\text{gaps}&\text{cut sizes}&\text{contracted rows}\\\hline
(z,z,r,q)&(1,1,3)&(a+b+c-2z,d,e)\\
(z,r,z,q)&(1,2,2)&(a+b-z,c+d-z,e)\\
(z,r,q,z)&(1,3,1)&(a+b-z,c,d+e-z)\\
(r,z,z,q)&(2,1,2)&(a,b+c+d-2z,e).
\end{array}
\]
The first two orders have different factor degrees; in the third both
upper gaps are interior, and in the fourth both are at ends.
Distinct gaps identify the two quadratic factors and orient a recovered
three-row factor once the gap word is fixed.
In the first order only the first two singleton rows can exchange, and
\Cref{lem:staircase-adjacent-swap} excludes it.
In the other three orders, the highest coefficient is a three-row
function with the fixed asymmetric gap word \((r,q)\).  Its ordered
rows in the table therefore agree.  Combined with the recovered
components, they fix every original row.

\emph{Two minima and equal other gaps.}
Write \(q>z\).  Up to rotation the four orders are
\(P=zzqq\), \(Q=zqzq\), \(R=zqqz\), \(S=qzzq\).
The safe-cut degrees are \((1,1,3)\) for \(P,R\), and
\((1,2,2)\) for \(Q,S\).  The highest coefficients have the
contracted triples
\[
\begin{array}{c|c}
P&(a+b+c-2z,d,e;q,q)\\
Q&(a+b-z,c+d-z,e;q,q)\\
R&(a+b-z,c,d+e-z;q,q)\\
S&(a,b+c+d-2z,e;q,q).
\end{array}
\]
For \(P/P\), the top coefficient fixes the orientation of the recovered
three-row component: its nontrivial reversal would require
\(a+b=2z\).  Only the first-two exchange remains, and is rigid.
For \(R/R\), rotate the mate to align its middle triple.  An endpoint
exchange is allowed by the top coefficient only when \(a=e\) or
\(b=d\), giving identity or full rotation.
For \(S/S\), align the two recovered quadratic components by rotation.
A single component reversal is an endpoint exchange.  Reversing both
and comparing the top triples gives identity or full rotation.

For \(Q/Q\), the safe \(z\)-cut fixes the singleton \(a\).
If \(a=q\), the \(q\)-restriction is
\[
 -h_zh_{q+b-z}R_z(c,d)h_e.
\]
It recovers \(\{c,d\}\) and \(\{q+b-z,e\}\).  The row multiset
also recovers \(\{b,e\}\), so subtraction of sums of squares fixes
\(b,e\) individually.
If \(a\ne q\), the \(q\)-cut has two irreducible quadratics and the
unique singleton \(h_e\), recovering \(e\).  The highest
\(z\)-coefficient fixes its contracted triple in order: reversal
with the same last row \(e\) forces the first row to equal \(e\)
and changes nothing.  Hence \(b\) is fixed.  Removing the known
\(R_z(a,b)\) from the \(q\)-cut then recovers \(\{c,d\}\).
In both cases only \(c,d\) can exchange.  The safe constant term
\(h_aR_q(b,c)R_q(d,e)\) shows that a nontrivial exchange requires
\(b=e\).  For this remaining comparison, the full difference is
\begin{align*}
 h_zh_q\bigl[&h_{a+b+c-q-z}R_q(b,d)
             -h_{a+b+d-q-z}R_q(b,c)\\
 &+h_{a+b-z}(h_ch_{b+d-q}-h_dh_{b+c-q})\bigr].
\end{align*}
If \(c>d\), its private largest index is \(a+b+c-q-z\); if
\(d>c\), it is \(a+b+d-q-z\).  The coefficients are nonzero,
since \(a>z\) and \(b,c,d>q+z\).  Thus \(c=d\).

For \(Q/S\), the safe cut forces the middle row of \(S\) to equal
\(a\), the singleton of \(Q\).  If \(a\ne q\), the \(q\)-cut
has nonlinear degrees \(2,2\) on \(Q\) and degree three on \(S\),
impossible.  If \(a=q\), it compares an ordinary quadratic primitive
\(B\) with a middle-contact primitive \(M\), of augmentations zero
and minus one.  This is also impossible.
Only \(P/R\) remains, and \Cref{lem:staircase-a2} gives A2 without
any extra contact hypothesis.

\emph{Three minima.}
The unique \(q>z\) is at an end or in the interior.  End/end is
\Cref{lem:staircase-four-plus-one}, and end/interior is
\Cref{lem:staircase-seam-crossing}.  After independently rotating the
profiles, interior/interior has gap word \((z,q,z,z)\) on both sides.
\Cref{lem:staircase-interior} gives identity or the separated
commutation family B, including the balanced top pair and every
contact specialization.

\emph{Four minima.}
All gaps equal \(z\).  Set \(\alpha_i=L_i-z>0\).
If \(r_\alpha\) is the ribbon Schur function, then
\[
 s_D=h_z^5\Phi_z(r_\alpha),\qquad
 \Phi_z(h_j)=h_{j+z}/h_z\quad(j>0).
\]
This substitution is injective in the localization at \(h_z\), and
commutes with reversal.  By
\cite{Billera2006DecomposableCompositions}, a nonrotation ribbon equality
requires reversal choices in at least two nonsymmetric composition
factors.  Each has at least two rows and three cells.  The identity
\(\ell(\alpha\circ\beta)=\ell(\alpha)+|\alpha|(\ell(\beta)-1)\)
forces at five rows an outer factor \((1,2)\) or \((2,1)\) and a
two-row inner factor; row-preserving one-part factors are absorbed into
the inner parts.  Their profiles are B with \(w=z\).  Thickness gives
\(p,m>2z\).

The coverage just proved can be recorded without equation-space tasks:
\begin{equation}\label{eq:staircase-inventory}
\begin{array}{c|c}
\text{minimum-gap multiplicity}&\text{possible nonrotation classes}\\\hline
1&\mathrm B\text{ with a uniquely smaller gap}\\
2,\ \text{other gaps distinct}&\varnothing\\
2,\ \text{other gaps equal}&\mathrm{A2}\\
3&\mathrm{A1}\text{ or }\mathrm B\\
4&\mathrm B\text{ with equal gaps}.
\end{array}
\end{equation}
A1 and A2 are sufficient by the two branch lemmas.  For B put
\(E=(p,m;z)\).  Its two profiles are
\(E\cdot_w(E\odot_wE)\) and
\((E\odot_wE)\cdot_wE\), the latter rotated if necessary.
Both functions equal
\(s_Es_{E\odot_wE}-h_ws_{E\odot_wE\odot_wE}\).
This proves the classification and all its sufficiency assertions.
\end{proof}

\begin{proof}[Completion of the proof of \Cref{thm:staircase-base}]
By \Cref{lem:staircase-closure}, an arbitrary connected mate of a thick
five-row staircase has the same full board pattern.  The minimum-square
cofactor reading in \Cref{lem:readings} recovers its positive gap multiset,
so it is thick as well.  Lemma~\ref{lem:staircase-thick-classification}
therefore leaves exactly \eqref{eq:thick-templates}.
The dictionary to \eqref{eq:staircase-families} is
\[
\begin{array}{c|c}
\mathrm{M1}&(p,m,z,w)=(a+d,b+d,d+t,d)\\
\mathrm{M2}&(p,m,z,w)=(a+d,b+d,d,d+t)\\
\mathrm{M3},\mathrm{M4}&(p,m,z)=(a+d,b+d,d).
\end{array}
\]
For B use M1 when \(z\ge w\), with
\((a,b,t,d)=(p-w,m-w,z-w,w)\), and M2 when \(w>z\), with
\((a,b,t,d)=(p-z,m-z,w-z,z)\).  Both rows and gaps agree, including the
long row \(p+m-w\), and the domain becomes exactly
\(a,b>t+d\).  For A1 and A2 take \((a,b,d)=(p-z,m-z,z)\).
The four members are distinct under the stated nonrotation inequalities.
The sorted row multisets have multiplicities \((1,3,1)\), \((1,2,2)\),
and \((2,2,1)\) for A1, A2, and B, respectively, so these families cannot
mix.  In A1 and A2 the three row values recover \(p,m,k=m+p-z\), then
\(z\).  In B the repeated rows recover \(\{p,m\}\), the largest row
\(k\) gives \(w=p+m-k\), and the gap multiset recovers \(z\).
Exchanging \(p,m\) rotates the same quadruple.  Thus no additional mate
can arise from another parameter choice; at \(z=w\) the M1 and M2
descriptions give the same quadruple.

For unit boards, \Cref{lem:staircase-closure,lem:homogenization} put both
images in the thick domain.  All bare template parameters on a
\(\Gamma\)-image are odd, so \(a,b,t\) in the dictionary are even
and \(d\) is odd.  Replace them by \(a/2,b/2,t/2,(d-1)/2\).
Their inequalities give exactly the integer domains in the theorem,
including \(d=0\).  Injectivity returns the full original class.
\end{proof}

\begin{corollary}[Reduced minors for the contact families]
\label{cor:staircase-contact-minors}
Every nonrotation equality between thick five-row profiles with a contact
is represented by one of the following reduced Toeplitz-minor
coincidences, of orders three and two respectively.
\end{corollary}

\begin{proof}
The B profiles in \eqref{eq:thick-templates} are separated: every row
exceeds both gap values, and every interval index is at least each of
its constituent row lengths.  Thus \Cref{lem:staircase-thick-classification} leaves
only A1 and A2.  In those two families write \(q=p\), \(b=m\),
and set \(\sigma=b-z\), \(\delta=q-z\).  Then
\(0<\delta<q<\sigma\).  Rotate both displayed profiles and normalize
the last column beta to zero.  Their common row-beta list is
\[
 X=(q+4\sigma,q+3\sigma,q+2\sigma,q+\sigma,q).
\]
Their decreasing column-beta lists are the decreasing arrangements of
\[
 \{j\sigma:j\in I\}\ \cup\ \{\delta+j\sigma:j\in J\},
\]
with masks
\[
\begin{array}{c|cc|cc}
 &I&J&I'&J'\\\hline
\mathrm{A1}&\{0,3\}&\{0,1,2\}&\{0,1\}&\{0,2,3\}\\
\mathrm{A2}&\{0,2,3\}&\{0,1\}&\{0,1,2\}&\{0,3\}.
\end{array}
\]
The primed masks describe the second member.  The sets \(I,J\) may
intersect: an element in their intersection is a common shift level in
two different residue sequences, not a repeated column.  Complements
below are taken in \(\{0,1,2,3,4\}\).

In the localization at \(h_q\), put
\[
 U(t)=\sum_{j=0}^4h_{q+j\sigma}t^j,\qquad
 V(t)=\sum_{j=0}^4h_{z+j\sigma}t^j,\qquad
 C(t)=\frac{V(t)}{U(t)}=\sum_{j=0}^4c_jt^j\pmod{t^5}.
\]
Let \(T(C)_{ij}=c_{i-j}\) for \(0\le i,j\le4\), with
\(c_j=0\) for \(j<0\), and define \(T(U)\) similarly.
In increasing row-beta order, columns from the two levels are shifted
columns of \(T(U)\) and \(T(V)\).  Multiplication by
\(T(U)^{-1}\), followed by expansion along the coordinate columns,
gives, including the column-ordering signs,
\[
 s_D=
 \begin{cases}
 -h_q^5\det T(C)[I^c,J],&D\text{ in A1},\\
 h_q^5\det T(C)[I^c,J],&D\text{ in A2}.
 \end{cases}
\]
The primed masks give the same formula for the mate.
For \(\iota(j)=4-j\), both rows of the mask table satisfy
\[
 I'=\iota(J^c),\qquad J'=\iota(I^c).
\]
Transposition and simultaneous reversal of the row and column orders
therefore identify the reduced minors.  Explicitly, they are
\[
 \det\begin{pmatrix}
 c_1&c_0&0\\c_2&c_1&c_0\\c_4&c_3&c_2
 \end{pmatrix}
 =
 \det\begin{pmatrix}
 c_2&c_0&0\\c_3&c_1&c_0\\c_4&c_2&c_1
 \end{pmatrix}
\]
for A1, and
\[
 \det\begin{pmatrix}c_1&c_0\\c_4&c_3\end{pmatrix}
 =\det\begin{pmatrix}c_3&c_0\\c_4&c_1\end{pmatrix}
\]
for A2.  Row and column index sets in these minors are in increasing
order.  Since \(U(0)=h_q\), the required inverse exists in the
localization, and the localization map is injective.  No specialization
at \(h_q=0\) is used in this reduction.

The Toeplitz-minor correspondence and the row/column reversal are
classical; see \cite{GarciaGarcia2020ToeplitzMinors}.
Here the preceding classification supplies the additional conclusion
that equality forces one of these charts.  No claim that an arbitrary
contact profile admits a chart is needed.
\end{proof}

\subsection{Endpoint extensions and borders}\label{sec:staircase-borders}

The minimum-gap classification in \Cref{thm:staircase-base} now supplies
the complete list of possible nonliteral five-row children.  The
coefficients here test their endpoint extensions and borders
directly, including every tie of the largest candidate indices.

\begin{lemma}\label{lem:staircase-endpoint}
For every connected diagram \(C\) with \(n\ge5\) rows and full negative
set \(\mathcal T_n\), common betas allowed, and every integer \(t>0\),
\[
 s_{C^{\uparrow t}}=s_{C^{\downarrow t}}\quad\Longleftrightarrow\quad C=C^*.
\]
\end{lemma}

\begin{proof}
For five rows, unequal endpoint row lengths are separated by the row
multiset.  If they both equal \(a\) but the exterior gaps differ, the
largest proper interval span is
\(\kappa=N-a+\max(z_1,z_4)\).
After cancellation of the full-interval term, its shifted holder has
label \(\kappa+t\) and nonzero coefficient
\(-h_a h_{z_1}h_{z_2}h_{z_3}\) if \(z_4>z_1\), or
\(h_a h_{z_2}h_{z_3}h_{z_4}\) if \(z_1>z_4\).

For tied exterior data write
\(C=(a,x,y,z,a;b,c,d,b)\), \(q=a-b>0\).
The exact domain is
\(b,c,d\ge0,x>b+c,y>c+d,z>b+d\).
Put \(X=x+y-c\), \(Y=y+z-d\), and
\(K_s=h_a h_s-h_bh_{s+q}\).  Define
\[
\begin{split}
 F_{u,v}={}&h_{a+t}(h_{u+q}h_v-h_uh_{v+q})
 -h_a(h_{u+q+t}h_v-h_uh_{v+q+t})\\
 &+h_b(h_{u+q+t}h_{v+q}-h_{u+q}h_{v+q+t}).
\end{split}
\]
Expansion of the sixteen terms gives the full identity
\[
 D_C(t)=h_b(h_yF_{x,z}-h_cF_{X,z}-h_dF_{x,Y}).
\]
Its largest possible label is \(L=\max(X,Y)+q+t\), with coefficient
\[
 [h_L]D_C(t)=\begin{cases}
 h_bh_cK_z,&X>Y,\\
 -h_bh_dK_x,&Y>X,\\
 h_b(h_cK_z-h_dK_x),&X=Y.
 \end{cases}
\]
Each strict coefficient is nonzero, since \(s>b\) makes \(K_s\ne0\).
At the tie, nonsymmetry forces \(x\ne z\); the generator
\(h_{\max(x,z)+q}\) then occurs in just one summand and is above the
other coefficient labels.  Thus the coefficient is again nonzero.
If \(x=z,c=d\), the profile is symmetric.

For a zero-free source with \(n\ge6\), the maximal cofactor is connected
and has pattern \(\mathcal T_{n-1}\).
Put \(u=x_1-x_2,v=y_{n-1}\), \(\kappa=N-\min(u,v)\), and call the
normalized child \(K\).  Formula \eqref{eq:levels}, with its cofactor
bounds, gives
\[
 [h_{\kappa+t}]D_C(t)=\begin{cases}
 (-1)^n s_{K^{\downarrow v}},&u>v,\\
 (-1)^{n+1}s_{K^{\uparrow u}},&v>u,\\
 (-1)^{n+1}D_K(u),&u=v.
 \end{cases}
\]
The strict coefficients are nonzero.  At the tie induction leaves only
\(K=K^*\), which together with \(u=v\) means \(C=C^*\).
Finally homogenize, replacing the extension length by \(2t\), to cover
all common betas.  This induction concerns endpoint differences only;
it does not infer literal cofactor equality from equal polynomials.
\end{proof}

\begin{lemma}\label{lem:staircase-border-separation}
For every pair in \eqref{eq:staircase-families}, no positive connected
border over the \(F\) rotation orbit is Schur-equal to a positive
connected border over the \(H\) rotation orbit.
\end{lemma}

\begin{proof}
We give the endpoint and tied-border coefficients explicitly.  They follow
by substitution in \eqref{eq:interval-expansion}, combining equal labels
before extraction.  Put \(A=h_{a+d}\), \(B=h_{b+d}\), and use
\[
\begin{array}{c|c|c|c|c}
 &c&\ell&\kappa&\sigma\\\hline
\mathrm{M1}&d+t&a+b+d-t&2a+2b+d-2t&t\\
\mathrm{M2}&d&a+b+d&2a+2b+d-t&0\\
\mathrm{M3},\mathrm{M4}&d&a+b+d&a+3b+d&0.
\end{array}
\]
Set \(C=h_c\), \(K=C h_\ell-AB\ne0\); its negative is a connected
two-row Schur polynomial.
For \(\Delta_\epsilon(z)=s_{F^{\epsilon z}}-s_{H^{\epsilon z}}\), the
largest-label coefficients are
\[
\begin{array}{c|c|c}
 &[h_{\kappa+z}]\Delta_\uparrow(z)&[h_{\kappa+z}]\Delta_\downarrow(z)\\\hline
\mathrm{M1},\mathrm{M2}&C^2K&-C^2K\\
\mathrm{M3}&-h_d^2K&\text{not needed}\\
\mathrm{M4}&-Ah_dK&\text{not needed}.
\end{array}
\]
All other labels are smaller under the stated family inequalities.
For M1 and M2 the opposite extension directions are separated by row
multisets, because \(a\ne b\).  For M3 and M4 the distinct endpoint
lengths \(a+d,b+d,a+b+d\) separate every comparison except top against
top.  Thus all positive endpoint cross-comparisons are unequal.
By \Cref{lem:border}, this excludes borders with unequal margins.

At equal margins \((q,q)\), orient the children as displayed.
Their corner and square-marker coefficients cancel, and the linear
coefficient at \(h_{S+q}\) is
\(E_q=\Delta_\uparrow(q)+\Delta_\downarrow(q)\).
For M1 and M2 set \(q_0=\min(a,b)-\sigma>0\),
\(A_q=h_{a+d+q}\), \(B_q=h_{b+d+q}\).
Its largest label is \(L=\max(\kappa,\kappa+q-q_0)\), and
\begin{equation}\label{eq:staircase-tied-coefficient}
\begin{array}{c|c}
\text{regime}&[h_L]E_q\\\hline
q<q_0&C^2(A_qB-AB_q)\\
q>q_0,\ a>b&-BCK\\
q>q_0,\ a<b&ACK\\
q=q_0,\ a>b&AC(B^2-C h_{2b+d-\sigma})\\
q=q_0,\ a<b&BC(C h_{2a+d-\sigma}-A^2).
\end{array}
\end{equation}
For verification of the maximum, every remaining label in the interval
expansion is strictly dominated by one of
\(\kappa,2a+b+d-t+q,a+2b+d-t+q\).
Ordering these gives the three regimes.  At \(q=q_0\) both contributions
are retained, giving the last two rows rather than either strict formula.
The first binomial is nonzero for \(a\ne b,q>0\); the strict branches
contain \(K\ne0\); the tied branches have a generator strictly above
\(b+d>c\), or \(a+d>c\).  Every coefficient is nonzero.

For M3 and M4 the parent row multiset forces \(q=b\): after common
rows are canceled it gives
\(\{2b,a+b+q\}=\{a+2b,b+q\}\).
At that value the largest label of \(E_b\) is \(4b+d\), with coefficient
\(-Ah_dK\) for M3 and \(-A^2K\) for M4.  The inequality \(b>a+d\)
makes these isolated and nonzero.  All labels of \(E_q\) are below the
outer marker \(S+q\).
The formulas allow \(d=0\) with \(h_0=1\); alternatively \(\Gamma\)
sends \((a,b,t,d)\) to \((2a,2b,2t,2d+1)\) and doubles the border
margins.  Thus every unit case descends as well.
\end{proof}

\begin{proposition}\label{prop:staircase-higher-rows}
Every connected source with at least six rows and full negative set
\(\mathcal T_n\) is rotation-rigid against arbitrary connected mates.
Common beta entries are allowed.
\end{proposition}

\begin{proof}
For a zero-free six-row
source the maximal cofactor is a connected five-row staircase.
An arbitrary mate has a positive-row connected cofactor by
\Cref{lem:readings} and irreducibility.
\Cref{lem:staircase-closure,thm:staircase-base} identify its complete
class.  Cross-orbit borders are excluded by
\Cref{lem:staircase-border-separation}; after same-orbit alignment,
\Cref{lem:staircase-endpoint} excludes a nontrivial margin swap.
At larger row counts the cofactor is the connected staircase one order
down.  Induction and the same endpoint assertion apply without an upper
row bound.  Homogenization supplies common betas.  The five-row input is
the algebraic classification in \Cref{lem:staircase-thick-classification};
all subsequent steps use the displayed endpoint and border coefficients.
\end{proof}

\subsection{Verification}\label{sec:certificates}

Both finite-order classifications and their higher-row deductions use
algebraic arguments in the text.  No execution record, feasibility
answer, or exclusion certificate is a premise.

The supplement at \url{https://github.com/johnnyshields/papers} contains
auxiliary checks by signed permutations and column-subset dynamic
programming.  In formal-label checks, numerical substitution is followed
by collection: equal numerical indices do not remain independently
labeled generators.  The cancellation in \Cref{ex:mixed-stratum} is
permitted by the proof, not assigned to an additional search branch.
Finite sample boxes test identities, units, ties, and implementations;
they do not establish coverage of the parameter domains.

\section{Disconnected cofactors and rigidity}\label{sec:component-reductions}

\subsection{Endpoint coefficients}

The notation \(D_Q(t)=s_{Q^{\uparrow t}}-s_{Q^{\downarrow t}}\) continues
from \Cref{sec:boundary-tools}.  By \Cref{thm:two-negatives}, every
one-negative component and its endpoint extensions are rotation-rigid.
Thus
\begin{equation}\label{eq:endpoint-equality}
 D_Q(t)=0\quad\Longleftrightarrow\quad Q=Q^*,\qquad t>0.
\end{equation}
The same assertion holds for dense components by \Cref{lem:small-rows}.

\begin{lemma}\label{lem:one-negative-endpoint}
Let \(Q\ne Q^*\) be connected and zero-free, either dense of order
\(m\ge2\), or one-negative of order \(m\ge3\).  Then the nonzero,
shift-independent leading-coefficient assertion of \eqref{eq:endpoint-top}
holds for every \(t>0\).
\end{lemma}

\begin{proof}
The dense case and orders three and four are \Cref{lem:endpoint-top}.
For a one-negative component of order at least five, use \eqref{eq:levels}.
Every cofactor index is at most \(\kappa\), so an untied top has one
nonzero minor as coefficient.  At \(C_2=S-D_{m-1}\), put
\(d=D_{m-1}\).  The coefficient is \((-1)^{m+1}D_R(d)\), where
\[
 R=((C_2,\ldots,C_m);(D_1,\ldots,D_{m-1})).
\]
This child is connected and one-negative.  By
\eqref{eq:endpoint-equality}, vanishing would make \(R\) symmetric.
Together with the equal exterior gaps \(d\), that would make \(Q\)
symmetric.  All minors are fixed as \(t\) varies.
\end{proof}

\begin{lemma}\label{lem:two-negative-endpoint}
Let \(Q=(C,D)\) be connected and zero-free with four rows and exactly
two negative indices, normalized by \(C_1=S,D_4=0\).  For every \(t>0\),
\eqref{eq:endpoint-top} holds with
\(\kappa=\max(C_2,S-D_3)\) and a nonzero coefficient independent of \(t\).
\end{lemma}

\begin{proof}
In \eqref{eq:levels}, all cofactor indices are at most \(\kappa\).
The coefficient of \(h_{\kappa+t}\) is \(M_{13}\), \(-M_{24}\), or
\(M_{13}-M_{24}\), according as \(S-D_3\) is greater than, less than,
or equal to \(C_2\).  Each strict case is a minor with positive sorted
diagonal and is nonzero.

At a tie put \(d=D_3=S-C_2>0\) and
\(R=((C_2,C_3,C_4);(D_1,D_2,D_3))\).  It has three positive rows and
exactly two negatives.  A connected three-row board has at most one
negative; three singleton components would cost three.  Thus \(R\)
consists of a singleton of length \(s>0\) and a connected dense two-row
component \(T\).  Directly, \(M_{13}-M_{24}=-D_R(d)\).
For the two possible component orders,
\[
 D_R(d)=h_{s+d}s_T-h_s s_{T^{\downarrow d}},\qquad
 D_R(d)=h_s s_{T^{\uparrow d}}-h_{s+d}s_T,
\]
respectively.  Endpoint extension preserves the split.  Modulo \(h_s\),
either expression is, up to sign, \(h_{s+d}s_T\ne0\): the two coordinate
variables are distinct and the connected degree-two factor has no
coordinate-variable divisor.  This proves the tied case.
\end{proof}

\begin{lemma}\label{lem:small-negative-endpoint}
If \(Q\) is connected and zero-free, has at least five rows and at most four negatives, then \(Q\) and its positive endpoint extensions are rotation-rigid against arbitrary connected mates.  For nonsymmetric \(Q\), \eqref{eq:endpoint-top} also holds.
\end{lemma}

\begin{proof}
The rigidity assertions follow from \Cref{thm:two-negatives,thm:column,prop:three-corners,thm:square,prop:hook}; their only four-row exceptions are \eqref{eq:four-family}.  Endpoint extensions preserve negative positions.

Use \eqref{eq:levels}.  Only a tied top level needs consideration.  Its coefficient is an endpoint difference of the maximal cofactor \(R\), with the common exterior gap as extension parameter.  If \(R\) is connected, it has at most four negatives.  The only possible nonrotation ambiguity is the four-row three-negative family
\[
 (w,b,b,b+w-c;c,c,w),\qquad
 (w,b,b+w-c,b;c,w,c),\qquad c<w.
\]
Opposite endpoint extensions retain the same overlap word.  The displayed words, including their reversals, are distinct, so these extensions cannot pass between the two exceptional rotation orbits.  Equality therefore implies symmetry of \(R\), hence of the tied parent \(Q\).

If \(R\) is disconnected, \Cref{lem:component-budget} leaves orders \(1+3\) or \(2+2\) when \(R\) has four rows, and only \(1+4\) when it has five rows.  In the unequal-order cases, factorization separates the extended first component from the unchanged last component by ordinary degree; opposite endpoint extensions cannot agree.  In the \(2+2\) case, equality forces the two factors to exchange.  Two-row rigidity then identifies the components as reverse partners, making \(R\), and thus \(Q\), symmetric.  At six or more cofactor rows a split costs at least five negatives.  This exhausts the tied cases without using the threshold theorem.
\end{proof}

\subsection{A saturated negative column}

\begin{corollary}\label{cor:singleton-two-negatives}
Let \(Q\) be connected and zero-free of order \(m\ge2\), with at most
two negative indices in its full row board.  For all positive integers
\(s,u,v,g\), every parent in \eqref{eq:singleton-context} is rotation-rigid
against arbitrary connected mates.
\end{corollary}

\begin{proof}
The component and its positive endpoint extensions are rotation-rigid by
\Cref{thm:two-negatives}; endpoint extension preserves every negative
position.  For nonsymmetric \(Q\), \eqref{eq:endpoint-top} follows from
\Cref{lem:one-negative-endpoint}
in the dense and one-negative cases.  With two negatives, its order is
at least four: use \Cref{lem:two-negative-endpoint} at order four and
\Cref{lem:small-negative-endpoint} at every larger order.
Thus every hypothesis of \Cref{lem:singleton-lift} holds.
\end{proof}

\begin{theorem}\label{thm:saturated-column}
Let \(F\) be connected with \(n\ge6\) nonempty rows.  Suppose that,
after a possible rotation, the first column of its full row board has
all \(n-2\) negative entries permitted by connectedness, and at most two
negative entries lie outside that column.  Then every connected skew
diagram \(H\) with \(s_H=s_F\) is \(F\) or \(F^*\), up to
translation.  Common beta entries are allowed.
\end{theorem}

\begin{proof}
The hypothesis is preserved by \Cref{lem:homogenization}, so first work
zero-free.  Use the indicated orientation, and write the full beta lists
as \(X,Y\), normalized by \(y_n=0\).  The column condition is
\(x_3<y_1\).  The maximal cofactor has a positive first singleton row,
and its remaining block
\[
 Q=(x_3,\ldots,x_n)/(y_2,\ldots,y_{n-1})
\]
has \(m=n-2\ge4\) positive rows and at most two negatives.  Its diagonal
entries are subdiagonal entries of the original board, hence positive.
A disconnection after row \(k\) would require a negative lower-left
rectangle with \(k(m-k)\ge m-1\ge3\) entries.  Thus \(Q\) is connected
and zero-free.

Put
\[
 s=x_2-y_1,\qquad u=x_1-x_2,\qquad
 v=y_{n-1},\qquad g=y_1-x_3.
\]
These parameters are positive.  Normalize \(Q\) by subtracting \(v\)
from both lists.  Its span is \(S=x_3-v\), so the offset in
\eqref{eq:singleton-context} is \(z=v+S+g=y_1\).  Substitution recovers
exactly the lists of \(F\).  Apply \Cref{cor:singleton-two-negatives}.
Finally, homogenization transports an arbitrary full-board equality at
its common row count, preserves the hypothesis, and descends the rotation
conclusion by injectivity and rotation covariance.
\end{proof}

The negative count is \(n-2\), \(n-1\), or \(n\), and is unbounded.
For every \(n\ge6\), the partition \((3,1^{n-3})\) is realizable by
\Cref{thm:negative-realizability}; it satisfies the column hypothesis and
has \(n\) negatives.  The lower row bound cannot be dropped: the first
source at \((b,g,q)=(2,3,1)\) in \eqref{eq:five-families} has three
first-column negatives and two outside, but has a nonrotation mate.

\subsection{Two nonlinear components}

\begin{lemma}\label{lem:two-component-lift}
Suppose the zero-free maximal cofactor has two connected components \(Q,R\), each with at least two rows.  Assume that each component and its endpoint extensions are rotation-rigid against arbitrary connected mates.  Assume also that every nonsymmetric component satisfies \eqref{eq:endpoint-top}, with a nonzero coefficient independent of the extension parameter.  Then the parent is rotation-rigid against arbitrary connected mates.
\end{lemma}

\begin{proof}
Normalize \(Q=(C,D)\), \(C_1=A,D_p=0\), and \(R=(E,F)\), \(E_1=S,F_r=0\).  Factorization recovers their functions and hence their literal diagrams up to order and rotation.  Every parent of the ordered components has lists
\begin{equation}\label{eq:two-component-lists}
 \begin{aligned}
 z&=v+S+g,\\
 X&=(z+A+u,z+C_1,\ldots,z+C_p,v+E_1,\ldots,v+E_r),\\
 Y&=(z+D_1,\ldots,z+D_p,v+F_1,\ldots,v+F_r,0),
 \end{aligned}
 \qquad u,v,g>0.
\end{equation}
Its span is \(N=A+S+u+v+g\).  At \(N-\min(u,v)\), the coefficient after corner removal, up to one common sign, is
\begin{equation}\label{eq:two-component-marker}
 \begin{cases}
 s_{Q^{\uparrow u}}s_R,&u<v,\\
 s_Qs_{R^{\downarrow v}},&v<u,\\
 s_{Q^{\uparrow u}}s_R+s_Qs_{R^{\downarrow u}},&u=v.
 \end{cases}
\end{equation}
Divisibility by the old irreducible factors identifies the regimes.  If \(s_Q=s_R\), cancel that factor: a strict marker has one endpoint-extension contribution, whereas a tied marker has two.  At their common largest index, the tied coefficient is exactly twice the strict one, since both extensions have the same signed maximal cofactor.  At a tie the span recovers \(g\).

If the two old factors differ, rotate to align their order and subtract two proposed tied markers.  Reduction modulo \(s_Q\) shows that \(s_Q\) divides the difference of the two proposed top extensions of \(Q\).  Such a difference, if nonzero, is homogeneous of ordinary degree \(p\) and weighted degree \(|Q|+u\).  Divisibility by \(s_Q\), of ordinary degree \(p\) and weighted degree \(|Q|\), would force a scalar quotient and hence \(u=0\), a contradiction.  Thus the top extensions agree.  The same argument modulo \(s_R\) equates the bottom extensions.  Endpoint rigidity fixes each component orientation.  For equal old factors, cancel their common factor.  Of the four orientation choices, a whole-parent rotation pair gives \(s_{Q^{\uparrow u}}+s_{Q^{\downarrow u}}\); the other two give \(2s_{Q^{\uparrow u}}\) and \(2s_{Q^{\downarrow u}}\).  Endpoint rigidity separates these three values unless the component is symmetric.

For strict margins rotate to have \(u<v\).  The marker identifies the ordered \(Q\), its orientation, and \(u\).  Suppose two remaining candidates have \(v_-<v_+\).  Subtract the common initial last-column terms.  At \(K=\max(N-v_-,S+v_+)\), the coefficient of their difference is
\begin{equation}\label{eq:two-component-separator}
 \mathbf1_{K=N-v_-}(-1)^{p+r}s_Qs_{R_-^{\downarrow v_-}}
 +\mathbf1_{K=S+v_+}(-1)^{r+1}s_T.
\end{equation}
The second cofactor has subdiagonal indices
\[
 C_1-D_1,\ldots,C_p-D_p,E_2^+-F_1^+,\ldots,E_r^+-F_{r-1}^+>0.
\]
Thus \(T\) is connected, and its span \(N-v_+\) is below \(K\).  All indices of the product coefficient are also below \(K\).  Strict maxima are nonzero, and at a tie irreducibility separates \(s_T\) from the product.  The values \(v,g\) are fixed.

Only reversal of \(R\) remains.  It is literal if \(R=R^*\).  Otherwise, write \(\kappa_R,c_R\) for its offset and nonzero coefficient in \eqref{eq:endpoint-top}.  For \(0\le i\le p\), let \(A_i\) be the determinant with column list \(D\) and row list \((A+u,C_1,\ldots,C_p)\) with entry \(i\) deleted, counting from zero.  Then \(A_0=s_Q\), \(A_1=s_{Q^{\uparrow u}}\).  Put
\[
 w_0=g+A+u,\qquad w_i=g+C_i\quad(i\ge1).
\]
Laplace expansion expresses the orientation difference as a signed sum of \(A_i\mathcal E_R(w_i,v)\).  For \(v<w_0\) the \(i=0\) term has a private largest index.  For \(v>w_0\), the coefficient at \(S+v\) has the private moving index \(\kappa_R+w_0\).  At \(v=w_0\) the \(i=0\) term vanishes, leaving, up to sign,
\begin{equation}\label{eq:two-component-tie}
 T=\sum_{i=1}^p(-1)^{p-i}A_i D_R(w_i).
\end{equation}
Write
\[
 L=A+u,\qquad t_0=\kappa_R+g+A,\qquad
 B_Q=\det(h_{C_i-D_j})_{2\le i\le p,\ 1\le j<p}.
\]
The sorted diagonal of \(B_Q\) is positive, so \(B_Q\ne0\), and
\([h_L]A_1=(-1)^{p+1}B_Q\).
For \(i\ge2\), the top index of \( D_R(w_i)\) is below \(t_0\).  Also \(t_0>A\), whereas every index in \([h_L]A_i\) is below \(A\).  Extracting the largest indicated coefficients in \eqref{eq:two-component-tie} gives
\begin{equation}\label{eq:two-component-tie-table}
 \begin{array}{c|c}
 t_0>L &[h_{t_0}]T=(-1)^{p-1}A_1c_R\ne0\\
 t_0=L &[h_L^2]T=B_Qc_R\ne0\\
 t_0<L &[h_Lh_{t_0}]T=B_Qc_R\ne0.
 \end{array}
\end{equation}
No row-order restriction was placed on \(Q\) beyond \(p\ge2\).  This excludes the last orientation change.
\end{proof}

\subsection{Minimum realizing row count six}

\begin{theorem}\label{thm:six-negative-minimum}
Let \(F\) be connected with \(\nu(F)=6\) and
\(n_{\min}(\rho(F))=6\).  Every connected skew diagram \(H\) satisfying
\(s_H=s_F\) belongs to \(\{F,F^*\}\).  The two members are distinct,
and common beta entries are allowed.
\end{theorem}

\begin{proof}
First work zero-free with six rows.  The maximal cofactor has five positive
rows and is disconnected by \Cref{lem:above-minimum-border}.
Three components cost at least seven negatives.  The component budget
\eqref{eq:component-budget} therefore leaves only
\[
\begin{array}{c|c|c}
\text{component orders}&\text{cross cost}&\text{internal negatives}\\\hline
1+4&4&2\\
2+3&6&0.
\end{array}
\]
For \(1+4\), the component and its endpoint extensions are rigid by
\Cref{thm:two-negatives}; \Cref{lem:two-negative-endpoint} supplies
\Cref{lem:singleton-lift}'s remaining hypothesis.  For \(2+3\), both
components and their endpoint extensions are dense and rigid by
\Cref{lem:small-rows}.  Their endpoint coefficients satisfy
\Cref{lem:one-negative-endpoint}, so \Cref{lem:two-component-lift} applies.
Both lifting lemmas compare with arbitrary connected mates.

For \(n>6\), the maximal cofactor \(R\) is connected, has one fewer row,
and retains the negative partition.  Induction applies to it and to its
positive endpoint extensions.  An equal mate is zero-free by
\Cref{lem:readings}; factorization of the maximal coefficient identifies
its connected cofactor with \(R\) or \(R^*\).  After rotation aligns the
literal cofactor, \Cref{lem:border} recovers the unordered margins.
A strict swap equates \(R^{\uparrow t}\) and \(R^{\downarrow t}\).
They are not literal equals, so induction makes them rotations; undoing
the common extension gives \(R=R^*\), and the swap rotates the parent.

\Cref{lem:homogenization} descends the zero-free conclusion to full boards
at their common row count.  Finally, the negative partitions in this
stratum are
\[
 (4,2),\ (2,2,1,1),\ (4,1,1),\ (3,1,1,1),\ (3,3),\ (2,2,2).
\]
Each differs from its conjugate, so \Cref{lem:negative-partition} gives
\(F\ne F^*\).  No negative partition was imposed on the mate.
\end{proof}

\section{Dense endpoint differences}\label{sec:dense-endpoints}

The final placement comparison needs more than a nonzero leading coefficient.  It needs the endpoint difference of a dense four-row component not to consist solely of that term.

\begin{lemma}\label{lem:three-variable-endpoint}
If \(T\ne T^*\) is dense with three rows, then
\[
 s_{T^{\uparrow t}}(x,y,z)-s_{T^{\downarrow t}}(x,y,z)\ne0
 \qquad(t>0).
\]
\end{lemma}

\begin{proof}
Write its normalized lists as
\(A_1>A_2>A_3>B_1>B_2>0\), with last column beta zero, and put
\[
 \lambda=(A_1-A_3-2,A_2-A_3-1,0),\qquad
 \mu=(B_1-2,B_1-B_2-1,0).
\]
The three-variable alternant formula gives
\begin{equation}\label{eq:three-variable-factorization}
 s_T(x,y,z)=(xyz)^{A_3-B_1+2}s_\lambda(x,y,z)s_\mu(x,y,z).
\end{equation}
Top extension adds \(t\) to the first part of \(\lambda\), and bottom extension adds \(t\) to the first part of \(\mu\).

For fixed \(t>0\), the ratio
\(s_{(p+t,q,0)}/s_{(p,q,0)}\) determines \((p,q,0)\).  To see this, write \((p,q,0)=(a+b,b,0)\) and \(H_j=h_j(x,y)\).  At \(z=0\), the ratio is \(H_{a+t}/H_a\).  Setting \(x=1\) gives
\[
 \frac{1-y^{a+t+1}}{1-y^{a+1}},
\]
whose first nonconstant term has degree \(a+1\), determining \(a\).  The branching rule gives
\begin{equation}\label{eq:branching-ratio}
 (xy)^{-b}s_{(a+b,b)}(x,y,z)
 =\sum_{i=0}^{a}\sum_{j=0}^{b}(xy)^{-j}H_{a-i+j}z^{i+j}.
\end{equation}
For \(b<b'\), the first difference in the cross-multiplied ratios occurs at \(z^{b+1}\): the extra summand at that degree has \(i=0,j=b+1\).  After multiplying by \((xy)^{b+1}\), its coefficient is
\[
 H_{a+t}H_{a+b+1}-H_aH_{a+t+b+1}
 =(xy)^{a+1}H_{t-1}H_b\ne0.
\]
Thus \(b\) is determined as well.  Equality of the two extended functions in \eqref{eq:three-variable-factorization} forces \(\lambda=\mu\), or \(A_1=A_3+B_1=A_2+B_2\), exactly the beta reflection equations for \(T=T^*\).
\end{proof}

\begin{lemma}\label{lem:dense-four-lower-part}
Let \(Q\ne Q^*\) be dense with normalized four-row lists
\begin{equation}\label{eq:dense-four-lists}
 C_1=S>C_2>C_3>C_4>D_1>D_2>D_3>D_4=0.
\end{equation}
Put \(\kappa=\max(C_2,S-D_3)\).  There is a nonzero polynomial \(c_Q\), independent of \(t>0\), such that
\begin{equation}\label{eq:dense-four-top}
 \max\operatorname{index} D_Q(t)=\kappa+t,
 \qquad [h_{\kappa+t}] D_Q(t)=c_Q,
\end{equation}
and
\begin{equation}\label{eq:dense-four-remainder}
  D_Q(t)\ne h_{\kappa+t}c_Q.
\end{equation}
\end{lemma}

\begin{proof}
In \eqref{eq:levels}, all cofactor indices are at most \(\kappa\).  For unsigned minors \(M_{ij}\) of \(Q\),
\[
 c_Q=\begin{cases}
 -M_{24},&C_2>S-D_3,\\
 M_{13},&C_2<S-D_3,\\
 M_{13}-M_{24},&C_2=S-D_3.
 \end{cases}
\]
In the strict cases this is a nonzero dense three-row function, and remains nonzero in three variables.  At the tie put \(\alpha=S-C_2=D_3\), and let \(T\) be the normalized maximal cofactor of \(Q\).  Then \(c_Q=- D_T(\alpha)\).  The child \(T\) is dense with three rows.  Since \(Q\) is its border with equal margins, nonsymmetry of \(Q\) implies nonsymmetry of \(T\).  \Cref{lem:three-variable-endpoint} shows that \(c_Q(x,y,z)\ne0\).

Every dense four-row determinant vanishes in three variables.  Indeed, over \(\mathbb Q(x,y,z)\),
\[
 h_j(x,y,z)=\sum_{a\in\{x,y,z\}}
 \frac{a^{j+2}}{\prod_{b\in\{x,y,z\}\setminus\{a\}}(a-b)}
 \qquad(j\ge0).
\]
Thus a matrix \((h_{C_i-D_j}(x,y,z))\) with positive indices has rank at most three, by separating the factors \(a^{C_i+2}\) and \(a^{-D_j}\) in each summand.  This applies to both endpoint extensions of \(Q\), whereas
\[
 h_{\kappa+t}(x,y,z)c_Q(x,y,z)\ne0.
\]
Their difference therefore cannot be just its leading term.  This proves \eqref{eq:dense-four-remainder} as well as the nonvanishing in \eqref{eq:dense-four-top}.
\end{proof}

\begin{lemma}\label{lem:dense-cross-resonance}
For a dense four-row \(Q\), write \(f=s_Q\), \(f_x=s_{Q^{\uparrow x}}\), and \(f^{\,y}=s_{Q^{\downarrow y}}\).  If \(x>\ell>0\), \(u>0\), and \(v=x+u\), then
\begin{align}
 h_{\ell+u}f_x-h_{\ell+v}f-h_\ell D_Q(v)&\ne0,
 \label{eq:cross-resonance-one}\\
 h_{\ell+u} D_Q(x)-h_\ell D_Q(v)&\ne\pm h_{\ell+v}f.
 \label{eq:cross-resonance-two}
\end{align}
\end{lemma}

\begin{proof}
For symmetric \(Q\), the first assertion is factorization of two different connected endpoint functions against coordinate factors, and the second is immediate.  Suppose \(Q\ne Q^*\).  Write
\[
 B_Q=\det(h_{C_i-D_j})_{2\le i\le4,\ 1\le j\le3},
 \qquad [h_S]f=-B_Q.
\]
For \eqref{eq:cross-resonance-one}, the possible largest indices are
\(S+x,\ell+v,\kappa+v\).  If \(\ell<\kappa\), a tie of the first and third has coefficient
\(-h_{\ell+u}B_Q-h_\ell c_Q\), nonzero modulo \(h_\ell\).  If \(\ell>\kappa\), a tie of the first and second has coefficient \(-h_SB_Q-f\), nonzero modulo \(h_S\).  If \(\ell=\kappa\), the remaining coefficients are \(-f-h_\ell c_Q\) or
\[
 -h_SB_Q-f-h_\ell c_Q.
\]
The first is nonzero modulo \(h_\ell\); the second contains the diagonal monomial of \(f\), all of whose indices are strictly below \(\kappa\), so that monomial cannot cancel.  A strict maximum has a nonzero coefficient.

The left side of \eqref{eq:cross-resonance-two} has largest index \(\kappa+v\), with coefficient \(-h_\ell c_Q\).  If \(\ell<\kappa\), a matching largest index on the right is possible only at \(S=\kappa+v\), where its coefficient is not divisible by \(h_\ell\).  If \(\ell>\kappa\), the right side has a larger index.  If \(\ell=\kappa\), the right side has either a larger index, a coefficient \(f\) not divisible by \(h_\ell\), or a nonzero square of its top generator.  None can agree with the left side.  Only \eqref{eq:dense-four-top}, not the lower remainder, has been used.
\end{proof}

\section{A dense four-row component and two singletons}\label{sec:three-components}

Fix \(Q\) as in \eqref{eq:dense-four-lists}, and retain \(f,f_t,f^{\,t}, D_Q\) from \Cref{sec:dense-endpoints}.  Write \(\Phi(x,y)=\Phi_Q(x,y)\).

\begin{lemma}\label{lem:three-component-forms}
Let a connected zero-free diagram have maximal-cofactor function \(h_\ell h_m f\), where \(\ell,m>0\) and \(f=s_Q\) is dense of order four.  Every connected Schur-equal diagram, after a possible rotation, has one of the following two forms, allowing \(Q\) or \(Q^*\).  Put
\[
 L=\ell+g_1,\quad R=m+g_2,\quad
 \ell,m,g_1,g_2,u,v>0,\quad N=S+L+R+u+v.
\]
The end placement \(E\), with component order \((\ell),(m),Q\), has lists
\begin{equation}\label{eq:end-lists}
 \begin{aligned}
 X_E&=(N,N-u,S+R+v,v+C_1,v+C_2,v+C_3,v+C_4),\\
 Y_E&=(N-u-\ell,S+R+v-m,v+D_1,v+D_2,v+D_3,v,0).
 \end{aligned}
\end{equation}
The middle placement \(M\), with order \((\ell),Q,(m)\), has lists
\begin{equation}\label{eq:middle-lists}
 \begin{aligned}
 X_M&=(N,N-u,v+R+C_1,v+R+C_2,v+R+C_3,v+R+C_4,v+m),\\
 Y_M&=(N-u-\ell,v+R+D_1,v+R+D_2,v+R+D_3,v+R,v,0).
 \end{aligned}
\end{equation}
Both are connected and zero-free.  All possible component orders are represented, including both singleton orders.
\end{lemma}

\begin{proof}
The row-count reading gives seven rows.  Factorization of \([h_N]s_D=h_\ell h_m f\) recovers exactly two one-row factors and one connected four-row factor.  Nonnegative-index rigidity in \cite{Shields2026Cofactor} identifies the latter as \(Q\) or \(Q^*\), against arbitrary mates.  Up to rotation, it lies at an end or between the singletons.  Reading the exterior row and column gaps gives \(u,v>0\); the two component separations give \(g_1,g_2>0\).  These are exactly \eqref{eq:end-lists} and \eqref{eq:middle-lists}.  Conversely their subdiagonal entries are the positive rows of the indicated components, and their separating entries are negative.  The strict lists and positive diagonal give the asserted diagrams.  No hypothesis about a mate's negative partition has been made.
\end{proof}

For these forms, set
\begin{equation}\label{eq:auxiliary-parent-functions}
 \begin{aligned}
 A_L&=h_{\ell+u}h_{m+L}-h_\ell h_{m+L+u},\\
 U_Q(x)&=h_{\ell+u}f_x-h_\ell f_{x+u},\\
 Z_Q(y)&=h_{m+v}f^{\,y}-h_m f^{\,y+v}.
 \end{aligned}
\end{equation}
First-column and last-row expansion give the full parent polynomials
\begin{align}
 P_E={}&A_L\Phi(R,v)-h_mh_{\ell+u}\Phi(L+R,v)
             +h_\ell h_m\Phi(L+R+u,v),\label{eq:end-polynomial}\\
 P_M={}&h_{\ell+u}h_{m+v}\Phi(L,R)-h_\ell h_{m+v}\Phi(L+u,R)
 \notag\\
 &-h_{\ell+u}h_m\Phi(L,R+v)+h_\ell h_m\Phi(L+u,R+v).
 \label{eq:middle-polynomial}
\end{align}
Here \(A_L\) is a dense two-row Schur function.  For \(x>\ell\), the function \(U_Q(x)\) is the connected five-row determinant with lists
\begin{equation}\label{eq:intermediate-U}
 X=(S+x+u,S+x,C_2,C_3,C_4),\qquad
 Y=(S+x-\ell,D_1,D_2,D_3,0).
\end{equation}
Its negative partition is \((1,1,1)\).  The corresponding \(Z_Q(y)\) has partition \((3)\), and is the rotation of the \(U\)-diagram built on \(Q^*\) with singleton data \(m,v\).  \Cref{thm:column} applies to these intermediates.  Two \(U\)-diagrams with equal functions are literally equal, since rotation changes the negative partition; equality of a \(U\)- and a \(Z\)-diagram is rotation.  Their lists then recover their parameters.

\subsection{End and middle placements do not mix}

\begin{lemma}\label{lem:end-middle-separation}
No end placement in \eqref{eq:end-lists} is Schur-equal to a middle placement in \eqref{eq:middle-lists}.
\end{lemma}

\begin{proof}
Let \(q=\min(u,v)\).  The coefficient at the next index \(N-q\), after corner removal, is
\begin{equation}\label{eq:first-marker-table}
 \begin{array}{c|c|c}
 &E&M\\ \hline
 u<v&-h_mh_{\ell+u}f&-h_mh_{\ell+u}f\\
 v<u&-h_\ell h_m f^{\,v}&-h_\ell h_{m+v}f\\
 u=v&-h_m(h_{\ell+u}f+h_\ell f^{\,u})
     &-f(h_mh_{\ell+u}+h_\ell h_{m+u}).
 \end{array}
\end{equation}
Every coefficient index is below \(N-q\).  Divisibility by \(f\) separates the last two end regimes from all middle regimes.  Divisibility by \(h_\ell h_m\) distinguishes those two end regimes from one another.  After division by \(f\), setting every positive-index generator equal to one distinguishes a strict middle marker from its tied marker.

Thus a cross-placement equality requires \(u<v\) on both sides, after rotating the middle parent.  The original singleton multiset \(\{\ell,m\}\), together with \(\{\ell+u,m\}\), determines which singleton was extended: the change in the sum of squares is \(2\ell u+u^2\).  Therefore \(u,\ell,m,N,S\) agree.

Subtract
\[
 J=f(h_\ell h_mh_N-h_mh_{\ell+u}h_{N-u}).
\]
At \(K=\max(N-v,N-u-L)\), the next coefficients are
\begin{equation}\label{eq:third-marker-table}
 \begin{array}{c|c|c}
 &N-v>N-u-L&N-v<N-u-L\\ \hline
 E&E_1=-h_\ell h_m f^{\,v}&E_2=A_Lf\\
 M&M_1=-h_\ell h_{m+v}f&M_2=h_mU_Q(L).
 \end{array}
\end{equation}
At equality of the indices use \(E_0=E_1+E_2\), \(M_0=M_1+M_2\).  These coefficients are nonzero.  The factors \(A_L\) and \(U_Q(L)\) are irreducible of ordinary degrees two and five.  In particular \(E_0\) is divisible by neither \(f\) nor \(h_m\).  The nine comparisons reduce as follows:
\begin{equation}\label{eq:cross-marker-exclusions}
 \begin{array}{c|ccc}
 &M_1&M_2&M_0\\ \hline
 E_1&f\text{-divisibility}&\text{irreducibility of }U&\text{first tie below}\\
 E_2&A_L\text{ is not a monomial}&f\text{-divisibility}&f\text{-divisibility}\\
 E_0&f\text{-divisibility}&h_m\text{-divisibility}&\text{second tie below}.
 \end{array}
\end{equation}
For \(E_1=M_0\), reduction modulo \(h_\ell\) forces \(m=\ell\).  The common \(K\) fixes \(v\), and the middle tie gives \(L=v-u\).  Opposite core orientations would require
\(h_{\ell+u}f_L=h_{\ell+v}f\), contrary to factorization.  Equal orientations would contradict \eqref{eq:cross-resonance-one}.

For \(E_0=M_0\), the tied indices give common \(v,L\), with \(v=L+u\).  The remaining equation is
\[
 h_{\ell+u}(h_{m+L}f-h_mf_L)
       +h_\ell h_m(f_v-f_E^{\,v})=0,
\]
where the subscript \(E\) specifies the end core orientation.  Reduction modulo \(h_m\) forces \(m=\ell+u\).  Opposite orientations again give an impossible product equality; equal orientations give \eqref{eq:cross-resonance-one}, up to sign.  Both ties are excluded.
\end{proof}

\subsection{End placements}

\begin{lemma}\label{lem:end-placement-rigidity}
Two end placements have equal Schur functions only when their aligned diagrams are identical.
\end{lemma}

\begin{proof}
The first marker in \eqref{eq:first-marker-table} distinguishes its three margin regimes.

If \(u=v\), the marker recovers the singleton order: for unequal lengths, exactly \(h_m\) among the two singleton variables divides it.  Subtraction recovers \(f^{\,u}\) and fixes the core orientation.  The span fixes \(L+R\).  The last two terms of \eqref{eq:end-polynomial} are then common, leaving \(A_L\Phi(R,v)\).  Its connected factors of degrees two and five recover \(L,R\).

Suppose \(v<u\).  The first marker fixes \(v\) and the core orientation, but leaves the singleton order.  Put \(T=L+R+u=N-S-v\), and subtract \(h_\ell h_m\Phi(T,v)\).  The negative of the coefficient at \(K=\max(N-u,m+L+u)\) is
\begin{equation}\label{eq:end-bottom-marker}
 \begin{cases}
 h_mh_{\ell+u}f,&N-u>m+L+u,\\
 h_\ell\Phi(R,v),&N-u<m+L+u,\\
 h_mh_{\ell+u}f+h_\ell\Phi(R,v),&N-u=m+L+u.
 \end{cases}
\end{equation}
The first type is divisible by \(f\), the second has singleton content \(h_\ell\), and the tied type has neither singleton divisor when \(\ell\ne m\).  If \(\ell=m\), a proposed equality of the second and tied types fixes \(g_2=T-K\), hence the same \(\Phi(R,v)\), and leaves a nonzero extra term.

Within the first type, \(u=N-K\) and the shifted singleton pair fix the order; the degree-two/degree-five factorization then fixes \(L,R\).  Within the second, the marker fixes \(\ell,m,R\) and \(L+u\).  Canceling the common term involving \(h_\ell h_{m+L+u}\) leaves equality of
\[
 h_{\ell+u}\bigl(h_{m+L}\Phi(R,v)-h_m\Phi(R+L,v)\bigr).
\]
The bracket is connected of order six, so factorization fixes \(u,L\).  At the tie \(u=N-K\), \(g_2=T-K\).  If the singleton order changed, with \(a_0=S-D_1\), the row multiset would give
\[
 \{\ell+u,a_0+m+g_2\}=\{m+u,a_0+\ell+g_2\}.
\]
For \(\ell\ne m\), this forces \(u=a_0+g_2\), whereas the marker tie gives \(u=S+v+g_2>a_0+g_2\).  The order is fixed.

Suppose now \(u<v\).  Use the end row of \eqref{eq:third-marker-table}.  The types \(E_1,E_2,E_0\) are respectively divisible by only \(h_\ell\), only \(f\), or neither, among those two polynomials.  For \(E_1\), the marker fixes \(v\) and the core orientation, after which \(A_L\Phi(R,v)\) fixes \(L,R\).  For \(E_0\), the tie fixes \(v=N-K\), \(L=v-u\), and then \(R\); its coefficient fixes orientation.

For \(E_2\), the marker fixes \(L=N-u-K\), while \(R+v\) is known.  If \(v_-<v_+\), subtract the common first three last-column terms.  At \(K'=\max(N-v_-,S+v_+)\), the difference has coefficient
\begin{equation}\label{eq:end-moving-bottom}
 \mathbf1_{K'=N-v_-}h_\ell h_m f_-^{\,v_-}
 -\mathbf1_{K'=S+v_+}s_T.
\end{equation}
Deleting the first core row and last column from the larger-margin parent gives \(T\).  Its subdiagonal indices are
\(\ell,m,C_2-D_1,C_3-D_2,C_4-D_3>0\), so it is connected of order six.  Its span \(N-v_+\) and all product-coefficient indices are below \(K'\).  An irreducible degree-six function cannot equal a product with two singleton factors.  Thus \(v,R\) are fixed.

It remains to compare \(Q\) and \(Q^*\) at fixed margins in this last case.  Write \(\mathcal E(w,v)=\Phi_Q(w,v)-\Phi_Q(v,w)\).  Its leading coefficient is \(- D_Q(w)\) at \(S+v\) for \(w<v\), or \( D_Q(v)\) at \(S+w\) for \(w>v\).  The end-parent difference is
\begin{equation}\label{eq:end-orientation}
 A_L\mathcal E(R,v)-h_mh_{\ell+u}\mathcal E(L+R,v)
                       +h_\ell h_m\mathcal E(L+R+u,v).
\end{equation}
Put \(w_0=L+R+u\).  If \(v<w_0\), the last term has a private largest index.  If \(v>w_0\), the coefficient at \(S+v\) has the private index \(\kappa+w_0\).  At \(v=w_0\), that term vanishes and the top coefficient becomes
\[
 T_E=h_mh_{\ell+u} D_Q(L+R)-A_L D_Q(R).
\]
Compare \(t_0=\kappa+L+R\) and \(Z=m+L+u\).  If \(t_0>Z\), the first term has a private index.  If \(t_0<Z\), its \(h_Z\)-coefficient is \(h_\ell D_Q(R)\ne0\).  At equality, \(u=\kappa+R-m\) and
\[
 [h_Z]T_E=h_mh_{\ell+u}c_Q+h_\ell D_Q(R).
\]
Its possible largest indices are \(\ell+u\) and \(m+u\).  Unequal values give a private contribution; equal values give coefficient \(2h_\ell c_Q\ne0\).  Thus a nonsymmetric core cannot reverse at fixed margins.
\end{proof}

\subsection{Middle placements}

\begin{lemma}\label{lem:middle-placement-rigidity}
Two middle placements have equal Schur functions only when they are identical or whole-parent rotations.
\end{lemma}

\begin{proof}
First fix \(\ell,m,u,v,L+R\), and compare \(L_-<L_+\), \(R_->R_+\), allowing either core orientation.  At
\(K=\max(S+L_++u,S+R_-+v)\), expansion of \eqref{eq:middle-polynomial} gives
\begin{equation}\label{eq:middle-movement}
 [h_K](P_+-P_-)=
 \mathbf1_{K=S+L_++u}h_\ell Z_{Q_+}(R_+)
 -\mathbf1_{K=S+R_-+v}h_mU_{Q_-}(L_-).
\end{equation}
All coefficient indices are below \(K\).  A strict maximum is nonzero.  At a tie, factorization forces \(\ell=m\) and equality of the two connected five-row functions.  Their single-row/single-column rigidity forces rotation, and \eqref{eq:intermediate-U} gives
\[
 u=v,\qquad R_+=L_-,\qquad Q_+=Q_-^*.
\]
Since \(L+R\) was fixed, this is whole-parent rotation.

When \(u=v\), the first marker fixes that margin and \(L+R\).  Align the singleton order by rotating the parent; \eqref{eq:middle-movement} then fixes \(L,R\) up to whole-parent rotation.  Core orientation at fixed parameters is treated below.

For \(u<v\), the middle row of \eqref{eq:third-marker-table} applies.  The type \(M_1\) is divisible by \(f\); the type \(M_2\) is divisible by \(h_m\); \(M_0\) is not divisible by \(h_m\) unless \(\ell=m\).  A proposed equality between a tied type and \(M_2\), with \(\ell=m\), fixes the same \(L\), and the tie gives \(v=L+u\).  Cancellation leaves either an immediate nonzero term or
\[
 h_{\ell+u} D_Q(L)-h_\ell D_Q(L+u)
       =\pm h_{\ell+v}f,
\]
excluded by \eqref{eq:cross-resonance-two}.  Thus the types agree.

Within \(M_1\), the value \(v=N-K\) is fixed, and \eqref{eq:middle-movement} fixes \(L,R\).  Within \(M_2\), the value \(L=N-u-K\) and \(U_Q(L)\) are recovered.  The intermediate rigidity fixes the core orientation.  Since \(R+v\) is known, write
\[
 P_M=h_{m+v}\mathcal A_Q(L,R)-h_m\mathcal A_Q(L,R+v),
 \quad
 \mathcal A_Q(L,y)=h_{\ell+u}\Phi_Q(L,y)-h_\ell\Phi_Q(L+u,y).
\]
The second term is common and \(\mathcal A_Q(L,R)\) is connected of order six.  Factorization of the first term fixes \(v,R\).  Within \(M_0\), the tie fixes \(v=N-K\), \(L=v-u\), and \(R\).  Subtracting its known product term recovers \(U_Q(L)\), hence orientation.

Finally compare the two core orientations at fixed parameters, assuming \(Q\ne Q^*\).  At \(K=\max(S+L+u,S+R+v)\), the coefficient of the difference is
\begin{align}
 [h_K](P_Q-P_{Q^*})={}&
 \mathbf1_{K=S+L+u}h_\ell\bigl(Z_Q(R)-Z_{Q^*}(R)\bigr)\notag\\
 &+\mathbf1_{K=S+R+v}h_m\bigl(U_Q(L)-U_{Q^*}(L)\bigr).
 \label{eq:middle-orientation}
\end{align}
A strict maximum is nonzero by intermediate rigidity.  At \(L+u=R+v\), the coefficient simplifies to
\begin{equation}\label{eq:middle-final-tie}
 T_M=h_mh_{\ell+u} D_Q(L)-h_\ell h_{m+v} D_Q(R).
\end{equation}
If \(L=R\), then \(u=v\), and
\[
 T_M=(h_mh_{\ell+u}-h_\ell h_{m+u}) D_Q(L).
\]
It vanishes only for \(\ell=m\), when reversing the core is whole-parent rotation.

Assume \(L>R\), so \(v>u\), and put
\(p=\kappa+L\), \(a=\ell+u\), \(b=m+v\).
Every strict maximum among \(p,a,b\) in \eqref{eq:middle-final-tie} is private.  If \(a=b\), remove the common generator and use the private top of \( D_Q(L)\).  If \(a=p>b\), the first product has a nonzero square coefficient; this also handles the triple tie.  The only remaining comparison is \(b=p>a\), for which
\[
 [h_p]T_M=h_mh_{\ell+u}c_Q-h_\ell D_Q(R).
\]
The two tie equations imply \(u=\kappa+R-m\).  The possible largest indices in this coefficient are \(\ell+u\) and \(\kappa+R=m+u\).  If \(\ell\ne m\), one is private.  If \(\ell=m\), vanishing would require
\[
  D_Q(R)=h_{\kappa+R}c_Q,
\]
contrary to \Cref{lem:dense-four-lower-part}.  This excludes the final nonrotation possibility.
\end{proof}

\begin{theorem}\label{thm:dense-three-component-rigidity}
A connected zero-free diagram whose maximal cofactor consists of a dense four-row component and two singletons is rotation-rigid against arbitrary connected mates.
\end{theorem}

\begin{proof}
\Cref{lem:three-component-forms} exposes every mate in an end or middle form.  \Cref{lem:end-middle-separation} compares the two placements, and \Cref{lem:end-placement-rigidity,lem:middle-placement-rigidity} compare each placement internally, including both core orientations.  The conclusion is therefore not restricted to mates sharing the source's negative partition.
\end{proof}

\section{The sharp rigidity theorem}\label{sec:threshold-proof}

\begin{proof}[Proof of \Cref{thm:negative-threshold}]
An empty diagram contributes \(1\), whereas every nonempty connected component, including a single row or column, has a nonconstant irreducible Schur function \cite{Barekat2009CoincidencesCorrigendum}.  Since \(s_D\) is nonconstant and irreducible, \(s_H=s_D\) forces \(H\) to have exactly one nonempty component.  Thus \(H\) is connected.

First work with zero-free diagrams at their minimum realizing row count.  Put \(m=n-1\).  By \Cref{lem:above-minimum-border}, the maximal cofactor is disconnected.

At \(n=7\), its component orders sum to six.  Four or more components cost at least twelve negatives; among three-component orders only \(1+1+4\) costs at most nine.  The complete budget is
\begin{equation}\label{eq:threshold-budget-table}
 \begin{array}{c|c|c|l}
 \text{orders}&\text{cross cost}&\text{internal budget}&\text{rigidity input}\\ \hline
 1+5&5&\le4&\text{five-row core: \Cref{lem:small-negative-endpoint,lem:singleton-lift}}\\
 2+4&8&\le1&\text{dense two-row and \(\le1\)-negative four-row cores}\\
 3+3&9&0&\text{two dense three-row cores}\\
 1+1+4&9&0&\Cref{thm:dense-three-component-rigidity}.
 \end{array}
\end{equation}
The second and third rows satisfy every hypothesis of \Cref{lem:two-component-lift}: component and endpoint rigidity follow from \Cref{thm:two-negatives,lem:small-rows}, and their endpoint-leading coefficients follow from \Cref{lem:one-negative-endpoint}.  In particular the \(3+3\) case uses only the endpoint-leading property, not \Cref{lem:dense-four-lower-part}.  The first row uses arbitrary-mate rigidity of the five-row core and its endpoint extensions as well as its endpoint-leading property, not rigidity alone.

For \(n\ge8\), three components cost at least \(2m-3\ge11\).  Two components both of size at least two cost at least \(2(m-2)\ge10\).  Thus only \(1+(n-2)\) remains.  Its cross cost is \(n-2\), and its internal budget is at most \(11-n\).  For \(n=8,9,10,11\), this gives respectively cores of orders \(6,7,8,9\) with at most \(3,2,1,0\) internal negatives.  \Cref{lem:small-negative-endpoint,lem:one-negative-endpoint,lem:singleton-lift} apply.  For \(n\ge12\), even the least cross cost exceeds nine.  This covers every minimum-order case without a partition census.

Now induct on the actual row count above the minimum.  By \Cref{lem:above-minimum-border}, the maximal cofactor is connected, retains \(\rho(D)\), and has one fewer row.  It and its endpoint extensions therefore satisfy the inductive rigidity assertion.  Factorization forces the arbitrary mate's maximal cofactor to be connected and identifies it up to rotation.  After alignment, \Cref{lem:border} recovers the margins.  A surviving swap equates opposite endpoint extensions.  They cannot be translates; if they are rotations, undoing the extension makes the cofactor symmetric and the swap rotates the parent.  This completes the zero-free induction.

For full boards, use the common row count and \eqref{eq:gamma} on both source and mate.  Their negative partitions, and hence the source's \(n_{\min}\), are unchanged.  The zero-free conclusion descends by injectivity and rotation covariance.  Finally \Cref{prop:ten-negative-quartet} supplies the sharpness assertion.
\end{proof}

\begin{corollary}\label{cor:ten-negative-high-minimum}
A connected diagram with exactly ten negative indices and minimum
realizing row count at least eight is rotation-rigid against arbitrary
connected mates, with common betas allowed.
\end{corollary}

\begin{proof}
At minimum row count eight, a disconnected seven-row cofactor has either
orders \(1+6\), with at most four internal negatives, or \(2+5\), both
dense.  Three components cost at least eleven negatives.
Apply \Cref{lem:small-negative-endpoint,lem:singleton-lift,lem:two-component-lift}.
At minimum row counts nine through twelve only \(1+(n-2)\) remains,
with at most \(12-n\) internal negatives, so the same singleton argument
applies.  Above twelve even the least cross cost exceeds ten.
Connected bordering and homogenization finish exactly as in the proof of
\Cref{thm:negative-threshold}.
\end{proof}

\begin{proof}[Completion of \Cref{thm:classification}(v)]
The partitions of six, up to conjugation, are
\[
 (6),\quad (5,1),\quad (4,2),\quad (4,1,1),\quad (3,3),\quad (3,2,1).
\]
Their minimum realizing row counts are respectively \(8,7,6,6,6,5\),
by \Cref{thm:negative-realizability}.  The first two pairs are rigid by
\Cref{thm:negative-threshold}; the next three by
\Cref{thm:six-negative-minimum}.  The staircase is classified by
\Cref{thm:staircase-base,prop:staircase-higher-rows}.
As in the proof of parts (i)--(iii), a Schur-equal skew diagram is connected
by irreducibility.  Thus all comparisons include arbitrary skew mates.
\end{proof}

\section{Sharpness and failures beyond the hypotheses}\label{sec:threshold-boundaries}

\subsection{Seven-row ten-negative families}

\begin{proposition}\label{prop:ten-negative-quartet}
Let \(a,b,d\) be integers with \(a\ge1\), \(d\ge0\), and \(b>a+d\).
The profiles \eqref{eq:ten-family-a-profiles} and \eqref{eq:ten-profiles}
define two families of four distinct connected Schur-equal diagrams
\(F_i,F_i^*,H_i,H_i^*\), \(i=A,B\), of respective areas
\(3a+11b+7d\) and \(4a+11b+7d\).  In either family their negative
partitions are, respectively,
\begin{equation}\label{eq:ten-patterns}
 (3,3,2,1,1),\quad(5,3,2),\quad(4,2,2,1,1),\quad(5,3,1,1).
\end{equation}
All have ten negatives and minimum realizing row count seven.  Both
families have proper legal marked-composition presentations.
\end{proposition}

\begin{proof}
Define family A by
\begin{equation}\label{eq:ten-family-a-profiles}
 \begin{aligned}
 F_A={}&(a+d,b+d,a+b+d,2b+d,2b+d,a+2b+d,3b+d;\\
      &\hspace{12mm}d,a+d,b+d,b+d,a+b+d,2b+d),\\
 H_A={}&(a+d,b+d,a+b+d,2b+d,a+2b+d,3b+d,2b+d;\\
      &\hspace{12mm}d,a+d,b+d,a+b+d,2b+d,b+d),
 \end{aligned}
\end{equation}
and family B by
\begin{equation}\label{eq:ten-profiles}
 \begin{aligned}
 F_B={}&(a+d,b+d,a+b+d,2b+d,a+2b+d,a+2b+d,3b+d;\\
      &\hspace{12mm}d,a+d,b+d,a+b+d,a+b+d,2b+d),\\
 H_B={}&(a+d,b+d,a+b+d,2b+d,a+2b+d,3b+d,a+2b+d;\\
      &\hspace{12mm}d,a+d,b+d,a+b+d,2b+d,a+b+d).
 \end{aligned}
\end{equation}
The profile inequalities \eqref{eq:profile-domain} follow from the
parameter bounds.  Both have common column list
\(Y=(6b,5b,4b,3b,2b,b,0)\).  Their row lists are
\begin{equation}\label{eq:ten-family-a-beta-lists}
 \begin{aligned}
 X_{F_A}&=(a+6b,6b,a+5b,5b,4b,a+3b,3b)+d,\\
 X_{H_A}&=(a+6b,6b,a+5b,5b,a+4b,4b,2b)+d,
 \end{aligned}
\end{equation}
and
\begin{equation}\label{eq:ten-beta-lists}
 \begin{aligned}
 X_{F_B}&=(a+6b,6b,a+5b,5b,a+4b,a+3b,3b)+d,\\
 X_{H_B}&=(a+6b,6b,a+5b,5b,a+4b,4b,a+2b)+d.
 \end{aligned}
\end{equation}
Here adding \(d\) means adding it to every entry.  Since
\(0\le d<a+d<b\), counting columns strictly above each row beta gives
\eqref{eq:ten-patterns}; rotation gives the conjugate partitions.
\Cref{thm:negative-realizability} gives minimum seven, and summing the
rows gives the areas.  In each family the last-row lengths differ and
both exceed the common first-row length, while the displayed profiles
are not identical.  Thus the four members are distinct.  They are
zero-free exactly when \(d>0\).

For the marked compositions of McNamara and van Willigenburg
\cite{McNamara2009Classification}, take
\begin{equation}\label{eq:ten-chart}
 \begin{gathered}
 D_A=(3,4,4;2,3),\qquad D_B=(2,3,3,3;1,2,2),\\
 E=(a+d,b+d;d),\qquad W=(a+d),
 \end{gathered}
\end{equation}
in attachment case (c).  Put \(L=a+d\), \(M=b+d\),
\(\beta=b-1\), and use integer cell coordinates
\[
 \begin{aligned}
 E&=\{(0,j):\beta\le j<\beta+L\}\cup\{(1,j):0\le j<M\},\\
 W_{\rm sw}&=\{(1,j):0\le j<L\},\qquad
 W_{\rm ne}=\{(0,j):\beta\le j<\beta+L\}.
 \end{aligned}
\]
The inequality \(\beta\ge L\) makes each occurrence maximal on its
content diagonals.  Their clearance is \(b-a-d+1\ge2\).
Deleting the northeast occurrence leaves a row; deleting the southwest
occurrence leaves two intervals overlapping in \(d+1\) cells.
The middle part is the single row \(\{(1,j):L\le j<M\}\), whose
hatted part is empty.  These verify Hypotheses I--IV, respectively.
The southwest attachment is horizontal and the northeast attachment
vertical.  Both outer diagrams are connected and proper.

Put \(\tau=(-1,b-1)\), \(\eta=(1,1)\), and
\[
 \mathcal U=\bigcup_{j=0}^{5}(E+j\tau)
             \ \cup\ \bigcup_{j=1}^{4}(E+\eta+j\tau).
\]
For family A the northwest-ribbon rule gives
\begin{equation}\label{eq:ten-family-a-unions}
 \begin{aligned}
 D_A\circ_W E&=\mathcal U\cup(E+2\eta+3\tau)=H_A,\\
 D_A^*\circ_W E&=\mathcal U\cup(E+2\eta+2\tau)=F_A.
 \end{aligned}
\end{equation}
Its supplemental interfaces are already covered by these copies.
The normalized left endpoints of the seven rows are
\((6-r)(b-1)\), \(r=0,\ldots,6\); the last copy extends rows
\(r=4,5\), or \(r=5,6\), to lengths \(a+2b+d\) and \(3b+d\).
For family B the supplemental interface gives
\begin{equation}\label{eq:ten-cell-unions}
 \begin{aligned}
 D_B\circ_W E&=\mathcal U\cup(E+2\eta+2\tau)
                       \cup(W_{\rm ne}+2\eta+3\tau)=F_B,\\
 D_B^*\circ_W E&=\mathcal U\cup(E+2\eta+3\tau)
                       \cup(W_{\rm ne}+2\eta+\tau)=H_B.
 \end{aligned}
\end{equation}
Reading the intervals gives the displayed profiles.  Outer reversal
and whole-parent rotation give all four constructed members.

Their equalities also have a direct determinant proof.  Localize at
\(h_d\), and put
\[
 U(t)=\sum_{j=0}^{6}h_{jb+d}t^j,\qquad
 V(t)=\sum_{j=0}^{6}h_{a+jb+d}t^j,\qquad
 \frac{V(t)}{U(t)}\equiv\sum_{j=0}^{6}c_jt^j\pmod{t^7}.
\]
Toeplitz elimination expresses each original determinant as \(h_d^7\)
times a minor of \(T(c)_{ij}=c_{i-j}\), where \(c_k=0\) for \(k<0\).
For family A the two index pairs are
\begin{equation}\label{eq:ten-family-a-minors}
 (\{3,5,6\},\{0,1,2\}),\qquad
 (\{4,5,6\},\{0,1,3\}).
\end{equation}
For family B the equality is
\begin{equation}\label{eq:ten-minor-equality}
 \det\begin{pmatrix}
 c_3&c_2&c_1&0\\c_4&c_3&c_2&c_0\\
 c_5&c_4&c_3&c_1\\c_6&c_5&c_4&c_2
 \end{pmatrix}
 =
 \det\begin{pmatrix}
 c_2&c_1&c_0&0\\c_4&c_3&c_2&c_1\\
 c_5&c_4&c_3&c_2\\c_6&c_5&c_4&c_3
 \end{pmatrix},
\end{equation}
with index pairs \((\{3,4,5,6\},\{0,1,2,4\})\) and
\((\{2,4,5,6\},\{0,1,2,3\})\).
In each case reflection \(i\mapsto6-i\), followed by transposition,
identifies the two minors.  This is the classical Toeplitz-minor reversal
of \cite{GarciaGarcia2020ToeplitzMinors}.
The determinant ring embeds in the localization, so the equalities hold
in the original ring.  At \(d=0\), the localized element is already
\(h_0=1\).
\end{proof}

\subsection{Six-row quartets}

\begin{proposition}\label{prop:six-row-quartets}
For every \(1\le q\le b<g\), let
\begin{equation}\label{eq:six-row-list}
 X=(b+5g,b+4g,b+3g,b+2g,b+g,b)
\end{equation}
and take the following column lists:
\begin{equation}\label{eq:six-column-lists}
 \begin{array}{c|l|l}
 \nu&Y(F_\nu)&Y(H_\nu)\\ \hline
 7&(3g+q,2g,g+q,g,q,0)&(3g,2g+q,g+q,g,q,0)\\
 8&(4g+q,2g+q,g+q,g,q,0)&(3g+q,2g+q,2g,g+q,q,0)\\
 9&(4g+q,3g+q,g+q,g,q,0)&(3g+q,3g,2g+q,g+q,q,0).
 \end{array}
\end{equation}
For each row of the table, \(F_\nu,F_\nu^*,H_\nu,H_\nu^*\) are four distinct connected Schur-equal diagrams.  Their minimum realizing row count is six.  Their source partitions are
\begin{equation}\label{eq:six-patterns}
 \begin{array}{c|c|c}
 \nu&\rho(F_\nu)&\rho(H_\nu)\\ \hline
 7&(4,2,1)&(4,2,1)\\
 8&(4,2,1,1)&(4,3,1)\\
 9&(4,2,2,1)&(4,3,2).
 \end{array}
\end{equation}
No upper bound on these connected classes is asserted.
\end{proposition}

\begin{proof}
Put \(s=b-q\), \(A=s+g\), \(B=s+2g\).  The profiles are
\begin{equation}\label{eq:six-profiles}
 \begin{array}{c|l}
 F_7&(B,B+q,B,A+q,A,b;A,A+q,A,b,s)\\
 H_7&(B+q,B,B,A+q,A,b;A+q,A,A,b,s)\\
 F_8&(A,B,B,A+q,A,b;s,A,A,b,s)\\
 H_8&(B,B,A+q,A,A,b;A,A,b,s,s)\\
 F_9&(A,A,B,A+q,A,b;s,s,A,b,s)\\
 H_9&(B,A+q,A,A,A,b;A,b,s,s,s).
 \end{array}
\end{equation}
They satisfy \eqref{eq:profile-domain}, including \(s=0\).  Comparing \eqref{eq:six-row-list} and \eqref{eq:six-column-lists} gives \eqref{eq:six-patterns}; \Cref{thm:negative-realizability} gives minimum six.  The first row distinguishes each \(F_\nu\) from \(H_\nu\).  Every displayed last row has length \(b\), and every displayed first row is longer, so neither is a rotation of either displayed source.  This gives fourfold distinctness independently of the determinant equality.

Reverse the six matrix rows and index them by \(k=0,\ldots,5\).  Define
\[
 U_j=(h_{b+(k-j)g})_{k=0}^5,\qquad
 V_j=(h_{b+(k-j)g-q})_{k=0}^5.
\]
For the down-shift \(L\) on six-vectors, \(LU_j=U_{j+1}\) and \(LV_j=V_{j+1}\), because \(b<g\).  Use the shift identity from the proof of \eqref{eq:five-families}:
\begin{equation}\label{eq:shift-identity}
 \sum_{i=1}^{6}\det(z_1,\ldots,Lz_i,\ldots,z_6)=0.
\end{equation}
For the seed word \((V_2,U_2,V_1,U_1,V_0,U_0)\), four summands have repeated columns, and the other two give
\[
 \det(V_3,U_2,V_1,U_1,V_0,U_0)
 =\det(U_3,V_2,V_1,U_1,V_0,U_0).
\]
For the seed word \((V_3,V_2,V_1,U_1,V_0,U_0)\), the same four vanishing positions give
\[
 \det(V_4,V_2,V_1,U_1,V_0,U_0)
 =\det(V_3,V_2,U_2,V_1,V_0,U_0).
\]
These are the seven- and eight-negative equalities.

The nine-negative equality uses a different elimination.  Set
\(u_k=h_{b+kg}\), \(v_k=h_{b+kg-q}\), and let \(T_v=(v_{i-j})_{0\le i,j\le5}\), with \(v_j=0\) for \(j<0\).  In the localization at \(v_0=h_{b-q}\), define
\[
 \sum_{k=0}^{5}c_kz^k\equiv
 \frac{\sum_{k=0}^{5}u_kz^k}{\sum_{k=0}^{5}v_kz^k}\pmod{z^6},
 \qquad C_j=\sum_{k=j}^{5}c_{k-j}e_k.
\]
Then \(T_v^{-1}V_j=e_j\), \(T_v^{-1}U_j=C_j\).  The two transformed determinants are
\begin{equation}\label{eq:nine-minor}
 \det(e_4,e_3,e_1,C_1,e_0,C_0)
 =c_2c_4-c_1c_5
 =\det(e_3,C_3,e_2,e_1,e_0,C_0).
\end{equation}
Restoring \(T_v\) proves equality.  The original polynomials embed in the localization, so the identity holds before localization as well.  In particular \(b=q\) is included, with \(v_0=h_0=1\).  The reduced two-by-two Toeplitz-minor interpretation belongs to the correspondence of \cite{GarciaGarcia2020ToeplitzMinors}; it is not one further application of \eqref{eq:shift-identity}.
\end{proof}

\subsection{A non-dense placement resonance}

\begin{proposition}\label{prop:nondense-resonance}
Assume \(1\le q<b<g\), set \(s=b-q\), and let
\[
 Q=(b,s+g,s+g;s,s).
\]
This connected three-row component is one-negative and rotation-rigid, but satisfies
\begin{equation}\label{eq:nondense-resonance}
 h_b s_{Q^{\uparrow(g-q)}}-h_{s+g}s_Q
 =h_s\bigl(s_{Q^{\uparrow g}}-s_{Q^{\downarrow g}}\bigr).
\end{equation}
At \(t=g\), its endpoint difference has a nonzero lower remainder below its moving leading term.  Its maximal cofactor is \(h_s^2\), a product of two singletons.  This resonance permits the distinct middle and end placements in the eight-negative quartet of \Cref{prop:six-row-quartets}.
\end{proposition}

\begin{proof}
The matrix of \(Q\) is
\[
 \begin{pmatrix}
 h_b&h_{b+g}&h_{b+2g}\\
 h_s&h_{s+g}&h_{s+2g}\\
 0&h_s&h_{s+g}
 \end{pmatrix}.
\]
It has exactly one negative index, so \Cref{thm:two-negatives} gives arbitrary-mate rigidity.  Direct expansion gives
\begin{align*}
 f=s_Q={}&h_bh_{s+g}^2-h_bh_sh_{s+2g}
           -h_{b+g}h_sh_{s+g}+h_{b+2g}h_s^2,\\
 s_{Q^{\uparrow(g-q)}}={}&h_{s+g}^3-2h_sh_{s+g}h_{s+2g}
                           +h_s^2h_{s+3g},\\
  D_Q(g)={}&h_{b+g}h_{s+g}^2-h_{b+2g}h_sh_{s+g}
                -h_bh_{s+g}h_{s+2g}+h_bh_sh_{s+3g}.
\end{align*}
These prove \eqref{eq:nondense-resonance}.  Here \(\kappa=s+2g\), and
\[
 [h_{\kappa+g}] D_Q(g)=h_bh_s\ne0.
\]
After removing this term, the largest remaining index is \(b+2g\), with coefficient \(-h_sh_{s+g}\ne0\).  Thus even nondivisibility by the moving top generator survives the resonance.  On the other hand,
\[
 [h_{b+2g}]f=h_s^2,
\]
so ordinary degree two must not be substituted for irreducible nonlinearity of the maximal cofactor.

In \eqref{eq:six-column-lists}, the maximal cofactors of \(F_8,H_8\) have component orders \((s,K,s)\) and \((K,s,s)\), where
\[
 K=((b+2g,b+g,b);(g+q,q,0)),\qquad K^*=Q.
\]
Both corner coefficients are \(-h_s^2s_K\).  The remaining tied comparison is exactly \eqref{eq:nondense-resonance}, with \(u=q\), \(x=g-q>s\), and \(v=g=x+u\).  All these parameters are strictly positive.  This is an interior end-to-middle equality with the same rigid nonlinear component, not an ambiguity in recovering its function.
\end{proof}

For seven negatives the four-row maximal-cofactor components themselves can be nonrotation mates: their profiles are
\[
 (A,A+q,A,b;s,b,s),\qquad(A+q,A,A,b;b,s,s),
\]
with \(s=b-q>0\), \(A=s+g\), and negative partition \((2,1)\).  They belong to the four-row exceptional family \eqref{eq:four-family}.  For nine negatives the maximal cofactors instead have four components, with orders \((s,s,R,s)\) and \((R,s,s,s)\), where \(R=(A,b;s)\).  Thus the nine-negative construction is not another three-component resonance.

\begin{remark}\label{rem:ribbon-staircase}
The staircase \((4,3,2,1)\) has ten negatives but minimum realizing row count six.  A classical ribbon example is
\[
 \begin{aligned}
 F&=(a,b,a,a+b,a+b,b;0,0,0,0,0),\\
 H&=(a,a+b,a+b,b,a,b;0,0,0,0,0),
 \end{aligned}
 \qquad a,b>0,\quad a\ne b.
\]
For \(\gamma=(a,b)\), these are \(\gamma\cdot\gamma^{\odot3}\) and \(\gamma^{\odot3}\cdot\gamma\).  The ribbon product identity, in the factor-reversal theory of Billera, Thomas and van Willigenburg \cite{Billera2006DecomposableCompositions}, gives
\[
 s_F=r_\gamma r_{\gamma^{\odot3}}-r_{\gamma^{\odot4}}=s_H.
\]
Their first and last row lengths also distinguish the rotation orbits.  This ribbon construction is different from the seven-row sharpness families.  It supplies no generation theorem for all sources with the staircase pattern, in particular for nonribbon sources with that pattern.
\end{remark}

\section{Ten negatives at minimum row count seven}\label{sec:ten-negative}

\subsection{Five-negative cofactor components}

\begin{lemma}\label{lem:five-negative-leading}
Let \(Q=(C,D)\ne Q^*\) be connected and zero-free with five rows and
five negative indices, normalized by \(C_1=S,D_5=0\).  There are an
integer \(0<k_Q<S\) and a nonzero polynomial \(c_Q\), independent of
\(t>0\), such that
\begin{equation}\label{eq:five-negative-leading}
 \max\operatorname{index}D_Q(t)=k_Q+t,\qquad
 [h_{k_Q+t}]D_Q(t)=c_Q.
\end{equation}
Every index in \(c_Q\) is at most \(k_Q\).
\end{lemma}

\begin{proof}
Put \(\kappa=\max(C_2,S-D_4)\).  In \eqref{eq:levels}, all fixed
cofactor indices are at most \(\kappa\).  An untied top has one
nonzero minor as coefficient.  At a tie put \(q=S-C_2=D_4\).
The coefficient is, up to sign, \(D_R(q)\), where \(R\) is the
maximal cofactor of \(Q\).

The four positive rows of \(R\) contain five negatives.  By
\Cref{lem:negative-partition,lem:component-budget}, its components are
two singletons and a dense two-row diagram.  If the nonlinear component
is at an end, factorization separates its extended function from its
unextended function, so \(D_R(q)\ne0\).  If it is in the middle,
\[
 D_R(q)=(h_{\ell+q}h_m-h_\ell h_{m+q})s_T,
\]
which vanishes exactly when \(\ell=m\).  In every nonvanishing case,
\eqref{eq:five-negative-leading} holds with \(k_Q=\kappa\).

In the remaining case the profile has the form
\[
 Q=(b+q,x,y,z,b+q;b,c,d,b),\qquad
 x=b+c+u,\quad z=b+d+w,\quad y=c+d-\sigma,
\]
where \(b,q,c,d,u,w,\sigma>0\) and \(\sigma<\min(c,d)\).
The two negative separating entries give \(u,w>0\); the dense two-row
middle component gives the inequalities on \(\sigma\).
Put \(X=x+y-c\), \(Y=y+z-d\), and
\[
 \mathcal A_{r,s}(t)=-\det\begin{pmatrix}
 h_b&h_{b+q}&h_{b+q+t}\\
 h_r&h_{r+q}&h_{r+q+t}\\
 h_s&h_{s+q}&h_{s+q+t}
 \end{pmatrix}.
\]
Expansion along the first row and last column gives
\begin{equation}\label{eq:five-negative-alternants}
 D_Q(t)=h_b\bigl(h_y\mathcal A_{x,z}(t)
 -h_c\mathcal A_{X,z}(t)-h_d\mathcal A_{x,Y}(t)
 +h_\sigma\mathcal A_{X,Y}(t)\bigr).
\end{equation}
Let \(K_j=h_{b+q}h_j-h_bh_{j+q}\).  The coefficient at
\(h_{\max(X,Y)+q+t}\) is
\[
 h_b\begin{cases}
 h_cK_z-h_\sigma K_Y,&X>Y,\\
 -h_dK_x+h_\sigma K_X,&Y>X,\\
 h_cK_z-h_dK_x,&X=Y.
 \end{cases}
\]
For \(X>Y\), its \(h_{Y+q}\)-coefficient is
\(h_b^2h_\sigma\ne0\); the reversed case is analogous.
At \(X=Y\), nonsymmetry forces \(x\ne z\): otherwise \(c=d\)
and \(u=w\).  The larger of \(h_{x+q},h_{z+q}\) then has a nonzero
coefficient.  Every fixed index in \eqref{eq:five-negative-alternants}
is at most \(k_Q=\max(X,Y)+q\), and this positive offset is below
\(S\).  The same coefficient works for every \(t>0\).
\end{proof}

The offset in \Cref{lem:five-negative-leading} need not be the ordinary
boundary value \(\max(C_2,S-D_4)\).

\begin{lemma}\label{lem:singleton-five-rigid}
Let \(Q\) be a connected zero-free five-row five-negative diagram whose
complete connected class is \(\{Q,Q^*\}\).  Every connected zero-free
parent with maximal cofactor consisting of a singleton and \(Q\) is
rotation-rigid.
\end{lemma}

\begin{proof}
Factorization recovers the singleton and \(Q\) up to rotation.  The
marker \eqref{eq:singleton-marker} and the separator
\eqref{eq:margin-separator} recover the margins as in the proof of
\Cref{lem:singleton-lift}.  Once the core has been recovered, fixing
its orientation when \(v\le u\) requires only
\(D_Q(t)=0\) exactly when \(Q=Q^*\), supplied by
\Cref{lem:five-endpoint}; arbitrary-mate rigidity of the endpoint
extensions is unnecessary.

For \(u<v\), after the margins are fixed, only core reversal remains.
Put \(a=s+g\) and
\(\mathcal E(w,v)=\Phi_Q(w,v)-\Phi_Q(v,w)\).
The difference is
\[
 h_{s+u}\mathcal E(a,v)-h_s\mathcal E(a+u,v).
\]
A zero summand leaves the other nonzero.  For \(v<a+u\) the second
nonzero summand has a private largest index.  For \(v>a+u\), the
coefficient at \(S+v\), up to sign, is
\[
 h_{s+u}D_Q(a)-h_sD_Q(a+u).
\]
Its second term has the private index \(k_Q+a+u\), larger than both
\(k_Q+a\) and \(s+u\), by \Cref{lem:five-negative-leading} and
\(a>s\).  Thus a nonsymmetric core cannot reverse.  For a symmetric
core the reversal is literal.
\end{proof}

\subsection{A one-negative four-row component}

\begin{lemma}\label{lem:one-negative-four-lower}
Let \(Q\ne Q^*\) be connected and zero-free with four rows and one
negative index.  Normalize its lists as
\begin{equation}\label{eq:one-negative-four-lists}
 C_1=S>C_2>C_3>D_1>C_4>D_2>D_3>D_4=0.
\end{equation}
Put \(\kappa=\max(C_2,S-D_3)\), and let \(c_Q\ne0\) be the
fixed leading coefficient supplied by \Cref{lem:one-negative-endpoint}.
Then
\begin{equation}\label{eq:one-negative-four-remainder}
 D_Q(t)\ne h_{\kappa+t}c_Q\qquad(t>0).
\end{equation}
\end{lemma}

\begin{proof}
Every diagonal index is strictly below \(\kappa\).  If the first and
last row lengths differ, the minimum-square monomials of the two endpoint
extensions have different square weights.  Their difference is twice
\(t\) times that row-length difference.  The smaller minimum cannot
occur in the other polynomial, and all its indices are below
\(\kappa+t\).  It supplies a nonzero lower term.

Suppose the end rows agree, so \(S=C_4+D_1\), and put
\(r=D_1-C_4>0\).  If \(r=1\) or \(2\), specialize to three
variables.  The rank-three formula in the proof of
\Cref{lem:dense-four-lower-part} gives zero also at indices \(-1,-2\).
Thus the one missing entry already has its correct value in that matrix,
and both four-row endpoint determinants vanish.  At a strict boundary,
\(c_Q\) is a nonzero three-row Schur function and stays nonzero.
At a tied boundary it is, up to sign, \(D_T(D_3)\), where \(T\) is
the connected three-row maximal cofactor.  This difference is nonzero by
\Cref{lem:small-rows}: otherwise the opposite endpoint extensions would
be rotations and \(T\), hence its equal-margin border \(Q\), would
be symmetric.  A skew Schur function with at most three rows has Schur
expansion supported on partitions of length at most three, and
specialization to three variables is injective on their span.
Consequently \(c_Q\) remains nonzero.

For \(r\ge3\), specialize to the alphabet \((1,1)\), so
\(h_j\mapsto j+1\) for \(j\ge0\).  The virtual matrix
\((C_i-D_j+1)\) has rank at most two.  Correcting its one negative-index
entry from \(1-r\) to zero raises rank by at most one, so both
four-row endpoint determinants vanish.  A strict leading coefficient is
a three-row one-negative determinant; its value is, up to sign,
\(r-1\) times two positive row and column beta gaps, and is nonzero.
At a tie, with \(T=(C_2,C_3,C_4)/(D_1,D_2,D_3)\), the same correction gives
\begin{equation}\label{eq:one-negative-two-variable}
 D_T(t)(1,1)=(r-1)t\bigl((D_2-D_3)-(C_2-C_3)\bigr).
\end{equation}
If the inner gaps agreed, that relation, \(S-C_2=D_3\), and
\(S=C_4+D_1\) would give all three reflection equations for
\(Q=Q^*\).  Hence \(c_Q\) is nonzero on this specialization too.
In either case a vanishing specialized endpoint difference cannot equal
its nonzero specialized leading term, proving
\eqref{eq:one-negative-four-remainder}.
\end{proof}

\begin{lemma}\label{lem:one-negative-cross-resonance}
For a connected zero-free one-negative four-row \(Q\), both
\eqref{eq:cross-resonance-one} and \eqref{eq:cross-resonance-two} hold
with their stated hypotheses \(x>\ell>0\), \(u>0\), \(v=x+u\).
\end{lemma}

\begin{proof}
Write \(f=s_Q\), \(f_x=s_{Q^{\uparrow x}}\), and
\(f^{\,y}=s_{Q^{\downarrow y}}\).  For symmetric \(Q\), the first
inequality follows from factorization of distinct connected endpoint
functions against coordinate factors; the second is immediate.
For nonsymmetric \(Q\), normalize its lists as
\eqref{eq:one-negative-four-lists} and use \(\kappa,c_Q\) from
\Cref{lem:one-negative-endpoint}.  The cofactor
\[
 B_Q=\det(h_{C_i-D_j})_{2\le i\le4,\ 1\le j\le3},
 \qquad [h_S]f=-B_Q,
\]
is connected and one-negative: its subdiagonal indices are
\(C_3-D_1>0\) and \(C_4-D_2>0\).  It is therefore irreducible.

For \eqref{eq:cross-resonance-one}, the possible largest indices are
\(S+x,\ell+v,\kappa+v\).  If \(\ell<\kappa\), a tie of the first
and third has coefficient \(-h_{\ell+u}B_Q-h_\ell c_Q\), nonzero
modulo \(h_\ell\).  If \(\ell>\kappa\), a tie of the first and
second has coefficient \(-h_SB_Q-f\), nonzero modulo \(h_S\).
For \(\ell=\kappa\), the remaining coefficients are
\(-f-h_\ell c_Q\) or \(-h_SB_Q-f-h_\ell c_Q\).
The first is nonzero modulo \(h_\ell\); the second retains the
coefficient-one diagonal monomial of \(f\), all of whose indices are
below \(\kappa\).  Every strict maximum is nonzero.

The left side of \eqref{eq:cross-resonance-two} has largest index
\(\kappa+v\), with coefficient \(-h_\ell c_Q\).
For \(\ell<\kappa\), the right side can have that largest index only
at \(S=\kappa+v\), where its coefficient is not divisible by
\(h_\ell\).  For \(\ell>\kappa\), it has a larger index.
For \(\ell=\kappa\), it has either a larger index, a coefficient
\(f\) not divisible by \(h_\ell\), or a nonzero square of its top
generator.  None agrees with the left side.
\end{proof}

\begin{proposition}\label{prop:one-negative-three-component-rigidity}
A connected zero-free diagram whose maximal cofactor consists of a
one-negative four-row component and two singletons is rotation-rigid
against arbitrary connected mates.
\end{proposition}

\begin{proof}
Factorization of the maximal coefficient recovers the two singleton
functions and the connected four-row function.  By
\Cref{thm:two-negatives}, the latter determines \(Q\) up to rotation.
Normalize it as in \eqref{eq:one-negative-four-lists}.
Every mate is therefore an end or middle placement in
\eqref{eq:end-lists} or \eqref{eq:middle-lists}, with this one-negative
\(Q\).  Both singleton orders are included, with \(L=\ell+g_1\), \(R=m+g_2\)
and all six parameters positive.  The expansions
\eqref{eq:end-polynomial} and \eqref{eq:middle-polynomial} remain
identities: the changed entry stays inside their core determinants.

For \(x>\ell\), the intermediate \(U_Q(x)\) in
\eqref{eq:intermediate-U} now has precisely the negative sites
\[
 (3,1),\ (4,1),\ (5,1),\ (5,2).
\]
Thus \(\rho(U)=(2,1,1)\), while \(\rho(Z)=(3,1)\).
\Cref{prop:hook} gives arbitrary-mate rigidity of both.  Equal
U-type functions are literally equal, since rotation changes their
negative partition; equality of a U-type and a Z-type function forces
rotation.  Their lists recover the parameters as before.

The first and third markers are still
\eqref{eq:first-marker-table} and \eqref{eq:third-marker-table}.
The irreducibility and coordinate-divisibility exclusions in
\eqref{eq:cross-marker-exclusions} apply unchanged.  Its two remaining
ties are excluded by \Cref{lem:one-negative-cross-resonance}, so end and
middle placements cannot mix.

Within end placements, \eqref{eq:end-bottom-marker} recovers the
singleton order and the margins in the bottom-dominant regime.  The
moving-bottom complement in \eqref{eq:end-moving-bottom} remains
connected: its subdiagonal entries are
\(\ell,m,C_2-D_1,C_3-D_2,C_4-D_3>0\).
Irreducibility separates it from a product with singleton factors.
The fixed-core-orientation comparison \eqref{eq:end-orientation} uses
only the nonzero fixed endpoint leader, supplied here by
\Cref{lem:one-negative-endpoint}.  Hence aligned end placements agree
only literally.

Within middle placements, the tied movement coefficient in
\eqref{eq:middle-movement} is settled by the hook rigidity of U and Z.
The comparisons of marker types use
\Cref{lem:one-negative-cross-resonance}.  The remaining fixed-orientation
coefficient is \eqref{eq:middle-final-tie}.  Its strict maxima and square
coefficients are nonzero by the endpoint leader.  The final tie, with
\(\ell=m\), would require
\(D_Q(R)=h_{\kappa+R}c_Q\), contrary to
\Cref{lem:one-negative-four-lower}.  Middle placements therefore agree
only literally or under whole-parent rotation.
\end{proof}

\subsection{Singletons over exceptional five-row cores}

In this subsection let \(F_i,H_i\), \(i=1,2\), be
\eqref{eq:five-families} with \(b=p,g=G\), where
\(0<q<p<G\), and put \(S=p+4G\).
Zero-freeness is equivalent to \(p>q\): at \(p=q\), the last row
beta equals the penultimate column beta.  The four choices
\(F_i,H_i,F_i^*,H_i^*\) include both opening types and both orientations.
Set
\[
 a_j=h_{p+jG},\qquad b_j=h_{p+jG-q},\qquad
 a_j(t)=h_{p+jG+t},\quad b_j(t)=h_{p+jG-q+t}.
\]
Negative subscripts denote zero.  The generators \(a_j,b_j\) with
\(j\ge0\) are jointly independent.
Define
\begin{equation}\label{eq:exceptional-five-coefficients}
 \begin{aligned}
 D_1&=a_0^2b_2-a_0a_1b_1-a_0a_2b_0+a_1^2b_0,\\
 D_2&=a_0b_0b_2-a_0b_1^2+a_1b_0b_1-a_2b_0^2,\\
 W&=a_0b_1-a_1b_0,\qquad K_1=a_0^2,\quad K_2=a_0b_0.
 \end{aligned}
\end{equation}
For each \(i\), write
\[
 U_i(t)=s_{F_i^{\uparrow t}}-s_{H_i^{\uparrow t}},\qquad
 V_i(t)=s_{F_i^{\downarrow t}}-s_{H_i^{\downarrow t}},\qquad
 C_i(t)=U_i(t)+V_i(t).
\]

\begin{lemma}\label{lem:exceptional-five-responses}
The nonzero leading coefficients of these differences are
\begin{equation}\label{eq:exceptional-five-response-table}
 \begin{array}{c|c|c}
 \text{difference}&\text{largest index}&\text{coefficient}\\\hline
 U_i(t),\ t>0&p+3G+t&-b_0D_i\\
 V_i(t),\ t>0&p+3G+t&b_0D_i\\
 C_i(t),\ t>q&p+3G-q+t&a_0D_i\\
 s_{H_i^{\uparrow t}}-s_{F_i^{\downarrow t}},\ t>0
   &p+4G-q+t&-K_iW\\
 s_{F_i^{\uparrow t}}-s_{H_i^{\downarrow t}},\ t>0
   &p+4G-q+t&-K_iW.
 \end{array}
\end{equation}
For \(U_i,V_i\) and the two mixed differences, every fixed index is at
most the leading offset.  These offsets are positive and less than \(S\).
\end{lemma}

\begin{proof}
For rows \(r=4,3,2,1,0\), the four Jacobi--Trudi matrices have row
vectors
\begin{equation}\label{eq:exceptional-five-matrices}
 \begin{array}{c|l}
 F_1&a_{r-3},b_{r-1},a_{r-1},b_r,a_r\\
 H_1&b_{r-2},a_{r-2},a_{r-1},b_r,a_r\\
 F_2&b_{r-3},b_{r-1},a_{r-1},b_r,a_r\\
 H_2&b_{r-2},a_{r-2},b_{r-1},b_r,a_r.
 \end{array}
\end{equation}
Shift the first row for an up extension and the last column for a down
extension.  Cofactor expansion gives \eqref{eq:exceptional-five-response-table}.
For example, the \(a_3(t)\)-cofactor difference in \(U_1\) is
\(-a_0D_2-W^2=-b_0D_1\), since
\(b_0D_1-a_0D_2=W^2\); in \(U_2\) it is \(-b_0D_2\).
The mixed \(b_4(t)\)-cofactor is the triangular minor \(-K_iW\).
The \(a_3(t)\)-terms of \(U_i,V_i\) cancel in their sum, leaving
\(b_3(t)\) with coefficient \(a_0D_i\).

In \(U_i,V_i\), every fixed index is at most \(p+3G\).
For \(t>q\), the index of \(b_3(t)\) is above this bound.
In the mixed comparisons, the common corner term cancels; every fixed
index is at most \(p+4G-q\), and \(b_4(t)\) is the moving leader.
Finally, the \(a_2\)-coefficients of \(D_1,D_2\) are
\(-a_0b_0,-b_0^2\), and the \(a_1\)-coefficient of \(W\) is
\(-b_0\).  These are nonzero.  The stated bounds exclude cancellation
of any leader with a fixed index.
\end{proof}

\begin{proposition}\label{prop:exceptional-singleton-context}
Let \(Q\) be one of \(F_i,H_i,F_i^*,H_i^*\), and let \(P_Q\) be
its singleton context \eqref{eq:singleton-context}, with
\(s,u,v,g>0\).  Its complete connected class is its rotation orbit,
except when \(Q\in\{F_i^*,H_i^*\}\) and
\begin{equation}\label{eq:exceptional-singleton-locus}
 (s,u,v,g)=(p-q,q,G,G-p).
\end{equation}
Then the class consists of the two contexts over \(F_i^*,H_i^*\) and
their rotations.  These are the two seven-row families of
\Cref{prop:ten-negative-quartet}, with \((a,b,d)=(q,G,p-q)\):
family B for \(i=1\), and family A for \(i=2\).
\end{proposition}

\begin{proof}
Write \(N=S+s+g+u+v\) for the parent span.  Factorization recovers the
singleton length \(s\) and the five-row core function.
\Cref{thm:five-base} leaves precisely the four core
choices above.  Rotate any mate to put the singleton first.
The marker \eqref{eq:singleton-marker} distinguishes \(u<v\),
\(v<u\), and \(u=v\).  If \(v\le u\), its bottom endpoint
function distinguishes the four choices, by
\Cref{lem:exceptional-five-responses,lem:five-endpoint}.
The remaining margin is recovered by the factorization argument in the
proof of \Cref{lem:singleton-lift}, after this literal core recovery.
Such a parent is rigid.

For \(u<v\), the marker fixes \(u\).  Subtracting
\[
 s_Q(h_sh_N-h_{s+u}h_{N-u})
\]
leaves the separator \eqref{eq:margin-separator}.  Its derivation uses
only the common core function and span, not a prescribed representative.
The product and connected complementary determinant remain nonzero, and
the tied coefficient is their same-sign Schur-positive sum.
Thus \(v\), then \(g\), is fixed even across different core choices.
Every remaining comparison has identical \(s,u,v,g\).

Put \(A=s+g>s\).  For distinct core choices \(Q,R\), write
\[
 \Delta^+(t)=s_{Q^{\uparrow t}}-s_{R^{\uparrow t}},\qquad
 \Delta^-(t)=s_{Q^{\downarrow t}}-s_{R^{\downarrow t}}.
\]
The up differences have positive fixed leading offsets by
\Cref{lem:exceptional-five-responses,lem:five-negative-leading}.
In \(\Phi_Q(x,y)-\Phi_R(x,y)\), the corner term cancels.  Off the
tie \(x=y\), the largest index is \(S+\max(x,y)\), with the
corresponding nonzero endpoint difference as coefficient.
If \(v<A+u\), the second border difference in
\eqref{eq:singleton-polynomial} has the private largest index
\(S+A+u\).  If \(v>A+u\), the coefficient at \(S+v\) is
\[
 h_{s+u}\Delta^+(A)-h_s\Delta^+(A+u),
\]
which has a private second moving index since \(A>s\).
Equality therefore requires
\begin{equation}\label{eq:exceptional-interface-tie}
 v=A+u.
\end{equation}
At this tie, the coefficient of \(h_{S+v}^2\) cancels: it is the
common maximal cofactor of the equal cores.  The remaining coefficient is
\begin{equation}\label{eq:exceptional-interface-coefficient}
 E=h_{s+u}\Delta^+(A)-h_s\bigl(\Delta^+(v)+\Delta^-(v)\bigr).
\end{equation}
For opposite orientations of the same core the parenthesized sum is zero,
and \(E\ne0\).  Four cross-orbit pairings remain.

Cancel the common rows in their row multisets.  The remaining two or
three entries give
\begin{equation}\label{eq:exceptional-context-rows}
 \begin{array}{c|c|c}
 \text{pair}&i=1&i=2\\\hline
 F_i,H_i&A=G&A=G\\
 F_i^*,H_i^*&v=G&v=G\\
 F_i,H_i^*&(A,v)=(G,G)\text{ or }(3G-q,2G)
          &(G,G)\text{ or }(3G-q,2G-q)\\
 H_i,F_i^*&(A,v)=(G,G)\text{ or }(2G,3G-q)
          &(G,G)\text{ or }(2G-q,3G-q).
 \end{array}
\end{equation}
For example, the first mixed row in family 1 is
\[
 \{A+G,3G-q,v\}=\{A,2G,2G-q+v\}.
\]
The entry \(3G-q\) matches either \(A\) or \(2G-q+v\), giving
exactly the two listed possibilities.  Family 2 replaces \(A+G\) by
\(A+G-q\) and \(2G\) by \(2G-q\).  Interchanging the exterior
positions gives the last row.  The first two rows follow by comparing
the products of the two remaining entries.
Because \(u=v-A>0\), the third row and every \((G,G)\) alternative
are impossible.

For \(F_i,H_i\), one has \(A=G\), \(v=G+u\), and \(s<G\).
The candidate largest indices in \(E\) are \(S,s+u,S-q+u\).
The last exceeds \(s+u\), so cancellation forces \(u=q\).
At that value its \(h_S\)-coefficient is
\[
 -D_i(h_{s+q}b_0+h_sa_0)\ne0.
\]
For \(F_i^*,H_i^*\), one has \(v=G\), \(A=G-u\), and
\(s+u<G\).  The leaders \(S-u,S-q\) force \(u=q\), and their
coefficient is
\[
 D_i(h_{s+q}b_0-h_sa_0).
\]
Equality of these coordinate monomials forces \(s=p-q\), hence
\eqref{eq:exceptional-singleton-locus}.

For \(H_i,F_i^*\) in family 1, the row conditions give
\((A,v,u)=(2G,3G-q,G-q)\).  The first term of \(E\) has leader
\(p+6G-q\), above the other leader \(p+6G-2q\).
For family 2, they give \((A,v,u)=(2G-q,3G-q,G)\).
The common leading coefficient, apart from the nonzero factor \(a_0\), is
\[
 h_sD_2-h_{s+G}b_0W.
\]
For \(s<p+G\) its private largest index is \(p+2G\); for
\(s>p+G\) it is \(s+G\).  At \(s=p+G\), the
\(a_2=h_{p+2G}\)-coefficient is \(-a_0b_0b_1\ne0\).
This includes the admissible tie \(p+q<G\).

All six unordered core pairs have been compared.  The remaining locus
is sufficient by the identities in \Cref{prop:ten-negative-quartet}.
Its four members are distinct, so this is the complete class.
\end{proof}

\subsection{The seven-row base and its borders}

\begin{proposition}\label{prop:ten-negative-base}
Every connected zero-free seven-row diagram with \(\nu=10\) and
\(n_{\min}=7\) is rotation-rigid, except for the two complete
four-member families in \Cref{prop:ten-negative-quartet}, with \(d>0\).
\end{proposition}

\begin{proof}
The maximal cofactor has six positive rows and is disconnected by
\Cref{lem:above-minimum-border}.  The complete budget is
\begin{equation}\label{eq:ten-negative-budget}
 \begin{array}{c|c|c}
 \text{component orders}&\text{cross cost}&\text{internal negatives}\\\hline
 1+5&5&5\\
 2+4&8&2\\
 3+3&9&1\\
 1+1+4&9&1.
 \end{array}
\end{equation}
Other three-component orders cost at least eleven, and four components
cost at least twelve.  Two-row connected components are dense, and a
three-row component has at most one negative.

For \(2+4\), apply \Cref{lem:two-component-lift} using
\Cref{thm:two-negatives,lem:small-rows,lem:one-negative-endpoint,lem:two-negative-endpoint}.
For \(3+3\), the same lifting lemma uses one dense and one one-negative
component, with their endpoint leaders from
\Cref{lem:one-negative-endpoint}.  The \(1+1+4\) case is
\Cref{prop:one-negative-three-component-rigidity}.
For \(1+5\), \Cref{thm:five-base} says that the five-row core is
rigid or belongs to one of its two exceptional classes.  Apply
\Cref{lem:singleton-five-rigid} in the first case and
\Cref{prop:exceptional-singleton-context} in the second.
Each argument compares with arbitrary connected mates.
\end{proof}

The five distinct gap values in either seven-row family are
\[
 d<a+d<b+d<a+b+d<2b+d.
\]
Family A repeats the third value and family B the fourth.  Their gap
multisets therefore determine the family and \(d,a,b\).  The four
ordered gap words within each class are pairwise different.

\begin{lemma}\label{lem:ten-negative-endpoints}
Let \(R\) be connected and zero-free with seven rows, ten negatives,
and minimum realizing row count seven.  For \(t>0\),
\(s_{R^{\uparrow t}}=s_{R^{\downarrow t}}\) if and only if \(R=R^*\).
Every positive endpoint extension of a member of either exceptional
seven-row quartet is rotation-rigid.  Endpoint extensions from its two
different rotation orbits cannot have equal functions.
\end{lemma}

\begin{proof}
Both extensions stay in the domain of \Cref{prop:ten-negative-base}
and retain the same ordered gap word.  The four gap words in an
exceptional class are distinct, so two distinct opposite endpoint
extensions cannot form an exceptional pair.  The base classification
makes them rotations, and undoing the common top extension gives
\(R=R^*\).  The converse follows from rotation.

A positive endpoint extension of an exceptional member has the same gap
multiset and larger area.  The gap multiset fixes its family's parameters,
so the extension cannot be exceptional; it is rigid by the base theorem.
Cross-orbit extensions cannot be identical or rotations, because the
original gap words are different even after reversal.  Equality of their
functions would force a common extension size by weighted degree, and
then contradict this rigidity.
\end{proof}

\begin{lemma}\label{lem:ten-negative-border-separation}
Every positive border over any member of either seven-row family in
\Cref{prop:ten-negative-quartet} is rotation-rigid against arbitrary
connected mates.  Common betas are allowed.
\end{lemma}

\begin{proof}
First take \(d>0\).  Factorization and \Cref{prop:ten-negative-base}
recover the connected seven-row child up to its four-member class.
\Cref{lem:border} recovers the unordered margins.  Unequal margins
would equate cross-orbit endpoint functions, excluded by
\Cref{lem:ten-negative-endpoints}.  Within one child orbit, the same
lemma excludes every margin swap other than whole-parent rotation.

At equal margins \((t,t)\), a core rotation only rotates the parent.
For opposite child orbits, canceling common parent rows leaves
\begin{equation}\label{eq:ten-negative-border-rows}
 \begin{array}{c|c}
 A&\{3b+d,3b+d+t\}=\{4b+d,2b+d+t\},\\
 B&\{a+3b+d,3b+d+t\}=\{4b+d,a+2b+d+t\}.
 \end{array}
\end{equation}
Each forces \(t=b\).

Put \(U_j=h_{jb+d}\), \(V_j=h_{a+jb+d}\), with negative
subscripts interpreted as zero.  With columns indexed \(7,6,\ldots,0\),
a symbol \(U_r\) below means the row \((U_{r-j})_j\), and similarly
for \(V_r\).  The four bordered matrices have rows
\begin{equation}\label{eq:ten-negative-border-matrices}
 \begin{array}{c|l}
 A,F&V_8,V_7,U_7,V_6,U_6,U_5,V_4,U_4\\
 A,H&V_8,V_7,U_7,V_6,U_6,V_5,U_5,U_3\\
 B,F&V_8,V_7,U_7,V_6,U_6,V_5,V_4,U_4\\
 B,H&V_8,V_7,U_7,V_6,U_6,V_5,U_5,V_3.
 \end{array}
\end{equation}
Their differences satisfy
\begin{equation}\label{eq:ten-negative-border-separators}
 \begin{aligned}
 [U_1^4V_0V_3V_5V_7]
 (s_{\B_{b,b}F_A}-s_{\B_{b,b}H_A})&=-1,\\
 [U_1^3V_0V_1V_3V_5V_7]
 (s_{\B_{b,b}F_B}-s_{\B_{b,b}H_B})&=-1.
 \end{aligned}
\end{equation}
For A the source contributes nothing, and the mate contributes the
positive column word \((7,1,2,5,3,8,4,6)\).
For B both share the positive word \((1,8,2,7,3,6,4,5)\), and only
the mate has the preceding additional word.  The forced \(U_1\)-sites
and the distinct \(V\)-indices exhaust the possibilities.
The alphabets are independent since \(0<a<b\), so these coefficients
exclude the tie.

For \(d=0\), homogenization sends the family parameters to
\((2a,2b,2d+1)\) and the two border margins to twice themselves.
The zero-free result descends by \Cref{lem:homogenization}.
\end{proof}

\begin{proof}[Proof of \Cref{thm:ten-negative-classification}]
The zero-free seven-row base is \Cref{prop:ten-negative-base}.
At eight rows the maximal cofactor is connected, has ten negatives, and
has minimum realizing row count seven by \Cref{lem:above-minimum-border}.
If that child is rigid, align it and apply \Cref{lem:border} with the
opposite-endpoint isolation in \Cref{lem:ten-negative-endpoints}.
If it is exceptional, use \Cref{lem:ten-negative-border-separation}.
Every eight-row source is therefore rigid.

Above eight rows, the child and its positive endpoint extensions have at
least eight rows and retain the same negative partition.  The inductive
rigidity assertion and \Cref{lem:border} give rigidity at the next
height.  This argument does not require endpoint extensions of an
arbitrary rigid seven-row source to be rigid.

Apply \Cref{lem:homogenization} to a full-board equality at its common
row count.  At height at least eight, rigidity descends immediately.
At height seven, a nonrotation image belongs to family A or B.
Rotate to a displayed orientation, whose first three gaps are
\(d',a'+d',b'+d'\).  All image gaps are odd, so \(d'\) is odd and
\(a',b'\) are even.  The inverse parameters
\[
 a=a'/2,\qquad b=b'/2,\qquad d=(d'-1)/2
\]
satisfy \(a\ge1,d\ge0,b>a+d\).  Reversing the homogenization gives
exactly the displayed full-board family and all four of its members.
Injectivity and rotation covariance exclude additional mates.
The boundary \(q=p\) in the five-row core corresponds to \(d=0\);
\(G=p+1\) gives separating gap one.  Both are included.
\end{proof}

Together with \Cref{cor:ten-negative-high-minimum}, this classifies all
connected sources with \(\nu=10\) and \(n_{\min}\ge7\).
The minimum-six stratum is not included.

\section{Concluding remarks}\label{sec:conclusion}

The classification through six negatives has only the seven displayed
exceptional schemes, all confined to four or five rows.  When a negative
pattern requires at least seven rows, nine negatives force rigidity.
At ten negatives under the same minimum-row condition, the two seven-row
families in \Cref{thm:ten-negative-classification} are the complete
exceptions; every such source with at least eight rows is rigid.  Minimum realizing row
count, unlike the negative partition itself, is constant on a connected
Schur class.  It must not be replaced by the actual row count.

A census of connected skew diagrams with eight through thirteen cells
gives eighteen Schur classes containing more than one rotation orbit.
It enumerates full profiles, identifies diagrams up to translation and
rotation, and counts each Schur class once rather than counting pairs of
diagrams.  Nine classes lie in this paper's classification domain; three
further classes are classical ribbon classes
\cite{Billera2006DecomposableCompositions}.  The remaining six are
nonribbon classes outside this domain, for which no geometric account is
asserted here.  This is a bounded frontier, not a proportion of the general
problem solved and not a claim that those six classes admit no legal
presentations.

The six-row quartets in \Cref{prop:six-row-quartets} exhibit three
remaining cases: a nonlinear component can itself have nonrotation mates,
a rigid non-dense component can change placement, and four components
require a further assembly comparison.  Their displayed members establish
no upper class bound.  \eqref{eq:ten-chart} supplies legal reversal
presentations for the two ten-negative families.

\begin{problem}\label{prob:quartet-exhaustion}
What are the complete connected classes containing the seven-, eight-,
and nine-negative six-row quartets in \Cref{prop:six-row-quartets}?
Four constructed members give a lower inclusion, not an upper bound.
\end{problem}

\begin{problem}\label{prob:placement-resonances}
Which non-dense components admit a mixed endpoint resonance producing
different cofactor placements?  A criterion must distinguish an
irreducible nonlinear cofactor from a product of singletons of the same
ordinary degree, as \Cref{prop:nondense-resonance} requires.
\end{problem}

\appendix
\section{The four-row calculation}\label{app:four}

\begin{proof}[Proof of \Cref{prop:four-rows}]
We may first work zero-free by \Cref{lem:homogenization}.  The maximal
cofactor has three positive rows.  Its factorization distinguishes the
three possibilities: one connected three-row component, a singleton
and a connected two-row component, or three singletons.
The arbitrary mate has the same order and zero-freeness by
\Cref{lem:readings}, and the same component inventory by irreducibility.

The connected case is \Cref{lem:small-rows,lem:border}.  In the
singleton-plus-two-row case, rotate to put the singleton first.  The
two-row component is dense; \Cref{lem:singleton-rigidity} applies.
Only its dense clause is used here, whose proof depends on
nonnegative-index rigidity, not on \Cref{thm:two-negatives} or the
present four-row classification.

Three singletons mean exactly
\(r_2>z_1+z_2\), \(r_3>z_2+z_3\), the upper Hessenberg case.
Use the eight terms of \eqref{eq:interval-expansion} and put
\(c=\min(z_1,z_2,z_3)\).  The generator \(h_c\) occurs only at minimum
subdiagonal sites.  Setting it to zero splits the board into connected
blocks of at most three rows.  Their irreducible polynomials recover the
unordered blocks up to rotation.  Their row counts distinguish the
\(1+3\) and \(2+2\) splits.

\emph{One minimum at an end.}
Orient the profile as \((p,q,r,s;c,a,b)\), with \(a,b>c\), and put
\(t=p-c>0\).  Only reversal of the three-row block remains.
The coefficient of \(-h_c\) is
\[
 h_{q+t}h_rh_s-h_{q+r+t-a}h_ah_s
 -h_{q+t}h_{r+s-b}h_b+h_Nh_ah_b.
\]
If \(a<b\), its private \(h_a\)-coefficient, compared with the reversed
block, forces \(s=q+r-a\).  After substitution the remaining difference is
\[
 h_r(h_{q+t}h_s-h_{s+t}h_q)
 -h_b(h_{q+t}h_{r+s-b}-h_{r+s+t-b}h_q).
\]
Its largest label \(r+s+t-b\) is unique and has nonzero coefficient;
\(r>a+b\) verifies the strict comparisons.  The case \(a>b\) reverses
this argument.  If \(a=b,q\ne s\), the largest remaining label is
\(r+t+\max(q,s)-a\), again with nonzero coefficient.  For \(q=s\)
the block was already symmetric.

\emph{One minimum in the middle.}
The profile is \((p,q,r,s;a,c,b)\), with \(a,b>c\).
Fix the order of the two recovered two-row blocks by global rotation.
The coefficient of \(-h_c\) is
\[
 h_{q+r-c}h_ph_s-h_{p+q+r-a-c}h_ah_s
 -h_{q+r+s-b-c}h_bh_p+h_Nh_ah_b.
\]
Flipping only the left block changes it, up to sign, by
\[
 h_s(h_ph_{q+u}-h_qh_{p+u})
 -h_b(h_ph_{q+v}-h_qh_{p+v}),
\]
where \(u=r-c>0\), \(v=r+s-b-c>u\).
If \(p\ne q\), the largest label \(\max(p,q)+v\) survives.
The right-block argument is its rotation.
For \(a<b\), the private \(h_a\)-coefficient fixes the right-block
orientation; \(b<a\) is symmetric.  For \(a=b\), comparing both-block
flips and taking the part of the \(h_c\)-coefficient free of \(h_a\)
gives
\[
 h_{q+r-c}h_ph_s=h_{p+s-c}h_qh_r.
\]
The merged labels exceed their summands.  Their equality forces
\(q+r=p+s\), then \(\{p,s\}=\{q,r\}\).  Either both blocks are
symmetric, or \(p=r,q=s\), giving global rotation.

\emph{Three minima.}
Subtract \(c\) from every row and gap.  The index substitution on the
homogeneous degree-four determinant is injective when the unit variable
is retained, and the result is a four-row ribbon.  Such a ribbon has no
nonrotation mate by \cite{Billera2006DecomposableCompositions}:
two nonsymmetric composition factors would require at least five rows,
using \(\ell(\alpha\circ\beta)=\ell(\alpha)+|\alpha|(\ell(\beta)-1)\),
since each has at least three cells and two rows.
Restoring \(c\) is injective and commutes with rotation.

\emph{Two minima.}
Let the gaps be \(\{c,c,d\}\), \(d>c\).  The cut recovers singleton
rows \(p,q>c\) and a block with rows \(a,b>d\) and gap \(d\).
Modulo global rotation all arrangements have one of the forms
\[
 L=(p,q,a,b;c,c,d),\qquad M=(p,a,b,q;c,d,c),
\]
allowing both singleton orders and both block orientations, subject to
the Hessenberg inequalities.  Put \(t=a+b-d\), \(N=p+q+t-2c\).
The coefficients of \(h_c^2\) and \(h_c\) are
\begin{align*}
 Q_L&=h_{p+q+a-2c}h_b-h_Nh_d,&
 Q_M&=h_{p+a-c}h_{q+b-c}-h_Nh_d,\\
 \ell_L&=-h_{p+q-c}h_ah_b+h_{p+q-c}h_th_d
          -h_ph_{q+a-c}h_b+h_ph_{q+t-c}h_d,\\
 \ell_M&=-h_{p+a-c}h_bh_q-h_ph_ah_{q+b-c}
          +h_{p+t-c}h_dh_q+h_ph_{q+t-c}h_d.
\end{align*}
All remaining indices exceed \(c\), and the constant terms agree.
For \(L\) against \(L\), quadratic equality fixes the block orientation;
swapping unequal singletons then has largest label
\(\max(p,q)+t-c\) with coefficient \(\pm h_{\min(p,q)}h_d\ne0\).
For \(M\) against \(M\), the quadratic terms agree only when
\(p=q\) or \(a=b\), giving rotation or identity.

For \(L(p,q,a,b)\) against \(M(p,a,b,q)\), quadratic equality forces
\(b=p+a-c\).  In the linear difference the largest possible labels
are \(L_1=q+b-c\) and \(L_2=p+t-c\).  Their coefficients are
\(h_ph_a\), \(-h_qh_d\), or at the tie \(h_ph_a-h_qh_d\).
Only the tie can vanish, and it forces \(p=d,q=a\), because \(a>d\).
Then \(b=a+d-c\), and the source inequality is \(a>c+d\).
Renaming \(d\) as \(w\) and \(a\) as \(b\) gives
\eqref{eq:four-family}.  Substituting these equalities in the eight
terms proves full polynomial equality.
For the other alignment \(M(p,b,a,q)\), quadratic equality instead
forces \(b=q+a-c\); the linear difference has unique largest label
\(p+t-c\), with coefficient \(-h_qh_d\ne0\).
These comparisons exhaust all arrangements, including coincident rows.

The zero-free exceptional parameters satisfy \(1\le c<w\), \(b>c+w\).
For a \(\Gamma\)-image its bare profile parameters \(c,w,b\) are odd.
Replacing them by \((c-1)/2,(w-1)/2,(b-1)/2\) gives exactly
\(0\le c<w\), \(b>c+w\); conversely substitution in the interval
formula gives the identity at \(c=0\) as well.
The different overlap words distinguish the two rotation orbits, and
the displayed repeated-row and gap data recover their parameters.
The three strictly negative indices follow from
\(2c-b<0\), \(c+w-b<0\), and the still smaller bottom-left index.
No other index is negative.  This proves the complete full-board statement.
\end{proof}

\section{First-order determination at a contact}\label{app:contact-first-jet}

The \(q\)-components and the augmentation \(\varepsilon\) are those
of Section~\ref{sec:staircase}.  The result below determines a component
function from its first jet; it does not recover that jet from a parent
polynomial.

\begin{theorem}[First-jet rigidity]\label{thm:staircase-first-jet}
Let \(A,B\) be \(q\)-components with at most four rows.  Then
\[
 s_A\big|_{h_q=0}=s_B\big|_{h_q=0},\qquad
 [h_q]s_A=[h_q]s_B
 \quad\Longrightarrow\quad s_A=s_B.
\]
Equivalently, equal restrictions of distinct full functions satisfy
\(\nu_{h_q}(s_A-s_B)=1\), where \(\nu_{h_q}\) is the order of
divisibility by the independent generator \(h_q\).
\end{theorem}

\begin{proof}
Put \(S_j(A)=[h_q^j]s_A\).  The inherited boundary preserves the private
full-interval term in \(S_0\), whose coefficient is
\((-1)^{n-1}\prod_i h_{g_i}\).  Thus \(S_0\ne0\) recovers the row
count, full interval, and gap multiset.  Equal restrictions have the same
row count, so we work at each order separately.

For four rows use \((a,b,c,d;r,s,t)\) and the indices of
\eqref{eq:staircase-four-expansion}.  The boundary gives
\(d,w,y,T>q\).  Since \(u>\max(a,b)\), \(v>\max(b,c)\), and
\(x>\max(u,v)\), the following list of contact sets is exhaustive.
A word such as \(ac\) means that exactly \(a,c\) among
\(a,b,c,u,v,x\) equal \(q\); comma-separated words denote
alternatives, not simultaneous contacts.  Equalities among the other
indices are allowed.  The independent generators are collected before
applying \(\varepsilon\).  In particular, the boundary excludes an
\(h_q^4\)-term; the only possible unknown coefficients after the
first jet are \(S_2,S_3\).
\[
\begin{array}{c|c|c|c}
\text{indices equal to }q&(\varepsilon(S_0),\varepsilon(S_1))&S_2&S_3\\\hline
\varnothing,a,b,c,u&(0,0)&0&0\\
v,uv&(1,-1)&0&0\\
x&(-1,1)&0&0\\
ab&(0,0)&R_t(c,d)&0\\
ac&(1,-2)&h_bh_d&0\\
bc&(1,-2)&h_ah_d&0\\
av&(0,1)&-h_sh_d&0\\
uc&(-1,2)&-h_rh_d&0\\
abc&(1,-1)&-h_th_{q+d-t}&h_d.
\end{array}
\]
All higher coefficients vanish.  The table is obtained by substitution
in the displayed eight terms, retaining every numerical coincidence.
It suffices to prove that the first jet determines the listed higher
coefficients whenever one of the last six cases occurs.

For \(av\), the scalar jet identifies the case.  Its first coefficient is
\[
 h_bh_ch_d-h_th_bh_{c+d-t}+h_sh_th_{q+d-t}.
\]
The last index is private and largest, with coefficient \(h_sh_t\).
Removing that term leaves \(h_bR_t(c,d)\).  Its monomial content
recovers \(b\); its primitive binomial recovers \(t,\{c,d\}\).
Since \(c<q<d\), the pair is oriented, and \(s=b+c-q\).
Thus \(S_2\) is recovered.

For \(uc\), again the scalar jet identifies the case, and
\[
 S_1=h_ah_bh_d+h_rh_th_{q+d-t}.
\]
The last index is private.  After removing its term, the only index above
\(q\) in the remaining monomial is \(d\).  The last index then
recovers \(t\), and its coefficient recovers \(r\).  Hence
\(S_2=-h_rh_d\) is fixed.

For \(bc\),
\[
 S_1=-h_rh_{a+q-r}h_d-h_th_ah_{q+d-t}.
\]
Every coefficient is negative.  This distinguishes it from \(ac\),
whose first coefficient has a private positive largest term, as below.
For \(bc\), \(a\ne q\), so
\Cref{lem:staircase-endpoint-coefficients} applies directly to \(S_0\).
It recovers \(a,r\) after a valid orientation, or both endpoint pairs
at a tie.  The middle rows are \(q,q\), so the recovered full interval
and gap multiset give \(d=T-a-2q+r+s+t\).  Hence
\(S_2=h_ah_d\) is determined.  No residual-classification theorem
is used in this first-jet argument.

For \(ac\), the first coefficient is
\begin{align*}
 S_1={}&-h_rh_{q+b-r}h_d-h_sh_{q+b-s}h_d
        -h_th_bh_{q+d-t}\\
       &+h_sh_th_{b+q+d-s-t}.
\end{align*}
Its last index \(y\) is strictly largest and has positive coefficient.
The known difference \(T-y=q-r\) recovers \(r\).  Put \(U=q+b-r\).
The restriction is
\[
 S_0=-h_rP_{(U,q,d;s,t)},\qquad U>q+s,\quad d>q+t.
\]
This middle-row three-row restriction determines the contracted profile
up to rotation.  To see this directly, its largest coefficient recovers
\(\{s,t\}\).  If the gaps differ, its two negative monomials are
identified by their unique gap factors below \(q\); swapping both
high-index pairs would require a positive shift in one equation and a
negative shift in the other.  If \(s=t\), their least high index and
common total recover \(\{U,d\}\), including coincident monomials.

The only possible distinct endpoint lift is therefore, with
\(\delta=q-r>0\),
\[
 (q,U-\delta,q,W;r,s,t)
 \quad\text{against}\quad
 (q,W-\delta,q,U;r,t,s).
\]
Their contracted full functions are rotations.  The first-coefficient
comparison is consequently between the two tail restrictions.  After
the common full-interval term is removed, their only possible largest
indices are \(q+W-t\) and \(q+U-s\), with coefficient contributions
\(-h_t h_{U-\delta}\) and \(h_s h_{W-\delta}\).
The barred specializations do not vanish: those two rows are not \(q\)
in the \(ac\) case.  At a tie, a crossed matching is impossible because
\(U-\delta>s\) and \(W-\delta>t\).  Hence equality forces
\(U=W,s=t\), and the lifts coincide.  This recovers \(S_2\).

For \(abc\),
\[
 S_1=-h_rh_{2q-r}h_d-h_sh_{2q-s}h_d
          +h_sh_th_{2q+d-s-t}.
\]
Its last term is private and largest.  The competing lower-degree
\(v\)-case has a single negative first-coefficient monomial, so cannot
match.  The \(uv\)-case has
\[
 S_1=-h_rh_ch_d-h_sh_ah_d+h_rh_th_{c+d-t}.
\]
If its private positive term matched, subtract it on both sides.
Each remaining \(abc\) monomial has exactly one index below \(q\),
whereas each remaining \(uv\) monomial has two.  They cannot coincide.
For two \(abc\) profiles, the gap multiset and \(T\) give \(d\).
The largest first-coefficient term recovers \(\{s,t\}\), hence \(r\).
Removing it and the known \(r\)-term leaves
\(-h_sh_{2q-s}h_d\), whose low index recovers \(s\), then \(t\).
Both \(S_2\) and \(S_3\) are determined.

The remaining case is \(ab\).  Its first coefficient is
\[
 S_1=-h_sR_t(q+c-s,d),\qquad c\ne q.
\]
The lower-degree competitors with scalar jet \((0,0)\) have first
coefficients
\[
 0,\quad s_{(b,c,d;s,t)},\quad h_aR_t(c,d),\quad
 h_dR_r(a,b),\quad -h_rR_t(c,d),
\]
for contact sets \(\varnothing,a,b,c,u\), respectively.
Under \(h_j=\epsilon^{j^2}\), a full Schur polynomial has positive
initial coefficient by \Cref{lem:readings}; multiplication by a monomial
preserves that sign.  Thus only the negative \(u\)-case can compete.

Write this candidate as \((a',b',C,D;s,r,t)\).  Equality of first
coefficients and the recovered gap multiset force
\[
 a'+b'-s=q,\qquad
 \{C,D\}=\{q+c-s,d\},\qquad C,D>q+t.
\]
The source restriction is
\(-h_r s_{(2q-r,c,d;s,t)}\).  Divisibility of the candidate restriction
by \(h_r\), after cutting its middle seam, forces \(a'=r\):
\(b'>r\) and \(C,D>q>r\), and the remaining two-row factor does not
vanish.  Consequently
\[
 r>s,\qquad q>2r,\qquad b'=B_0=q+s-r.
\]
Set
\begin{gather*}
 v_0=B_0+C-r,\quad x_0=q+C-r,\quad w_0=C+D-t,\\
 y_0=B_0+C+D-r-t,\quad T_0=B_0+C+D-s-t.
\end{gather*}
The candidate restriction divided by \(h_r\) is
\begin{align*}
 &h_{B_0}h_Ch_D-h_rh_{v_0}h_D-h_th_{B_0}h_{w_0}\\
 &\qquad+h_sh_{x_0}h_D+h_rh_th_{y_0}-h_sh_th_{T_0}.
\end{align*}
Here
\[
 T_0>y_0>w_0>\max(x_0,D),\qquad
 x_0>v_0>C>B_0>r>s.
\]
The three largest-index terms cannot cancel.  After canceling the common
\(h_D\) from the remaining potential pairings, their least indices and
\(x_0>v_0\) exclude every opposite-sign coincidence.  The norm is six.
The source restriction divided by \(h_r\) is a full three-row Schur
function of norm four: its full interval is private, and its diagonal
cannot equal an adjacent-merge term because the merged index exceeds both
rows.  Same-sign collisions preserve this norm.  The candidate is
therefore impossible.

For two \(ab\) profiles, the first coefficient recovers
\(s,t,\{c+q-s,d\}\), and the gap multiset recovers \(r\).
The source restriction divided by \(-h_r\) is a full three-row
polynomial, whose minimum-square monomial recovers
\(\{2q-r,c,d\}\).  Remove one occurrence of \(2q-r\).
The two multisets \(\{c,d\}\) and \(\{c+q-s,d\}\), with \(q-s>0\),
recover \(c,d\) in order: subtract their sums of squares.
Hence \(S_2=R_t(c,d)\) is fixed.

This exhausts the four-row table.  With three rows, the boundary gives
\(c,b+c-s,a+b+c-r-s>q\).  Degree at least two in \(h_q\) occurs only
when \(a=b=q\), with
\[
 S_2=h_c,\qquad S_1=-h_sh_{q+c-s},\qquad
 (\varepsilon(S_0),\varepsilon(S_1))=(0,-1).
\]
The other scalar jets are \((0,0),(-1,1),(1,-1)\).
Thus this case is identified and its first coefficient recovers \(s,c\).
Two-row polynomials have degree at most one, and a singleton is independent
of \(h_q\).  The theorem follows.
\end{proof}

\begin{remark}
The inherited inequality \(d>q+t\) preserves the full-interval term
and confines the possible \(q\)-indices to the contact table.
On the norm-four locus \(c=q,s=t,d=r,a=q+r-s\), one has
\(r>q+s\), \(b>r+s\), so only \(c\) equals \(q\).
For example, at \(A=(8,8,3,6;6,1,1)\) and \(q=3\),
\[
 [h_3]s_A=h_6h_8^2-h_6^2h_{10},\qquad
 [h_3^j]s_A=0\quad(j\ge2).
\]
The cancellation in the restriction does not conceal a higher
\(h_q\)-coefficient.  The theorem compares the original component
functions, not merely their primitive associates, and it does not identify
their first jets from a parent polynomial.  Indeed, for \(e>q\),
\[
 [h_q]s_{A\cdot_q(e)}
 =h_e[h_q]s_A-s_{A\odot_q(e)}\big|_{h_q=0}.
\]
The contraction term remains part of global compatibility.
The first-jet theorem is an additional component-reconstruction result;
the classification in Section~\ref{sec:staircase-assembly} uses
the separate compatibility lemmas.
\end{remark}

In congruence notation, the theorem reads
\begin{equation}\label{eq:farkas}
 s_A\equiv s_B\pmod{h_q^2}
 \quad\Longleftrightarrow\quad s_A=s_B
 \qquad\text{for \(q\)-components with at most four rows}.
\end{equation}

\section*{Declarations}

{\footnotesize
\noindent\textbf{Funding.} This research received no external funding.

\smallskip
\noindent\textbf{Competing interests.} The author has no relevant interests to disclose.

\smallskip
\noindent\textbf{Data availability.} No datasets were generated or analyzed.

\smallskip
\noindent\textbf{Computational verification.} The proofs do not depend on numerical computation.  Scripts used for verification are available at \url{https://github.com/johnnyshields/papers}.

\smallskip
\noindent\textbf{Generative AI use.} During the preparation of this work, the author used OpenAI's ChatGPT (GPT-6 Astra, GPT-5.6 Sol) and Anthropic's Claude (Opus 5.5, Opus 5) to assist with mathematical research, manuscript drafting, and the development of auxiliary verification code.  After using these tools, the author reviewed and edited as needed.  The author takes full responsibility for the content of the work.
\par}

\begin{thebibliography}{99}

\bibitem{Barekat2009CoincidencesCorrigendum}
F. Barekat, V. Reiner and S. van Willigenburg,
\emph{Corrigendum to ``Coincidences among skew Schur functions''},
Adv. Math. \textbf{220} (2009), no.~5, 1655--1656,
DOI: \href{https://doi.org/10.1016/j.aim.2008.10.010}{10.1016/\allowbreak j.aim.2008.10.010}.

\bibitem{Billera2006DecomposableCompositions}
L.~J. Billera, H.~Thomas and S.~van Willigenburg,
\emph{Decomposable compositions, symmetric quasisymmetric functions and equality of ribbon Schur functions},
Adv. Math. \textbf{204} (2006), no.~1, 204--240,
DOI: \href{https://doi.org/10.1016/j.aim.2005.05.014}{10.1016/\allowbreak j.aim.2005.05.014}.

\bibitem{GarciaGarcia2020ToeplitzMinors}
D.~Garc\'ia-Garc\'ia and M.~Tierz,
\emph{Toeplitz minors and specializations of skew Schur polynomials},
J. Combin. Theory Ser. A \textbf{172} (2020), 105201,
DOI: \href{https://doi.org/10.1016/j.jcta.2019.105201}{10.1016/\allowbreak j.jcta.2019.105201}.

\bibitem{Macdonald1995SymmetricFunctions}
I.~G. Macdonald,
\emph{Symmetric Functions and Hall Polynomials}, 2nd ed.,
Oxford Mathematical Monographs, Oxford University Press, 1995,
DOI: \href{https://doi.org/10.1093/oso/9780198534891.001.0001}{10.1093/\allowbreak oso/\allowbreak 9780198534891.001.0001}.

\bibitem{McNamara2009Classification}
P.~R.~W. McNamara and S.~van Willigenburg,
\emph{Towards a combinatorial classification of skew Schur functions},
Trans. Amer. Math. Soc. \textbf{361} (2009), no.~8, 4437--4470,
DOI: \href{https://doi.org/10.1090/S0002-9947-09-04683-2}{10.1090/\allowbreak S0002-9947-09-04683-2}.

\bibitem{Reiner2007Coincidences}
V.~Reiner, K.~M. Shaw and S.~van Willigenburg,
\emph{Coincidences among skew Schur functions},
Adv. Math. \textbf{216} (2007), no.~1, 118--152,
DOI: \href{https://doi.org/10.1016/j.aim.2007.05.006}{10.1016/\allowbreak j.aim.2007.05.006}.

\bibitem{Shields2026Cofactor}
J.~Shields,
\emph{Cofactor independence and dilation for skew Schur functions},
preprint, 2026,
DOI: \href{https://doi.org/10.5281/zenodo.22732784}{10.5281/\allowbreak zenodo.22732784}.


\end{thebibliography}
\end{document}
